% Modernised reading — fonds Alexandre Grothendieck, shelfmark 57, pages 1-314.
%
% This is an INTERPRETATION, not a transcription. It re-reads the pages in
% today's notation and vocabulary, states the hypotheses the manuscript leaves
% implicit, and drops the critical apparatus so that the mathematics reads
% straight through. Everything the brackets and underlines carried — a doubtful
% reading, a notation he used differently, a missing passage — moves into a
% footnote.
%
% It is held to being mathematically correct as it stands, not merely faithful
% to the page. Where the manuscript is loose or mistaken, the true statement is
% what appears here, and the footnote says what the page has. In any dispute of
% substance the transcription governs: whoever cites Grothendieck must cite the
% batch-NN.fr.tex files.
%
% Scope: ONE document for the whole shelfmark, its sixteen batches read as the
% workshop they are. The folder is the material of a treatise on the Picard
% scheme that its own page 65 plans in twelve chapters; the stations below are
% read against that plan. Several threads cross batch boundaries and are
% written once here: the separation of Pic (letter to Murre p. 20-21, the
% rigidity notes p. 31, the units lemma p. 115-123, the non-separated example
% p. 213-218, the open questions p. 312); the comparison of a local ring with
% its completion (typescript p. 2-12, cones p. 184-193, Hironaka letter
% p. 226-230, analytic draft p. 247-271); and the theorem of the cube
% (calculations p. 14-18, letter p. 20, rigidity p. 24-31).
%
% NOTATION fixed once, for three hundred and fourteen pages:
%   Pic_{X/S} the relative Picard functor (the fppf sheaf) and the scheme when
%     it is representable; Pic(X) the absolute group. The pages underline the
%     functor once or twice; the underlines carry nothing the subscript does
%     not, and are dropped.
%   A^vee the dual abelian scheme, always. The pages write A-hat, A', A^*, B^*.
%   G^0 the neutral component, pi_0(G) = G/G^0, G^tau, G[n] the n-torsion.
%     The pages write G_0, Gamma, Pi_0, {}_nG.
%   "lisse" for his "simple"; "schéma" for his "préschéma", with "séparé"
%     added wherever his "schéma" carried it.
%   S^h the henselisation: page 193 writes it S-tilde, which page 188 had
%     already used for the blow-up.
%   Two notions share the word « rigide » and are kept apart throughout: the
%     rigidity LEMMA (morphisms from a product, p. 24-40) and the RIGID GROUP
%     (prime-to-p torsion dense, p. 42-57). A third use, « rigidifié »
%     (a line bundle trivialised along a section), is always written so.
%
% Substantive departures from the pages, each footnoted where it occurs:
%   - p. 3 calls the cokernel Q a kernel;
%   - p. 5: the local Picard scheme of the cone over an elliptic curve E has
%     Q^0 = E and Q/Q^0 = Z/d, d the degree (as p. 187 and 192 compute);
%   - p. 17-18: the products « i != j » run over unordered pairs;
%   - p. 24: the additive decomposition needs G commutative, and X_1 needs
%     connected fibres; p. 28 needs G separated (supplied);
%   - p. 43: multiplication by n is étale, but surjective only on G^0 (Z);
%   - p. 67: the inequalities a), b) must be strict;
%   - p. 69: b_1 = b_3 = 2 dim Pic, hence b_2 = E - 2 + 4p;
%   - p. 96: the vanishing conditions bear on O_Y(-m), not O_Y(m);
%   - p. 98-99: Cor. 3 restricts to Y, not X; the mu_p line is a slip;
%     p. 99 means Pic^tau (not Pic) is proper;
%   - p. 101: U' = f^{-1}(U);
%   - p. 104-111: the Seshadri argument treats D effective;
%   - p. 137: the rigidification needs the inverse of g^*L;
%   - p. 150-154: Pic^0 abelian means ordinary multiple points (delta = r - 1)
%     on a tree, not only ordinary double points;
%   - p. 157: E of rank N+1 for P^N;
%   - p. 163: the last term of the tangent sequence is H^2(O_X);
%   - p. 169: the N.B. on singular loci of codim >= 2 needs S_2;
%   - p. 182: Div^L = P(E) needs geometrically integral fibres, not primary;
%   - p. 191: the condition fails for plane curves exactly when d >= 4;
%   - p. 195: the inequality is sigma(s) >= sigma(s_i), as the Lemma proves;
%   - p. 205: Gamma^tau is the union of the finite subgroups, not finite;
%   - p. 207: the criterion needs an openness hypothesis (A^1 example);
%   - p. 214: L_{D_1}|Y_1 has degree -1;
%   - p. 216-218: « Lemme f » (non-representability) is superseded by the
%     non-separated scheme of p. 218, which represents Pic;
%   - p. 247: Cor. 1.2 needs f proper;
%   - p. 269: the « = Z^I » is dropped;
%   - p. 274: the bracket defines tau-equivalence; it agrees with numerical
%     equivalence by the Matsusaka-Kleiman theorem;
%   - p. 284: h^1 = dim A + dim B holds in residue characteristic 2;
%   - p. 288: condition (iii) must read p F_1 in G_1, and the uniqueness of the
%     subgroup of order p holds in a supersingular elliptic factor;
%   - p. 303: the N.B. sequence is read as one of Picard groups.
%
% Announced and never made, recorded as absences: « Problem B » (p. 73, never
% defined); the families promised on p. 286-287; the proof of Lemme f (struck,
% then contradicted); the end of the theorem of p. 123; Corollaire 1.6
% (p. 255); the construction of the remark of p. 50.
%
% Not restated: the typescripts, reports and letters by others that the
% transcriptions note without reproducing (an IHES letter, Bourbaki drafts and
% « Tribu » reports, English typescripts on theta functions and analytic
% surfaces, a Koszul-complex exposé, a lim^(p) typescript); and the printed EGA
% proof sheets of pp. 153, 158, 160, whose pencil corrections the
% transcription does not attribute firmly.
%
% Thin readings, said where they occur: pages 24-27, 29-31 and 36-40 (fast
% pencil), batch 3 (pp. 41-60), batch 9 (pp. 161-180), batch 11
% (pp. 201-220), and the connective prose of batches 13-14.
%
% Pass: Opus 5.5 (claude-opus-5-5), 2026-10-04 — first pass, unchecked against
% the pages by a human.
\documentclass[11pt,a4paper]{article}
\input{../preamble/grothendieck}

\folder{57}
\pages{1}{314}
\dating{1962-[vers 1968]}
\watermark{Édition de démonstration\\interprétation personnelle de l'œuvre}
\foldertitle{Picard : tapuscrit annoté (s.d.), lettres (s.d., 1962) — lecture modernisée du dossier entier}

\begin{document}

\section*{Résumé}

\begin{resume}
Trois cent quatorze feuillets, pour l'essentiel de 1962, tournent autour d'un
seul objet : le \emph{schéma de Picard}, c'est-à-dire l'espace dont les points
sont les fibrés en droites d'une variété.

Un fibré en droites, c'est une façon de faire passer au-dessus de chaque point
une droite, et de recoller ces droites. Au-dessus d'un cercle il y en a deux :
le cylindre, où l'on recolle tout droit, et le ruban de Möbius, où l'on
recolle avec une torsion. Deux torsions se défont : les fibrés forment un
groupe, ici à deux éléments. Sur une courbe algébrique ce groupe, le groupe de
Picard, est beaucoup plus gros, et c'est la découverte d'Abel et de Jacobi
qu'il n'est pas un simple ensemble : ses éléments de degré zéro forment
eux-mêmes une variété, la jacobienne, qu'on peut parcourir continûment comme
on parcourt un tore. Le dossier pose la question pour une variété quelconque,
et surtout pour une \emph{famille} de variétés : existe-t-il un espace qui
paramètre tous les fibrés en droites, de façon qu'une famille de fibrés soit la
même chose qu'une application vers cet espace ?

Grothendieck sait depuis l'hiver 1961--1962 construire cet espace pour les
familles projectives. Ces feuillets sont ce qui vient après : l'atelier d'un
traité, dont la page 65 dresse le plan en douze chapitres avec un graphe de
dépendances, et qui ne fut jamais rédigé sous cette forme. On y trouve ce que
le plan demande. Des théorèmes d'existence qui ne supposent pas la variété
projective, obtenus en \emph{recollant} le fibré cherché à partir d'un objet
plus simple. Le théorème de finitude : à torsion près, les fibrés invisibles
aux nombres d'intersection ne forment qu'une famille bornée. La séparation,
c'est-à-dire l'unicité des limites. Le schéma d'Albanese, dual du précédent.
Les « lois de rigidité », qui disent qu'une application d'un produit de
variétés complètes vers un groupe est une somme d'applications des facteurs,
et d'où sort le théorème du cube.

La séparation est la question la plus parlante, et un exemple des pages 213 à
218 la rend visible. Une conique lisse dégénère, quand un paramètre $t$ tend
vers $0$, en deux droites qui se coupent. Sur la conique, un fibré n'a qu'un
invariant, son degré $n$ ; sur la paire de droites il en a deux, les degrés
$(p, q)$ sur chacune, avec $p + q = n$. Un fibré de degré $n$ sur la conique
générale a donc une infinité de limites possibles, et l'espace de Picard de la
famille est une droite dont l'origine a été dédoublée une infinité de fois.
L'espace existe, mais une suite y peut converger vers plusieurs points à la
fois : c'est ce qu'on appelle un schéma non séparé.

Le dossier s'ouvre et revient sans cesse sur une variante locale. Au lieu d'une
variété complète, on regarde un point singulier, et le groupe des fibrés en
droites au voisinage épointé de ce point. Pour le sommet du cône construit sur
une courbe elliptique, ce groupe est, à un groupe fini près, la courbe
elliptique elle-même : un invariant d'une singularité qui est une variété. Une
lettre à Hironaka de juillet 1962 compare ce que voient de ce groupe l'algèbre,
la géométrie formelle et la géométrie analytique complexe, et un brouillon
d'article la suit.

Ce qu'on cherchera ensuite sous des noms modernes : les exposés de Grothendieck
au Séminaire Bourbaki sur les schémas de Picard (1962) et l'exposé XIII de
SGA 6 sur la finitude ; le lemme de rigidité et le théorème du cube dans le
livre de Mumford sur les variétés abéliennes ; la représentabilité par des
espaces algébriques (Artin) ; la spécialisation du foncteur de Picard
(Raynaud) ; le schéma de Picard local (Boutot) ; l'extension vectorielle
universelle d'un schéma abélien et les fibrés à connexion ; la stratification
de l'espace des courbes stables par le graphe dual.

Les lettres disent autre chose encore, qui ne demande aucun bagage. Une théorie
s'y construit en public, par des questions données à d'autres. À Murre, en
juillet 1962, il expose un problème d'existence qu'il n'a pas su résoudre et
conclut : « Thus I now wrote you in the hope you may be interested to have a
try on this problem. As a general fact, our knowledge of non projective
existence theorems is exceedingly poor, and I hope this will change
eventually. » Dans la même correspondance il s'excuse de notes « very
sketchy », d'une écriture qu'il qualifie lui-même de « wretched », et suppose
que les notes de son correspondant valent mieux que les siennes ; le tapuscrit
sur le Picard local présente sa construction comme « the only one I could think
of, indeed! ». Le dossier se termine sur une liste de questions ouvertes,
annotée de « Oui » et de « Non » au fil des ans, où l'on voit une théorie
avancer en recensant d'abord ce qu'elle ne sait pas.

\keywords{Picard scheme, Picard functor, local Picard scheme, Greenberg functor, theorem of the cube, theorem of the square, rigidity lemma, abelian scheme, unramified functor, divisorial correspondence, quotient by a flat equivalence relation, torsion-dense group scheme, Albanese scheme, dual abelian scheme, Pic-tau, numerical equivalence, Neron-Severi group, Matsusaka theorem, Igusa inequality, flattening stratification, universal vector extension, line bundle with connection, Atiyah class, Grothendieck-Lefschetz theorem, Artin-Schreier sequence, Weil divisor and Cartier divisor, units of a local ring, effective descent, Weil restriction, generalized Jacobian, toric rank, dual graph of a nodal curve, stable curve, Brauer-Severi scheme, PGL-torsor, pseudo-morphism, fppf sheafification, relative effective divisor, cohomological flatness, class group of a cone, formal completion, henselization, component group, neutral component, constructible set, non-separated Picard scheme, Mittag-Leffler condition, formal functions theorem, exponential sequence, GAGA, resolution of singularities, Neron-Severi scheme, polarization, ample cone, Igusa surface, Hodge asymmetry, Tate module, Pontryagin duality, Kodaira-Spencer class, pinching, non-projective proper scheme}
\end{resume}

\section*{Le fil du dossier, et l'ordre des feuillets}

Le carton n'est pas un texte : c'est une liasse de chemises, chacune ouverte
par une feuille de garde où Grothendieck a écrit un titre, et l'ordre des
chemises n'est pas celui de la pensée. Le tapuscrit sur le Picard local, qui
ouvre le dossier, cite SGA~6 et les travaux d'Artin, et il est sans doute le
plus tardif ; les deux lettres datées sont de juillet 1962 ; le commentaire des
pages 70 et 71 suppose les travaux de Kleiman. Voici le fil, et les conventions
qui valent pour tout ce qui suit.

\subsection*{Les stations}

\begin{itemize}
  \item \textbf{Pages 2 à 12} --- le tapuscrit « Local Picard schemes » : le
    groupe de Picard du spectre épointé d'un anneau local complet, vu comme les
    points d'un schéma en groupes $Q$ sur le corps résiduel.
  \item \textbf{Pages 14 à 40} --- carré, cube, rigidité : les calculs du
    théorème du cube, une lettre non datée à Murre, puis le lemme de rigidité et
    ses conséquences pour les schémas abéliens.
  \item \textbf{Pages 42 à 63} --- les groupes rigides (ceux qu'engendre leur
    torsion première à $p$), puis le schéma d'Albanese comme dual de $\mathrm{Pic}$.
  \item \textbf{Pages 65 à 78} --- le plan du traité, le théorème de finitude
    (Matsusaka, Igusa, puis un commentaire sur Kleiman) et la lettre à Murre du
    18 juillet 1962, en deux états.
  \item \textbf{Pages 80 à 102} --- l'extension du dual d'une variété abélienne
    par les fibrés à connexion ; les critères d'équivalence par section
    hyperplane.
  \item \textbf{Pages 104 à 139} --- les outils de l'existence : le théorème de
    Seshadri, la séparation (unités et diviseurs à support dans les fibres), les
    théorèmes d'existence relatifs par descente.
  \item \textbf{Pages 141 à 182} --- les courbes, les plongements projectifs, les
    pseudo-morphismes, et la définition du foncteur de Picard.
  \item \textbf{Pages 184 à 223} --- la structure locale : le Picard des cônes et
    de leurs anneaux locaux, la composante neutre et le sous-groupe $G^\tau$, la
    partie constructible, et deux exemples où le schéma de Picard se dégrade.
  \item \textbf{Pages 226 à 271} --- la lettre à Hironaka du 6 juillet 1962 et
    le brouillon d'article qui la suit : $\mathrm{Pic}$ en géométrie formelle et
    analytique.
  \item \textbf{Pages 274 à 314} --- schémas de Néron--Severi et polarisations,
    exemples et contre-exemples, et la liste des questions ouvertes.
\end{itemize}

\subsection*{La notation}

On écrit $\mathrm{Pic}(X)$ pour le groupe des classes de faisceaux inversibles
sur $X$ et $\mathrm{Pic}_{X/S}$ pour le foncteur de Picard relatif, faisceau
pour la topologie fidèlement plate, ainsi que pour le schéma qui le représente
quand il existe. Les pages soulignent le foncteur une ou deux fois selon les
feuillets ; on ne garde pas cette distinction, que l'indice porte déjà.
$\mathrm{Pic}^0_{X/S}$ et $\mathrm{Pic}^\tau_{X/S}$ sont la composante neutre
et le sous-groupe des éléments dont un multiple est dans $\mathrm{Pic}^0$ ;
$\mathrm{NS}$ est le groupe de Néron--Severi, que les pages écrivent aussi
« NérSév », « N.S. » ou « Nér ».

Le dual d'un schéma abélien $A$ est noté $A^\vee$ partout ; les pages écrivent
$\hat{A}$, $A'$, $A^*$, selon les feuillets. Pour un schéma en groupes $G$, on
note $G^0$ la composante neutre, $\pi_0(G) = G/G^0$ le groupe des composantes,
et $G[n]$ le noyau de la multiplication par $n$ ; les pages écrivent $G_0$,
$\Gamma$ ou $\underline{\Pi}_0$, et ${}_nG$. Le mot « simple » des pages se dit
aujourd'hui « lisse » ; leur « préschéma » est notre « schéma », et leur
« schéma » notre « schéma séparé ».

\subsection*{Trois rigidités}

Le mot « rigide » sert dans le carton à trois choses qui n'ont rien à voir, et
les confondre rend les pages 24 à 57 inintelligibles. Le \emph{lemme de
rigidité} (pages 24 à 40) dit qu'un morphisme d'un produit $X_1 \times X_2$ de
variétés complètes vers un groupe se décompose ; c'est l'énoncé dont sortent le
théorème du cube et la commutativité des schémas abéliens. Un \emph{groupe
rigide} (pages 42 à 57) est un schéma en groupes commutatif dont la torsion
première à la caractéristique est dense ; c'est une propriété d'un groupe, et
elle sert à descendre des sous-groupes. Un faisceau inversible
\emph{rigidifié} le long d'une section est un faisceau muni d'une
trivialisation sur cette section ; c'est l'outil technique de la définition du
foncteur de Picard (pages 125 et 172). On garde les trois mots distincts.

\subsection*{Ce que le dossier annonce et ne fait pas}

La lettre à Murre des pages 72 à 78 nomme deux problèmes d'existence, A et B,
et ne définit que le premier. Les pages 286 et 287 annoncent quatre familles de
contre-exemples et n'en construisent aucune. Le « Lemme f » de la page 216
affirme que le foncteur de Picard d'un exemple n'est pas représentable ; sa
démonstration est barrée, et la page 218 montre au contraire qu'il l'est, par
un schéma non séparé. Le théorème de la page 123 s'arrête au milieu de son
énoncé. Ces absences sont dites à leur place.

\subsection*{Ce qui n'est pas lu ici}

Le carton contient des feuillets qui ne sont pas de Grothendieck ou ne portent
pas sa main : des tapuscrits d'autres auteurs, des rapports et des projets de
rédaction, des épreuves des \emph{Éléments}, une lettre administrative. Les
transcriptions les signalent sans les reproduire, et cette lecture ne les
restitue pas davantage. Plusieurs liasses sont écrites au crayon, très vite :
les pages 24 à 40, les pages 41 à 60, 161 à 180 et 201 à 220, et la prose de
liaison des pages 241 à 271. On n'en restitue que ce que la transcription
porte, en disant chaque fois où la lecture est mince.

\pagerange{2}{12}
\section*{Le schéma de Picard local (pages 2 à 12)}

Le tapuscrit est en anglais, à la première personne, et annoté à l'encre de sa
main.\footnote{Les numéros d'équation (1) à (3) sont ajoutés à la main ; le
numéro tapé « (34) » est corrigé en « (24) ». Le texte cite SGA~6, XII, la
représentabilité de Murre et « Artin's work » : il est postérieur aux autres
pièces du dossier, ce qui s'accorde avec la datation « vers 1968 » de
l'inventaire.} Il pose une question et esquisse une construction.

\subsection*{La question}

Soient $A$ un anneau local noethérien complet, de corps résiduel $k$,
$S = \operatorname{Spec}(A)$, et $U = S - \{s\}$ le spectre épointé. On
voudrait voir $\mathrm{Pic}(U)$ comme le groupe $Q(k)$ des points rationnels
d'un schéma en groupes $Q$ sur $k$, le \emph{schéma de Picard local}.

On suppose d'abord $U$ régulier, et qu'il existe une résolution qui ne touche
pas $U$ : un morphisme propre $f \colon X \to S$, induisant un isomorphisme
$f^{-1}(U) \simeq U$, avec $X$ régulier et égal à l'adhérence de $U$. Soit
$X_0$ la fibre spéciale, de sorte que $U = X - X_0$. Comme $X$ est régulier,
diviseurs de Weil et de Cartier coïncident, et l'on a une suite exacte
\[
D \longrightarrow \mathrm{Pic}(X) \longrightarrow \mathrm{Pic}(U) \longrightarrow 0 ,
\]
où $D \simeq \mathbf{Z}^I$ est le groupe des diviseurs à support dans $X_0$,
$I$ l'ensemble des composantes irréductibles de $X_0$. Comme $X$ est propre sur
l'anneau complet $A$, le théorème d'existence pour les faisceaux cohérents donne
\[
\mathrm{Pic}(X) \xrightarrow{\ \sim\ } \varprojlim_n \mathrm{Pic}(X_n), \qquad
X_n = X \times_S \operatorname{Spec}(A/\mathfrak{m}^{n+1}) .
\]

Supposons trouvés des schémas en groupes $P_n$ sur $k$, des isomorphismes
compatibles $\mathrm{Pic}(X_n) \simeq P_n(k)$, et des morphismes de transition
$P_m \to P_n$. Le pro-groupe $P = (P_n)$, ou sa limite si les transitions sont
affines, vérifie $\mathrm{Pic}(X) \simeq P(k)$, et l'application
$D \to \mathrm{Pic}(X)$ définit un morphisme $D_k \to P$ du groupe constant
$D_k$. Soit $Q$ son \emph{conoyau}.\footnote{Le tapuscrit écrit « let $Q$ be
this kernel », pour le conoyau qu'il vient de nommer.} Si $D'$ désigne l'image
de $D$ dans $P$, on a $\mathrm{Pic}(U) = P(k)/D'$, d'où une injection
\[
\mathrm{Pic}(U) \hookrightarrow Q(k)
\]
dont l'obstruction à la surjectivité est dans $H^1(k, D'_k)$. Elle s'annule si
$k$ est séparablement clos, ou si $D'$ est sans torsion --- par exemple si
$D_k \to P$ est injectif --- car un homomorphisme continu d'un groupe profini
vers un groupe discret sans torsion est nul. $Q$ est donc le bon
candidat.\footnote{« If $k$ is sep.\ closed, or if » est ajouté à l'encre
devant « $D'$ is torsion-free », où le tapuscrit écrit « is » pour « if ».}

\subsection*{La construction, quand $A$ a un corps de représentants}

Supposons choisi un relèvement $k \to A$ du corps résiduel. Le candidat
naturel, « the only one I could think of, indeed! », est
\[
P_n = \mathrm{Pic}_{X_n/k} ,
\]
représentable, d'après Murre, par un schéma localement de type fini sur $k$.
L'injection $\mathrm{Pic}(X_n) \hookrightarrow P_n(k)$ n'est bijective en
général que si $k$ est séparablement clos : l'obstruction est dans le groupe de
Brauer de l'anneau artinien $B_n = H^0(X_n, \mathcal{O}_{X_n})$, qui est local
si $X$ est connexe, de sorte que $\mathrm{Br}(B_n) = \mathrm{Br}(k_n)$, $k_n$
son corps résiduel. Ces obstructions sont compatibles et proviennent du groupe
de Brauer du corps résiduel $k'$ de
$H^0(X, \mathcal{O}_X) = \varprojlim H^0(X_n, \mathcal{O}_{X_n})$, qui est la
normalisée de $A$ ; si $A$ est normal, $k' = k$.

On décrit alors $Q(k)$ par descente galoisienne : si $A'$ est le complété du
hensélisé strict de $A$, $U'$ son spectre épointé, et $\pi = \mathrm{Gal}(k^s/k)$,
\[
Q(k) \simeq \mathrm{Pic}(U')^{\pi} .
\]
C'est, dit le tapuscrit, la situation familière pour les schémas de Picard
globaux. Une correction suit : pour obtenir cette formule quand $k$ n'est pas
séparablement clos, il faut remplacer le groupe constant $D_k$ par le groupe
constant tordu $\underline{D}$, image directe du faisceau constant $\mathbf{Z}$
par $\operatorname{Spec}(\prod_i k_i) \to \operatorname{Spec}(k)$, où $k_i$ est
la clôture séparable de $k$ dans le corps des fonctions de la composante
$X_0^i$.\footnote{Une première rédaction, encadrée et barrée à la main,
définissait $\underline{D}$ par la plus grande sous-algèbre séparable de
$H^0(X_0, \mathcal{O}_{X_0})$.}

Deux réserves sont notées. L'indépendance de $Q$ vis-à-vis de la résolution
devrait suivre de l'existence de résolutions dominant deux résolutions données,
ce qu'on a en caractéristique $0$ (Hironaka) ou en dimension $2$. Et la
structure algébrique dépend essentiellement du relèvement choisi du corps
résiduel ; l'exemple à étudier est le cône sur une courbe elliptique, pour
lequel le tapuscrit annonce que $Q$ est la courbe elliptique.\footnote{Plus
précisément, et comme le calculent les pages 184 à 192 : si la courbe
elliptique $E$ est plongée par un faisceau de degré $d$, on a
$\mathrm{Pic}(U) \simeq \mathrm{Pic}(E)/\mathbf{Z}\,\xi$, extension de
$\mathbf{Z}/d\mathbf{Z}$ par $E(k)$ ; c'est $Q^0$ qui est isomorphe à $E$, et
$Q/Q^0 \simeq \mathbf{Z}/d\mathbf{Z}$.}

\subsection*{La structure de $P$ et de $Q$}

Les transitions $P_{n+1} \to P_n$ sont affines (SGA~6, XII), donc
$P = \varprojlim P_n$ est un vrai schéma en groupes, d'algèbre de Lie
\[
\mathrm{Lie}(P) \simeq \varprojlim H^1(X_n, \mathcal{O}_{X_n}) \simeq H^1(X, \mathcal{O}_X) ,
\]
de dimension finie puisque $R^1 f_* \mathcal{O}_X$ est cohérent et concentré au
point fermé. Quitte à remplacer $X$ par un modèle qui le domine, on peut
espérer un idéal inversible $J$ de même support que $X_0$ tel que
$L = J/J^2$ soit ample sur $Z = V(J)$. Le système $(Z_n = V(J^{n+1}))$ définit
le même pro-objet, et pour $n$ grand les transitions
$\mathrm{Pic}_{Z_{n+1}/k} \to \mathrm{Pic}_{Z_n/k}$ sont des isomorphismes,
leurs noyau et obstruction étant contrôlés par $H^1$ et $H^2$ de
$L^{\otimes (n+1)}$ sur $Z$, qui s'annulent par le théorème de Serre. Le
pro-groupe est alors essentiellement constant et $P$ est localement de type
fini sur $k$. C'est certain si $A$ est normal de dimension $2$, par la
négativité de la matrice d'intersection des composantes de $X_0$ (Du Val,
Mumford).

On voudrait que $\underline{D} \to P$ soit une immersion fermée, ce qui revient
à l'injectivité de $D \to \mathrm{NS}(X_0)$ (le noyau de
$\mathrm{NS}(Z_n) \to \mathrm{NS}(X_0)$ étant fini). Grothendieck s'y attend
toujours ; en dimension $2$, c'est encore la négativité. Alors
$Q = P/\underline{D}$ est un schéma en groupes localement de type fini sur
$k$, avec
\[
P^0 \xrightarrow{\ \sim\ } Q^0, \qquad
Q/Q^0 \simeq \mathrm{NS}_{Z_n/k}/\underline{D} \ (n \gg 0), \qquad
\mathrm{Lie}(Q) \simeq H^1(X, \mathcal{O}_X) .
\]
La restriction $H^1(X, \mathcal{O}_X) \to H^1(U, \mathcal{O}_U) = H^2_{\mathfrak{m}}(A)$
est injective : avec $i \colon U \to X$, morphisme affine, le faisceau
$i_*\mathcal{O}_U/\mathcal{O}_X$ a une suite de composition de quotients
$L^{\otimes -n}$, $n \geq 1$, sans sections puisque $L$ est ample sur $Z$, de
dimension $\geq 1$. L'algèbre de Lie du schéma de Picard local se plonge donc
dans $H^2_{\mathfrak{m}}(A)$ ; si $A$ est de Cohen--Macaulay de dimension
$\geq 3$, elle est nulle, $Q$ est étale, et $\mathrm{Pic}(U)$ est dénombrable
dès que $k$ l'est.

\subsection*{Le foncteur des points}

Reste à décrire $Q$ comme foncteur sur les $k$-algèbres, ce que la
construction, liée à une résolution particulière, ne fait pas directement.
Pour $T$ affine sur $k$, on a
\[
P(T) = \varprojlim P_n(T) \hookleftarrow \varprojlim_n \mathrm{Pic}(X_{n,T})/\mathrm{Pic}(T) ,
\]
et $P$ est le faisceau étale associé au membre de droite, l'obstruction étant
dans $H^2(T_{\text{ét}}, \mathbf{G}_m)$ (pour $A$ normal).\footnote{La note
marginale précise que l'affinité de $T$ sert à annuler
$H^1(T, \mathcal{O}_T)$ et $H^2(T, \mathcal{O}_T)$ ; les exposants tapés 2 et
3 y sont corrigés à la main en 1 et 2.} Comme $P$ est localement de
présentation finie, il suffit de le connaître sur les $k$-algèbres de type
fini $B$. Posant $C = A \hat{\otimes}_k B = \varprojlim A_n \otimes_k B$,
anneau adique noethérien, et $X \hat{\otimes}_k B = X \otimes_A C$, propre sur
$C$, le théorème d'existence donne
$\mathrm{Pic}(X \hat{\otimes}_k B) \simeq \varprojlim \mathrm{Pic}(X_{n,B})$ :
sur les arguments noethériens, $P$ est le faisceau étale associé au préfaisceau
séparé $B \mapsto \mathrm{Pic}(X \hat{\otimes}_k B)/\mathrm{Pic}(B)$, et $Q$
celui associé à
\[
B \longmapsto \mathrm{Pic}(X \hat{\otimes}_k B)/\bigl(\mathrm{Pic}(B) + \underline{D}(B)\bigr) .
\]
Deux obstructions successives séparent $Q(B)$ de ce quotient : l'une dans
$H^1(\operatorname{Spec} B, \underline{D})$, nulle si $\underline{D}$ est
constant et $B$ normal ; l'autre dans $\mathrm{Br}(B)$, nulle si $X_0$ porte un
zéro-cycle de degré $1$.

Si les composantes de $X_0$ sont géométriquement irréductibles et $B$
géométriquement régulier sur $k$, alors $X \hat{\otimes}_k B$ est régulier, et
le quotient par $\underline{D}(B)$ est $\mathrm{Pic}(U \hat{\otimes}_k B)$, où
$U \hat{\otimes}_k B$ est l'image inverse de $U$ dans
$\operatorname{Spec}(A \hat{\otimes}_k B)$ --- objet qui ne dépend plus de la
résolution. Sur ces arguments, $Q$ est le faisceau étale associé à
$B \mapsto \mathrm{Pic}(U \hat{\otimes}_k B)/\mathrm{Pic}(B)$, à une
obstruction de Brauer près : une description de $Q(B)$ par de vrais groupes de
Picard locaux, sous des hypothèses très restrictives sur $B$.

\subsection*{Sans corps de représentants, sans régularité}

En inégales caractéristiques, le corps résiduel $k$ étant parfait de
caractéristique $p$, tout repose sur le problème suivant : pour $A$ local
artinien de corps résiduel $k$ et $X$ propre sur $A$, construire
« naturellement » un schéma en groupes $P$ localement de type fini sur $k$ et
un plongement $\mathrm{Pic}(X) \hookrightarrow P(k)$, d'obstruction dans
$\mathrm{Br}\, H^0(X, \mathcal{O}_X)$. Le candidat proposé passe par le schéma
en anneaux $\underline{A}$ sur $k$ de Greenberg, tel que
$A = \underline{A}(k)$ : on pose $X_B = X \otimes_A \underline{A}(B)$ et l'on
considère le faisceau fppf associé à $B \mapsto \mathrm{Pic}(X_B)$, dont les
critères d'Artin devraient décider la représentabilité.

Enfin, sans supposer $U$ régulier, Grothendieck conjecture l'existence d'un
schéma de Picard local localement de type fini, et propose deux voies : un
morphisme propre surjectif $X \to S$, isomorphe au-dessus de $U$, tel que
$\mathrm{Pic}(X) \to \mathrm{Pic}(U)$ soit surjectif ; ou une résolution qui ne
soit plus un isomorphisme au-dessus de $U$, avec des données de descente. Le
test proposé : pour $A$ de Cohen--Macaulay de dimension $\geq 3$ (plus
généralement $H^2_{\mathfrak{m}}(A) = 0$), $\mathrm{Pic}(U)$ devrait être
dénombrable et invariant par extension séparable du corps résiduel, ce qui
exprimerait que le schéma de Picard local est discret.

\pagerange{14}{18}
\section*{Carré, cube (pages 13 à 18)}

La feuille de garde porte « Carré, cube, rigidité… ». Les feuillets qui
suivent sont des calculs, sans prose, dont le cadre est celui d'un schéma
abélien $A$ sur $S$ et d'un faisceau inversible $L$ sur $A$. Pour des points
$x, y, \ldots$ de $A$ à valeurs dans un $S$-schéma $T$, on note
$L_x = x^*L$, faisceau inversible sur $T$, et $L_\emptyset = e^*L$ ; les
égalités qui suivent sont des isomorphismes canoniques de faisceaux inversibles
sur $T$.

\subsection*{Le théorème du cube en $n$ variables}

Posons $(\delta L)_{x,y} = L_{x+y} \otimes L_x^{-1} \otimes L_y^{-1}$,
symétrique en $x, y$, de sorte que $(\delta L)_{0,0} = L_\emptyset^{-1}$. Le
théorème du cube s'écrit
\[
L_{x+y+z} \simeq L_{x+y} \otimes L_{y+z} \otimes L_{x+z} \otimes L_x^{-1} \otimes L_y^{-1} \otimes L_z^{-1} \otimes L_\emptyset ,
\]
et la page 18 en tire par récurrence, pour $n \geq 2$,
\[
L_{x_1 + \cdots + x_n} \simeq \Bigl(\bigotimes_{i<j} L_{x_i + x_j}\Bigr) \otimes \Bigl(\bigotimes_i L_{x_i}\Bigr)^{-(n-2)} \otimes L_\emptyset^{\alpha_n},
\qquad \alpha_n = \frac{(n-1)(n-2)}{2} ,
\]
le produit portant sur les paires non ordonnées.\footnote{Les pages 17 et 18
écrivent $\prod_{i \neq j}$ ; les formules développées qu'elles encadrent pour
$n = 4$ et $n = 5$ ne prennent chaque paire qu'une fois. Les bornes des
produits de la récurrence de la page 18 sont lues sous réserve.} L'exposant se
trouve en spécialisant tous les $x_i$ en $0$ : on doit avoir
$1 = -\frac{n(n-1)}{2} + n + \alpha_n$ (page 16) ; la page 17 résout la
récurrence $\alpha_{n+1} = 2\alpha_n - \alpha_{n-1} + 1$, essaie d'abord une
solution linéaire, la barre, et trouve la quadratique
$\alpha_n = \frac12(n^2 - 3n + 2)$, encadrée ; on vérifie
$\alpha_3 = 1$, $\alpha_4 = 3$, $\alpha_5 = 6$, comme dans les formules
développées des pages 17 et 18.

\subsection*{La multiplication par $n$}

En faisant $x_1 = \cdots = x_n = x$, on obtient, à un faisceau venu de la base
près --- exactement si $L$ est rigidifié le long de la section unité ---, la
formule encadrée de la page 14 :
\[
n^*L \simeq L^{\otimes n} \otimes (\delta L)^{\otimes \frac{n(n-1)}{2}},
\qquad \delta L = 2^*L \otimes L^{-2} ,
\]
et, en écrivant $\delta L = L^{\otimes 2} \otimes \delta' L$ avec
$\delta' L = 2^*L \otimes L^{-4}$, la seconde formule encadrée\footnote{L'exposant
$n^2$ de cette formule est noyé sous une surcharge et lu sous réserve par la
transcription ; il est confirmé par le calcul
$n + 2 \cdot \frac{n(n-1)}{2} = n^2$, que la page écrit à côté.}
\[
n^*L \simeq L^{\otimes n^2} \otimes (\delta' L)^{\otimes \frac{n(n-1)}{2}} .
\]
Comme
$2^*L \simeq L^{\otimes 3} \otimes (-1)^*L$, on a
$\delta L = L \otimes (-1)^*L$, et l'on retrouve la forme familière
$n^*L \simeq L^{\otimes \frac{n^2+n}{2}} \otimes ((-1)^*L)^{\otimes \frac{n^2-n}{2}}$ ;
$\delta' L = (-1)^*L \otimes L^{-1}$ est nul si $L$ est symétrique.

\subsection*{Le Picard d'un produit}

La page 14 encadre encore, pour des $X_i$ propres et plats sur $S$ avec
$f_{i*}\mathcal{O}_{X_i} \simeq \mathcal{O}_S$ universellement, la suite exacte
de faisceaux (fidèlement plats quasi-compacts)
\[
0 \to \prod_i \mathrm{Pic}_{X_i/S} \longrightarrow \mathrm{Pic}_{(\prod X_i)/S} \longrightarrow \prod_{\{i,j\}} \mathrm{Corr}_S(X_i, X_j) \to 0 ,
\]
$\{i, j\}$ parcourant les parties à deux éléments de l'ensemble d'indices,
scindée si les $X_i$ ont des sections sur $S$ : c'est le théorème du cube sous
la forme « il n'y a pas de termes triples ». $\mathrm{Corr}_S(X_i, X_j)$ est
le foncteur des classes de correspondances divisorielles. Elle donne
$\mathrm{Pic}^\tau_{(\prod X_i)/S} \simeq \prod \mathrm{Pic}^\tau_{X_i/S}$
lorsque les groupes de correspondances sont sans torsion, condition que la
page ajoute en travers, soulignée deux fois.\footnote{« sans torsion » est
écrit à côté de l'encadré sans qu'un trait le rattache à un terme précis ;
c'est notre lecture qu'il porte sur les correspondances.}

La page 15 ébauche un foncteur $\mathrm{Pic}_{\mathrm{sp}}(X/Y)$ de classes de
faisceaux inversibles « à sections marquées », à cohomologie supérieure nulle
et dont l'image directe est localement libre de rang
$\chi_{X/Y}(L)$ ; le feuillet est surchargé et la définition ne se lit pas
au-delà. La page 16 ébauche le même calcul que la page 18 en termes de
points, avec un schéma d'incidence des $X_{[ijkl]}$.

\pagerange{19}{22}
\section*{Une lettre à Murre, non datée (pages 19 à 22)}

Lettre dactylographiée en anglais. Murre doit faire un exposé sur les foncteurs
non ramifiés, et pose des questions ; Grothendieck répond en trois points.

\emph{Le passage au quotient.} Soit $f \colon X \to Y$ un morphisme de
$S$-schémas, avec $X$ et $Y$ localement de présentation finie sur $S$, ou $Y$
localement noethérien et $X$ localement de type fini sur $Y$. Si la relation
d'équivalence $R = X \times_Y X$ est plate sur $X$, alors $f$ se factorise en
$X \to Z \to Y$, avec $X \to Z$ fidèlement plat localement de présentation
finie, $Z$ localement de présentation finie sur $Y$, et $Z \to Y$ un
monomorphisme ; autrement dit, le faisceau quotient fpqc $X/R$ est
représentable, et c'est ainsi qu'on le démontre.\footnote{Dans le langage
actuel, le quotient d'un schéma par une relation d'équivalence plate de
présentation finie est d'abord un espace algébrique ; c'est le monomorphisme
vers $Y$, séparé et localement quasi-fini, qui en fait un schéma.}

Il rapporte une application due à Raynaud : sur un anneau de valuation
discrète, si $G$ est un schéma en groupes de type fini et $H$ un sous-groupe
fermé et plat dont le quotient générique $G_t/H_t$ est quasi-affine, alors
$G/H$ est représentable, quasi-affine et plat, et affine si $H$ est invariant ;
en particulier un schéma en groupes (séparé) plat de type fini à fibre
générique affine est affine.\footnote{« Group scheme » est à entendre, en 1962,
au sens de schéma séparé.} Raynaud cherche à lever l'hypothèse de
quasi-affinité et à construire $G/H$ quasi-projectif.

\emph{Le théorème du cube.} Pour la non-ramification du foncteur
$\mathrm{Corr}$, une forme infinitésimale faible du théorème du carré suffit,
démontrée dans des notes manuscrites jointes « on correspondance
classes » --- « I agree the handwriting is wretched and the notes moreover very
sketchy ». Le théorème du cube découle formellement de la séparation de
$\mathrm{Corr}_S(X, Y)$ pour deux des trois facteurs, et des propriétés
formelles du foncteur de Picard, dont la moins triviale est la commutation aux
limites projectives d'anneaux artiniens.

\emph{La séparation de $\mathrm{Corr}_S(X, Y)$.} Elle est à peu près triviale
dès que le foncteur de Picard d'un des facteurs est séparé, ce qui est le cas
si ses fibres géométriques sont intègres. On se ramène au spectre d'un anneau
de valuation discrète, d'uniformisante $t$ : si $L$ est trivial sur la fibre
générique, il est défini par un diviseur de Cartier à support dans la fibre
spéciale $X_0$ ; celle-ci est un diviseur de Cartier ($t = 0$), intègre, donc
le diviseur est un multiple de $X_0$ --- l'argument suppose les fibres
géométriquement normales, donc $X$ normal --- et il est linéairement équivalent
à $0$.

Il est convaincu que $\mathrm{Corr}$ est toujours séparé (sous les hypothèses
habituelles de propreté, de platitude, et $f_*\mathcal{O} = \mathcal{O}$ pour
les deux facteurs). C'est facile par la théorie de la dimension si le foncteur
de Picard d'un facteur est représentable, et cela reste vrai s'il est
seulement « pré-ét-représentable », quotient d'un schéma localement de type fini
par une relation d'équivalence étale (ou quasi-finie et plate) : c'est le cas
si $Y$ est projectif, par la représentation de grands ouverts de
$\mathrm{Pic}_{Y/S}$ comme quotient d'un schéma de plongements de $Y$ dans un
$\mathbf{P}^r$ par le groupe projectif, qui opère librement --- c'est la
construction que reprennent les pages 157 à 161. Sans projectivité, le lemme de
Chow conserve la platitude sur un anneau de valuation discrète mais perd
l'égalité $H^0(X_0, \mathcal{O}_{X_0}) \simeq k(s)$. Une approche topologique
est proposée : montrer que l'image de la spécialisation du groupe fondamental
de la fibre générique géométrique dans celui de la fibre spéciale est d'indice
fini, au moins après abélianisation, ce qu'il croit démontrable par le foncteur
de Picard pour $X$ projectif.

La lettre se clôt sur l'envoi de notes contenant une esquisse de la
représentabilité des foncteurs non ramifiés, « although I do not think the
latter can be of any use to you, as I have a hard time myself to read them » :
les notes prises par Murre lors de leur discussion doivent être plus
détaillées, « the inverse is more plausible ».

\pagerange{24}{40}
\section*{Rigidité et schémas abéliens (pages 23 à 40)}

La chemise porte « Rigidité » et « Groupes V » ; le manuscrit est paginé 1 à 3
par l'auteur aux pages 24, 26 et 28. La main est rapide, et plusieurs
démonstrations ne se lisent que par fragments ; les énoncés passent.

\subsection*{Le lemme de rigidité}

\emph{Théorème (page 24).} Soient $S$ localement noethérien,
$\varphi_i \colon X_i \to S$ ($i = 1, 2$) deux $S$-schémas de type fini, $\varepsilon_1$ une section de $X_1$ sur $S$,
et $p_1, p_2$ les projections de $X_1 \times_S X_2$. On suppose que $X_2$ est
propre et plat sur $S$ avec $\varphi_{2*}\mathcal{O}_{X_2} \simeq \mathcal{O}_S$
universellement, et que $X_1$ est plat sur $S$ à fibres géométriquement
connexes. Alors pour tout schéma en groupes $G$ séparé sur $S$ et tout
$S$-morphisme $f \colon X_1 \times_S X_2 \to G$, on a
\[
f = (f_1 \circ p_1)\cdot(f_2 \circ p_2)
\]
pour des $S$-morphismes $f_i \colon X_i \to G$ uniquement déterminés par
$f_1 \circ \varepsilon_1 = e$, et alors $f_2 = f \circ (\varepsilon_1 \times \mathrm{id}_{X_2})$.\footnote{La
page écrit les hypothèses sous la forme : a) $p_1$ propre et
$p_{1*}\mathcal{O} \simeq \mathcal{O}_{X_1}$ ; b) $p_1$ universellement ouvert,
$X_1$ plat sur $S$ « à fibres géom.\ primaires » (lu sous réserve), avec pour
exemple, entre crochets, $\varphi_2$ propre et plat à $H^0$ trivial et
$\varphi_1$ à fibres géométriquement irréductibles. Nous retenons l'exemple,
qui est ce que la démonstration utilise : la connexité des fibres de $X_1$ est
ce qui fait que chaque composante connexe de $X_1$ rencontre
$\varepsilon_1(S)$. La page écrit aussi « $S$-homomorphismes » pour les $f_i$,
qui ne sont que des morphismes, et $f_2 \colon X_2 \to X$ pour $X_2 \to G$.}

La démonstration (page 25) considère
$f' = f \cdot (f_2 \circ p_2)^{-1}$, trivial sur $\varepsilon_1(S) \times_S X_2$.
Par le lemme 1 ci-dessous, $f'$ se factorise par $p_1$ au-dessus d'un voisinage
ouvert $U_1$ de $\varepsilon_1(S)$ ; par le lemme 2 et l'ouverture universelle
de $p_1$, l'ensemble des points de $X_1$ au-dessus desquels $f'$ est à image
finie est aussi fermé ; il rencontre chaque fibre $X_{1,s}$, connexe, en une
partie ouverte et fermée contenant $\varepsilon_1(s)$, donc c'est tout $X_1$.\footnote{Le bas de la page 25, très
surchargé, ne se lit que par fragments ; la marche que l'on donne est celle que
les lemmes 1 et 2 imposent.}

\emph{Lemme 1 (page 26).} Soient $f \colon X \to Y$ propre, $Y$ localement
noethérien, $f_*\mathcal{O}_X \simeq \mathcal{O}_Y$, $h \colon Z \to Y$ séparé
localement de type fini, $g \colon X \to Z$ un $Y$-morphisme, et $y \in Y$ tel
que $g(X_y)$ soit fini. Il existe un voisinage ouvert $U$ de $y$ et une section
$\sigma$ de $Z$ sur $U$, unique, avec $g = \sigma \circ f$ sur $f^{-1}(U)$.
Démonstration : $X' = g(X)$ est propre sur $Y$ et à fibre finie en $y$, donc
fini au-dessus d'un voisinage ; $g$ se factorise par le $Z' = \operatorname{Spec}(g_*\mathcal{O}_X)$
fini sur $Z$, et
$\mathcal{O}_Y \to (h g'')_*\mathcal{O}_{Z'} \subset f_*\mathcal{O}_X \simeq \mathcal{O}_Y$
montre que $Z' \to Y$ est un isomorphisme ; on prend
$\sigma = g''(hg'')^{-1}$.

\begin{tikzcd}[column sep=small, row sep=small]
 & Z' \arrow[dr, "g''"] & \\
X \arrow[ur, "g'"] \arrow[rr, "g"] \arrow[d, "f"'] & & Z \arrow[dll, "h"] \\
Y & &
\end{tikzcd}

\emph{Corollaire.} L'ensemble des $y$ tels que $g(X_y)$ soit fini est ouvert, et
$g$ se factorise au-dessus de lui par une section unique.

\emph{Lemme 2 (page 27).} Si $f \colon X \to Y$ est de type fini et
universellement ouvert, $h \colon Z \to Y$ localement de type fini et
$g \colon X \to Z$ un $Y$-morphisme, l'ensemble des $y$ tels que $g(X_y)$ soit
fini est fermé.\footnote{La démonstration, par réduction au cas noethérien,
irréductible et dominant, puis au spectre d'un anneau de valuation discrète,
n'est lisible que par mots isolés. Un « corollaire 1 » et une « remarque » de
la même page, qui portent sur la propreté de $X_1 \times_S X_2 \to S$ et sur
$f' = f - f_1 p_1 - f_2 p_2$, ne se lisent pas au-delà de leurs premiers
mots.}

\subsection*{Conséquences pour les schémas abéliens}

\emph{Corollaire (page 24).} Si $G$ est commutatif et $X_2$ a aussi une section
$\varepsilon_2$, $f$ s'écrit de façon unique
$f = f_1 p_1 + f_2 p_2 + \sigma \varphi$, avec $f_i \circ \varepsilon_i = e$,
$\sigma$ une section de $G$ et $\varphi$ la projection sur $S$.\footnote{La
page écrit additivement sans supposer $G$ commutatif ; pour un $G$ quelconque
la décomposition du théorème suffit, à condition de garder l'ordre des
facteurs.}

\emph{Corollaire 2 (page 28).} Soient $A$ un schéma abélien sur $S$ localement
noethérien et $G$ un schéma en groupes localement de type fini et séparé sur
$S$. Tout $S$-morphisme $f \colon A \to G$ s'écrit de façon unique
$f = u \cdot c$, avec $u$ un homomorphisme de groupes et $c = f \circ \varepsilon_A$
« constant ». En particulier, un morphisme qui respecte les sections unités est
un homomorphisme.\footnote{L'hypothèse de séparation, que la page n'écrit pas
ici, est celle du théorème dont le corollaire découle. On applique le lemme au
morphisme $(x, y) \mapsto f(xy)\,f(y)^{-1} f(x)^{-1}$, trivial sur les deux
axes.}

\emph{Corollaire 3.} Un schéma abélien est commutatif, et toute loi de
composition sur $A$ admettant $\varepsilon_A$ pour unité bilatère est la loi
donnée. La commutativité vient du corollaire 2 appliqué à $x \mapsto x^{-1}$ ;
pour la seconde assertion, si $m'$ est une telle loi,
$m'(x, y) - x - y$ est trivial sur les deux axes, donc nul.

\subsection*{La reformulation fonctorielle (pages 29 à 31)}

L'énoncé est repris pour un foncteur en groupes $G$ quelconque sur les
$S$-schémas, au lieu d'un schéma en groupes, et la page 29 conclut que le
foncteur des morphismes pointés $\mathrm{Hom}^\bullet_S(X, G)$ est \emph{non
ramifié} sur $S$.\footnote{Les conditions (i)--(iii) de la page 29 portent sur
la trivialité d'un homomorphisme sur un sous-schéma de même espace sous-jacent,
sur une fibre, puis sur une composante connexe ; elles sont très surchargées et
on n'en retient que la conclusion.} La page 30 donne la forme la plus nette.

\emph{Théorème (page 30).} Soient $f \colon X \to S$ propre et plat avec
$\mathcal{O}_S \simeq f_*\mathcal{O}_X$ universellement, $S$ localement
noethérien, $G$ un foncteur sur les $S$-schémas formellement représentable,
compatible au passage aux complétés et aux limites --- conditions vérifiées si
$G$ est représentable localement de type fini --- et $u \in G(X)$. Soit $U$
l'ensemble des $s$ tels que $u_s$ provienne de $G(k(s))$. Alors $U$ est ouvert,
et il existe un unique $v \in G(U)$ dont l'image dans $G(X_U)$ est
$u|_{X_U}$. La page ajoute, d'une autre encre, que $U$ est aussi fermé sous
une condition de type valuatif --- c'est le cas si $G$ est représentable par un
schéma séparé --- et c'est cette fermeture qu'utilisent les corollaires.\footnote{On
ne lit que « et \emph{fermé} si … (critère valuatif) ». Les conditions b) et c)
sur $G$ ne sont lues que partiellement.}

Les corollaires en tirent : si $S$ est connexe et que $u_s$ se factorise par
$k(s)$ en un point, $u$ se factorise par $S$ ; sous des sections unités, la
trivialité en une fibre entraîne la trivialité ; et la forme de Weil
(corollaire 5), pour $G$ commutatif : un morphisme $u \colon X \times_S Y \to G$, $Y$ à fibres
géométriquement connexes muni d'une section, s'écrit de façon unique
$u = u_1\,\mathrm{pr}_1 + u_2\,\mathrm{pr}_2 + v\,p$.

La page 31 pose, avec la mention « Demander Mumford », la question que la
lettre à Murre reprend : pour $X$, $Y$ propres et plats sur $S$, à $H^0$
trivial universellement, munis de sections, un faisceau inversible sur
$X \times_S Y$ trivialisé sur les deux axes et trivial au-dessus d'un ouvert
non vide est-il trivial ? Autrement dit, $\mathrm{Corr}_S(X, Y)$ est-il
toujours séparé ? La marge répond qu'il suffit que $\mathrm{Pic}_{X/S}$ ou
$\mathrm{Pic}_{Y/S}$ soit pré-ét-représentable, en particulier que $X$ ou $Y$
soit projectif.

\subsection*{Variétés à multiplication, et relèvement des lois de groupe}

La page 32 énonce, puis barre d'un réseau de traits, un théorème : sur $S$
normal, un schéma propre et lisse muni d'une section, dont toutes les fibres
sont des variétés abéliennes, est un schéma abélien. Il est subsumé par le
théorème final de la page 40, qui n'a pas besoin de la normalité. Un
théorème 2, barré lui aussi, redonne la décomposition
$u = u_1\,\mathrm{pr}_1 + u_2\,\mathrm{pr}_2$ d'un morphisme
$X_1 \times_S X_2 \to G$ sous des hypothèses de platitude, avec pour corollaire,
illustré d'un carré de multiplications, qu'un morphisme $u \colon X \to G$ d'un
schéma en groupes plat et lisse à fibres connexes vers un schéma en groupes,
tel que $u(e) = e$, est un homomorphisme.\footnote{Le corollaire est lu mot à
mot sous réserve, sous la rature.}

\emph{Lemme 1 (page 34).} Soit $G$ une variété propre et géométriquement
intègre sur $k$, munie d'un point rationnel $e$, d'une multiplication
admettant $e$ pour unité bilatère, et d'une application $y \mapsto y^\vee$
telle que $x y^\vee = e \iff x = y$. Alors la loi est associative et
commutative, et $G$ est une variété abélienne.\footnote{La page dit
« irréd.\ et \emph{propre} » ; le lemme de rigidité veut la variété
géométriquement intègre, ou que l'on passe à la clôture algébrique.} Pour la
commutativité, $f(a, b) = (ab)(ba)^\vee$ vaut $e$ sur $\{e\} \times G$, donc
partout par rigidité, d'où $ab = ba$ ; pour l'associativité, de même avec
$[(ab)c][a(bc)]^\vee$. L'inverse de $y$ est $y^\vee$, puisque $y y^\vee = e$.

\emph{Lemme 3 (page 36).} Soient $S = \operatorname{Spec}(A)$, $A$ artinien
local, $G$ lisse et propre sur $S$ muni d'une section $e$, et $S_0 \subset S$
défini par un idéal de carré nul $I$. Si $G_0 = G \times_S S_0$ est un schéma
abélien d'unité $e_0$, il existe sur $G$ une unique loi de groupe relevant
celle de $G_0$ et d'unité $e$ ; $G$ est un schéma abélien. L'obstruction au
relèvement de la multiplication $\pi_0$ est dans
\[
H^1\bigl(G_0 \times G_0, \pi_0^*T_{G_0} \otimes I\bigr) = H^1(G_0 \times G_0, \mathcal{O}) \otimes_k (t \otimes_k I) ,
\]
$t$ l'espace tangent en l'origine, puisque le fibré tangent d'un schéma abélien
est trivial. Par Künneth, elle est déterminée par ses restrictions à
$e \times G_0$ et $G_0 \times e$, qui sont les obstructions à relever
l'identité de $G_0$ : elles sont nulles. L'indétermination est dans
$H^0(G_0 \times G_0, \mathcal{O}) \otimes t \otimes I = t \otimes I$, et la
condition $e \cdot e = e$ la fixe.\footnote{La première rédaction, biffée,
introduisait explicitement l'idéal de carré nul ; la lecture définitive du
début est incertaine, et la dernière lettre de « $\mathcal{V} \otimes_k \mathfrak{m}$ »
(notre $t \otimes I$) est une lecture probable. Le cas d'un $A$ artinien
quelconque s'en déduit en filtrant par des idéaux de carré nul.} Plus
généralement, pour un morphisme $X \to G$, l'obstruction est dans
$H^1(X_0, \mathcal{O}) \otimes t \otimes I$ et l'indétermination dans
$H^0(X_0, \mathcal{O}) \otimes t \otimes I$, fixée par l'image d'une section si
$H^0(X_0, \mathcal{O}) = k$ ; le même raisonnement montre alors que $e$ est une
unité bilatère et que la loi relevée est commutative et associative.

La page 40 en tire, par unicité, l'existence sur un anneau local noethérien
complet, puis local par descente, et enfin sur une base connexe : l'ensemble
des points où la fibre est abélienne, déjà ouvert, est aussi fermé, ce qu'on
vérifie sur un anneau de valuation discrète.\footnote{Cette dernière page est
lue avec beaucoup d'incertitude, mot à mot ; l'enchaînement est clair, les
justifications non.}

\emph{Théorème (page 40).} Soient $S$ localement noethérien connexe et $G$ lisse
et propre sur $S$ muni d'une section $e$. Si une fibre est une variété
abélienne d'origine $e$, il existe sur $G$ une unique loi de groupe d'unité $e$,
et elle fait de $G$ un schéma abélien.

C'est le théorème qu'on trouve aujourd'hui dans le livre de Mumford sur la
théorie géométrique des invariants (1965), au chapitre des schémas abéliens.

\pagerange{42}{57}
\section*{Groupes rigides et descente des sous-groupes (pages 41 à 57)}

Manuscrit paginé 4 à 11 par l'auteur ; le début (ses pages 1 à 3) manque. La
main est rapide : les énoncés passent, la prose de liaison est souvent perdue,
et ce qui suit ne restitue que ce que la transcription porte.

\subsection*{La définition}

Un schéma en groupes commutatif $G$ de type fini sur $S$ est dit \emph{rigide}
si (i) il est plat sur $S$ ; (ii) pour tout $s \in S$ et tout nombre premier
$\ell$ distinct de la caractéristique de $k(s)$, la fibre géométrique
$\overline{G}_s$ est l'adhérence schématique de ses points de torsion
$\ell$-primaire ; (iii) une condition de finitude, que la page ne laisse pas
lire.\footnote{La condition (ii) est lue avec des lacunes, mais l'encadré
« ${}_{\ell^k}(\overline{G}_s)$ » et la proposition qui suit en fixent le
sens ; de (iii) on ne lit que « pour tout ouvert $U$ de $S$ … le schéma … est
fini sur $S$ ».} L'idée est qu'un groupe rigide est déterminé par sa torsion
première aux caractéristiques, qui est étale et ne bouge pas dans une famille.
En sont : les tores $\mathbf{G}_m$, les variétés abéliennes, les groupes
$\mathbf{Z}/\ell$ ; n'en sont pas $\mathbf{G}_a$ ni $\mathbf{Z}/p$ en
caractéristique $p$. La page affirme qu'une extension de groupes rigides est
rigide, qu'un quotient de rigide est rigide, et qu'une extension non triviale
d'une variété abélienne par $\mathbf{G}_a$ est rigide, du moins en
caractéristique $0$ --- ce qu'on vérifie : l'adhérence de la torsion est un
sous-groupe qui se surjecte sur la variété abélienne et ne contient pas
$\mathbf{G}_a$, sinon l'extension serait scindée. La page 50 propose de
renforcer la notion sur un corps (« très rigide ») en exigeant la propriété de
tous les sous-groupes fermés connexes, ce qui revient à une suite de
composition à quotients $\mathbf{Z}/\ell$, $\mathbf{G}_m$ et variétés
abéliennes.

\emph{Proposition (page 43).} Si $G$ est rigide et $n$ premier aux
caractéristiques résiduelles, tout ouvert $U$ de $G$ est le plus petit
sous-schéma fermé de $U$ qui contienne les $G[n^k] \cap U$ ; « démonstration
donnée plus haut », c'est-à-dire dans les pages manquantes.

\emph{Proposition (page 43).} Soit $G$ un schéma en groupes commutatif
localement de type fini sur $S$, et $n$ un entier inversible sur $S$. Alors la
multiplication par $n$ est non ramifiée, donc $G[n]$ est non ramifié sur $S$ ;
si $G$ est plat, $n_G$ est étale et $G[n]$ est étale sur $S$. On se ramène à
un corps algébriquement clos, où l'application tangente de $n_G$ est la
multiplication par $n$, inversible.\footnote{La page ajoute au-dessus de la
ligne que $n_G$ est « surjectif » (lecture incertaine). Ce n'est vrai que sur
la composante neutre : pour le groupe constant $\mathbf{Z}$, la multiplication
par $n$ n'est pas surjective.}

\subsection*{Le théorème de rigidité}

\emph{Lemme (page 44).} Soient $S$ localement noethérien connexe, $X$ non
ramifié sur $S$, $X_1$ et $X_2$ deux sous-schémas de $X$ finis et étales sur
$S$, et $s \in S$ tel que $X_{1,s} = X_{2,s}$. Alors $X_1 = X_2$. On se ramène
à $X_1 \subset X_2$ ; le lieu où ils coïncident est ouvert, et il est fermé
parce qu'un point du bord serait spécialisation d'un point où ils
coïncident.\footnote{Le cas local passe par le complété de $A$ et une
réduction au cas où $X$ est fini sur $S$ ; il n'est lu qu'en partie.}

\emph{Théorème de rigidité (page 45).} Soient $S$ localement noethérien
connexe, $G$ un schéma en groupes de type fini sur $S$, $G_1$ et $G_2$ deux
sous-schémas en groupes fermés rigides, et $s \in S$ tel que
$G_{1,s} = G_{2,s}$. Alors $G_1 = G_2$.

Les corollaires (pages 46 à 49) : deux homomorphismes d'un groupe rigide qui
coïncident en une fibre sont égaux (on applique le théorème au graphe) ; si
$S' \to S$ est fidèlement plat de type fini à fibres géométriquement connexes,
un sous-groupe fermé rigide de $G \times_S S'$ descend, de façon unique, en un
sous-groupe de $G$, puisque ses deux images inverses sur $S' \times_S S'$
coïncident sur la diagonale, et que toute composante connexe de
$S' \times_S S'$ la rencontre ; en particulier, pour une extension primaire
$k'/k$, un sous-groupe rigide de $G \otimes_k k'$ provient d'un unique
sous-groupe de $G$, et le foncteur $H \mapsto H \otimes_k k'$ est pleinement
fidèle sur les groupes commutatifs rigides de type fini.\footnote{Pour le
corollaire 3, la page « admet » qu'un sous-groupe rigide sur $k'$ est défini sur
une sous-algèbre de type fini et y reste rigide, et note en marge « précision à
prouver ». Une rédaction de la descente par limites inductives, page 47, est
encadrée et barrée.}

\emph{Remarque (pages 50 et 51).} L'hypothèse de rigidité n'est pas
superflue : sur un anneau de valuation discrète de caractéristique résiduelle
$p$, un schéma abélien $A$ peut avoir deux sous-groupes étales de degré $p$,
isomorphes à $(\mathbf{Z}/p)_S$, de même fibre spéciale et distincts. La page
construit l'exemple dans un produit $C' \times C$ de courbes elliptiques, avec
les sous-groupes engendrés par deux sections $(\xi, e)$ et $(\xi, \eta)$ ; la
construction elle-même ne se lit pas.

\subsection*{Extensions primaires (pages 52 à 57)}

Les pages suivantes traitent le cas où seule la composante neutre est rigide.
\emph{Proposition (page 52).} Si $k'/k$ est primaire, $G$ un groupe de type
fini sur $k$, et $H'$ un sous-groupe fermé de $G \otimes_k k'$ dont la
composante neutre est rigide, il existe une extension radicielle finie $k_1$ de
$k$ et un sous-groupe fermé $H_1$ de $G \otimes_k k_1$ tels que $H'$ et $H_1$
deviennent égaux sur $k' \otimes_k k_1$. La démonstration applique la rigidité
à la composante neutre, puis traite un groupe fini, pour lequel
$G \otimes_k k' \to G$ est bijectif. Les corollaires (le cas séparable, le
cas parfait, une version relative au-dessus d'un $S' \to S$ surjectif à fibres
irréductibles) ne sont lus qu'en partie.

\emph{Lemme (page 57).} Soient $u \colon G \to G_1$ un morphisme surjectif de
schémas en groupes de type fini et plats sur $S$ localement noethérien, donc
plat, et $N$ son noyau. Les sous-schémas fermés $H$ de $G$ stables par $N$
correspondent biunivoquement aux sous-schémas fermés $H_1$ de $G_1$, par
$H = u^{-1}(H_1)$ ; $H$ est plat, resp.\ un sous-groupe, si et seulement si
$H_1$ l'est. On le démontre par descente. La proposition de la page 55, qui
s'appuie sur ce lemme, décrit les familles plates de sous-groupes de $G$ dont
la composante neutre générique est rigide comme images inverses de familles de
sous-groupes radiciels d'un quotient fixe $G/H$ ; elle est lue par fragments.

\pagerange{59}{63}
\section*{Le schéma d'Albanese (pages 58 à 63)}

\subsection*{Le foncteur}

Soit $X$ fidèlement plat et quasi-projectif sur $S$. Pour un $S$-groupe $G$,
les données suivantes sont équivalentes (page 59) : un torseur $P$ sous $G$ et
un $S$-morphisme $X \to P$ ; une donnée de descente sur $G \times_S X$
relativement à $X \to S$, compatible à la structure de torseur ; un morphisme
$\varphi \colon X \times_S X \to G$ satisfaisant la condition de cocycle
\[
\varphi(x, y)\,\varphi(y, z) = \varphi(x, z) ,
\]
donc $\varphi(x, x) = e$, avec une condition d'effectivité automatiquement
satisfaite si $G$ est affine sur $S$ ou si $X \to S$ admet localement des
quasi-sections finies, par exemple s'il est projectif.

Pour $X$ projectif sur $S$ avec $\mathcal{O}_S \simeq f_*\mathcal{O}_X$
universellement, on note $T_X(P)$ le foncteur qui à un torseur $P$ sous un
schéma abélien associe $\mathrm{Hom}_S(X, P)$, et l'on cherche à le
représenter : un couple universel $(P, \psi)$. La page 60 l'esquisse par un
argument de sous-groupe minimal par lequel $\varphi$ se factorise. Pour
$P = A$ un schéma abélien, en faisant le changement de base $S' = X$ qui
munit $X$ d'une section,
\[
T_X(A) \simeq \mathrm{Hom}_{S'\text{-pointés}}(X_{S'}, A_{S'})
= \mathrm{Hom}_{S'\text{-pointés}}(X_{S'}, \mathrm{Pic}^0_{A^\vee_{S'}/S'})
\simeq \mathrm{Hom}_{S'\text{-gr}}(A^\vee_{S'}, \mathrm{Pic}_{X_{S'}/S'}) ,
\]
la troisième égalité passant par les correspondances divisorielles entre
$X_{S'}$ et $A^\vee_{S'}$, et l'on redescend sur $S$ « d'après la platitude et
la rigidité ». D'où, encadré\footnote{La page note $B = \hat{A} = \mathrm{Pic}^0_{A/S}$, de
sorte que $A = \mathrm{Pic}^0_{B/S}$ ; nous écrivons $A^\vee$. La fin de la
première ligne, prolongée dans la marge, est illisible.} :
\[
T_X(A) \simeq \mathrm{Hom}_{S\text{-gr}}(A^\vee, \mathrm{Pic}_{X/S}) .
\]

\subsection*{Albanese, dual de Picard}

Si $\mathrm{Pic}_{X/S}$ contient un plus grand sous-schéma abélien $B_{X/S}$,
les homomorphismes d'un schéma abélien dans $\mathrm{Pic}_{X/S}$ se factorisent
par lui, et la dualité des schémas abéliens donne
\[
\mathrm{Hom}_{S\text{-gr}}(A^\vee, \mathrm{Pic}_{X/S}) \simeq \mathrm{Hom}_{S\text{-gr}}(\mathrm{Alb}^0_{X/S}, A),
\qquad \mathrm{Alb}^0_{X/S} = B_{X/S}^\vee .
\]
Il y a donc un torseur canonique $\mathrm{Alb}^1_{X/S}$ sous
$\mathrm{Alb}^0_{X/S}$ et un morphisme $\psi_{X/S} \colon X \to \mathrm{Alb}^1_{X/S}$
tels que la donnée d'un torseur $P$ sous $A$ et d'un morphisme $X \to P$
équivaille à celle d'un homomorphisme $u \colon \mathrm{Alb}^0_{X/S} \to A$,
par $P = \mathrm{Alb}^1_{X/S} \times^{\mathrm{Alb}^0_{X/S}} (A, u)$. Ce sont
les schémas d'Albanese homogène et inhomogène de $X/S$ ; le triple
$(\mathrm{Alb}^0, \mathrm{Alb}^1, \psi)$ représente $T_X$ par construction, et
sa formation commute à tout changement de base.\footnote{Les pages écrivent
$\underline{A}^0_{X/S}$, $\underline{A}^1_{X/S}$ et
$\underline{B}_{X/S}$ ; le $\hat{B}$ du premier argument, page 61, est écrit
sur un $A$ biffé.}

Le N.B. de la page 62 propose de définir axiomatiquement le schéma d'Albanese
par la représentabilité de $T_X$ restreint aux schémas abéliens, stable par
changement de base, et juge plausible qu'il existe toujours sur un corps, même
quand $\mathrm{Pic}^0_{X/k}$ n'est pas propre, en prenant le plus grand
sous-schéma abélien de $\mathrm{Pic}_{X/k}$ --- dont il faudrait établir une
propriété universelle par changement de base.

\subsection*{Questions (page 63)}

Si $\mathrm{Pic}_{X/S}$ est lisse le long de la section unité, alors
$\mathrm{Pic}^0_{X/S}$ est ouvert dans $\mathrm{Pic}_{X/S}$, les fibres étant
géométriquement localement intègres (la marge renvoie à EGA~IV). S'il est de plus propre, c'est un
schéma abélien, et alors (a) les fibres de $\mathrm{Pic}_{X/S}$ sont lisses ;
(b) $\mathcal{O}_X$ est cohomologiquement plat sur $S$ en dimension $1$. La
réciproque est posée, avec trois points d'interrogation, et Grothendieck
conjecture plus généralement qu'un groupe sur une base réduite, à fibres
lisses, dont le faisceau $\omega$ le long de la section unité est localement
libre, est lisse le long de cette section. Il demande enfin si, pour $X/S$
propre et lisse en caractéristique $0$, la condition de platitude
cohomologique est automatique : il en a une démonstration algébrique quand $S$
est réduit, et ne sait rien quand $S$ est artinien.\footnote{Toute cette page
est lue mot à mot sous réserve, et la « condition » de la dernière question
n'est pas nommée. Si, comme le contexte l'indique, il s'agit de la platitude
cohomologique, la réponse est connue depuis : pour $X/S$ propre et lisse en caractéristique $0$, les
$R^q f_*\mathcal{O}_X$ sont localement libres et commutent au changement de
base (Deligne, 1968, par la dégénérescence de la suite spectrale de Hodge vers
de Rham).}

\pagerange{65}{65}
\section*{Le plan d'un traité (pages 64 et 65)}

Sous la chemise « Plan », une page intitulée « Picard » énumère, dans l'ordre
où elle les écrit et avec les numéros qu'on y lit en dernier :
\begin{itemize}
  \item 1. Théorèmes d'existence absolus ;
  \item 2. Théorèmes d'existence relatifs ;
  \item 3. Critères d'équivalence et théorèmes de finitude (« Néron--Igusa
    exclus ?? ») ;
  \item 4. $\mathrm{Pic}^\tau_{X/S}$, équivalence numérique, groupe de
    Néron--Severi ;
  \item 5. Critères de séparation, de propreté pour $\mathrm{Pic}_{X/S}$ ;
  \item (sans numéro) Critères de lissité pour $\mathrm{Pic}_{X/S}$ ;
  \item 8. Albanese ;
  \item (7) Dualité des schémas abéliens, notamment autodualité ;
  \item 9. Picard des variétés abéliennes (« Torelli ? ») ;
  \item 5. Variations sur le carré et le cube : correspondances divisorielles,
    multiplication par un entier dans $\mathrm{Pic}^\tau_{X/S}$ ;
  \item 10. Compléments au-dessus d'un corps : variétés non complètes,
    diviseurs sur les courbes, critères d'équivalence, cas normal, rappel des
    résultats spéciaux obtenus dans les pages précédentes ;
  \item 11. Phénomènes spéciaux en caractéristique $0$ ;
  \item 12. (?) Picard locaux et pro-groupes.
\end{itemize}
Les numéros sont remaniés à l'encre : des flèches en marge reportent les
critères de lissité entre 4 et 5, font des variations sur le cube le
chapitre 5, et échangent Albanese et la dualité.\footnote{Les numéros biffés
sont 6 sous le premier 5, 6 pour la lissité, 7 sous le 8 d'Albanese, 6 sous le
second 5. L'ordre définitif n'est pas sûr, et deux chapitres portent le
numéro 5. La transcription lit « progressions » pour le dernier mot du
chapitre 12 ; « pro-groupes » est notre conjecture, d'après les pro-groupes du
tapuscrit des pages 2 à 12.} Au bas de la page, un graphe des dépendances,
dont la colonne $2 \Rightarrow 3 \Rightarrow 4 \Rightarrow 5 \Rightarrow 7
\Rightarrow 8 \Rightarrow 9 \to 10$ est sûre, et dont les flèches obliques,
venues de 1 et de 4 vers le chapitre 6, le sont moins.

Le dossier fournit des matériaux pour presque tous ces chapitres : les
existences relatives (pages 125 à 139), la finitude (pages 67 à 71), les
critères d'équivalence (pages 96 à 100), $\mathrm{Pic}^\tau$ (pages 195 à 211),
la séparation (pages 115 à 123 et 312), Albanese (pages 59 à 63), les
correspondances et le cube (pages 14 à 31), les courbes (pages 141 à 154), les
Picard locaux (pages 2 à 12, 184 à 193, 226 à 271).

\pagerange{67}{71}
\section*{Le théorème de finitude (pages 66 à 71)}

\subsection*{La démonstration de Matsusaka (pages 67 et 68)}

Soient $X$ une surface projective lisse sur un corps algébriquement clos $k$,
$H$ un diviseur ample et $K$ le diviseur canonique. Riemann--Roch donne
$\chi(D) = \frac12 D\cdot(D - K) + \chi(\mathcal{O}_X)$, et la dualité
$\chi(D) \leq \ell(D) + \ell(K - D)$, où $\ell = \dim H^0$. Si
\[
\text{a)}\ (D - K)\cdot H > 0 \qquad \text{et} \qquad
\text{b)}\ \tfrac12 D\cdot(D - K) + \chi(\mathcal{O}_X) > 0 ,
\]
alors $\ell(K - D) = 0$, puisque $K - D$ est de degré négatif, et
$\ell(D) \geq \chi(D) > 0$.\footnote{La page écrit les deux inégalités larges.
Il faut a) stricte --- pour $D = K$, on a $(D - K)\cdot H = 0$ et
$\ell(K - D) = 1$ --- et b) stricte pour conclure $\ell(D) > 0$.} Ces
conditions ne dépendent que de la classe numérique de $D$ et sont satisfaites
par $D = nH$, $n$ grand. Si $\Delta$ est numériquement équivalent à $0$,
$D + \Delta$ les satisfait aussi, et est donc équivalent à un diviseur positif.

\emph{Proposition.} Soient $D$ satisfaisant a) et b) et $\mathfrak{M}$ la
famille des diviseurs positifs numériquement équivalents à $D$. Tout diviseur
numériquement équivalent à $0$ est linéairement équivalent à un $D' - D$ avec
$D' \in \mathfrak{M}$. Par Riemann--Roch, le polynôme de Hilbert des éléments
de $\mathfrak{M}$ est fixé ; par la finitude de Chow appliquée aux schémas de
Hilbert, $\mathfrak{M}$ est quasi-compacte. D'où, en chaîne :
$\mathrm{Pic}^{\equiv 0}_{X/k}$ est quasi-compact ;
$\mathrm{Pic}^{\equiv 0}_{X/k}/\mathrm{Pic}^0_{X/k}$ est fini ;
$\mathrm{Pic}^{\equiv 0}_{X/k} = \mathrm{Pic}^\tau_{X/k}$, et la torsion du
groupe de Néron--Severi est finie ; la forme d'intersection sur
$\mathrm{NS}/\text{torsion}$ est non dégénérée.\footnote{La page note
$\mathrm{Pic}^{\mathfrak{n}}$ ou $\mathrm{Pic}^n$ le sous-groupe
numériquement trivial ; le premier exposant de l'égalité
« $\mathrm{Pic}^{\tau} = \mathrm{Pic}^{\tau}$ » est peut-être un trait et non
un $\tau$.}

\subsection*{Igusa (page 69)}

Le résultat « plus profond » d'Igusa : soient $b_i$ les nombres de Betti
($\ell$-adiques) de $X$ et $E(X) = \sum (-1)^i b_i$, qui est aussi $c_2(X)$,
l'auto-intersection de la diagonale dans $X \times X$. On a $b_0 = b_4 = 1$ et
$b_1 = b_3 = 2p$, où $p = \dim \mathrm{Pic}_{X/k}$, d'où
$b_2 = E - 2 + 4p$ numériquement,\footnote{La page écrit $b_1 = b_3 = \dim \mathrm{Pic}_X$ et en tire, de façon
cohérente avec cette erreur, $b_2 = E - 2 + 2p$. Le premier nombre de Betti
est le double de la dimension de la variété de Picard (réduite).} et l'inégalité
\[
\rho \leq b_2, \qquad \rho = \operatorname{rg} \mathrm{NS}(X) .
\]

\subsection*{Un commentaire sur Kleiman (pages 70 et 71)}

La feuille, plus tardive, commente au crayon les travaux de Kleiman, avec des
renvois à une liste de critères numérotés (i) à (viii) que le dossier ne
contient pas. Kleiman démontre, pour $X$ projectif sur un corps, les critères
(viii) de $\tau$-équivalence à $0$ sans le « fourbis délicat » (i), (ii), puis
le cas propre par le lemme de Chow ; il montre directement que
$\mathrm{Pic}^{\mathrm{num}}_{X/S}$ est quasi-compact. Par le théorème de
l'indice sur une surface, il démontre le critère d'équivalence : si $X$ est
projectif, à composantes de dimension $\geq 3$, et $Y$ une section
hyperplane, l'image inverse de $\mathrm{Pic}^\tau_{Y/k}$ est
$\mathrm{Pic}^\tau_{X/k}$. Avec l'intégrité de Picard--Igusa, le théorème de
Néron--Severi et la résolution des surfaces d'Abhyankar, cela entraîne que les
groupes de Néron--Severi des fibres géométriques d'un $X/S$ propre de
présentation finie sur une base quasi-compacte sont « uniformément » de type
fini : torsion d'ordre borné, rang borné.

La page 71 se demande si ces résultats donnent le critère d'équivalence sous la
forme (ii) : on se ramène au cas où tout est représentable et plat, et le
morphisme induit sur les groupes de Néron--Severi modulo torsion est un
monomorphisme ; « il n'est pas du tout clair pourquoi il serait de type fini ».
Sa propre démonstration montre que
$\mathrm{Pic}_{X/S} \to \mathrm{Pic}_{Y/S}$ est génériquement quasi-affine sur
$S$, peut-être affine --- ce que la lettre à Murre de 1962, reproduite aux
pages 72 à 78, annonçait déjà.\footnote{Le sigle « NSL » de la page est lu ici « Néron--Severi modulo
torsion » ; c'est une interprétation. Le dernier mot de la page est lu
« topos », peut-être « type ».}

\pagerange{72}{78}
\section*{La lettre à Murre du 18 juillet 1962 (pages 72 à 78)}

La lettre existe en deux états dactylographiés : une copie au net (pages 72 à
74) et l'état antérieur, corrigé à la main et à la machine, et signé (pages 75
à 78). On lit le texte au net, et l'on dit en note où l'autre diffère, puisque
ces différences sont l'information.

\subsection*{Finitude}

Grothendieck a « substantially improved » les résultats de la dernière section
de son dernier exposé Bourbaki.\footnote{Vraisemblablement l'exposé n°~236,
« Les schémas de Picard : propriétés générales » (1962) ; le « first talk on
Picard schemes » est alors le n°~232, et le « first talk on descent » le
n°~190. L'identification est la nôtre.} Le résultat principal --- pour un
morphisme lisse et projectif à fibres géométriquement connexes, les morceaux
$\mathrm{Pic}^P_{X/S}$ indexés par un polynôme de Hilbert sont de type fini ---
a été étendu par Mumford au cas plat à fibres géométriquement intègres (au
moins normales). Avec ce résultat, « the proof of which is quite simple and
beautiful », il se débarrasse de la normalité et ne garde que les deux premiers
coefficients non triviaux du polynôme de Hilbert. Les deux clés :

(i) Si $X$, $Y$ sont propres sur $S$ noethérien et $f \colon X \to Y$ un
$S$-morphisme surjectif, les schémas de Picard existant, alors
$f^* \colon \mathrm{Pic}_{Y/S} \to \mathrm{Pic}_{X/S}$ est de type fini, et
même affine sur un corps : une partie de $\mathrm{Pic}_{Y/S}$ est
quasi-compacte si et seulement si son image l'est.

(ii) De même pour l'immersion d'un diviseur $X \to Y$, si $Y/S$ est projectif à
fibres dont toutes les composantes sont de dimension $\geq 3$ et $X$ le lieu
des zéros d'une section d'un faisceau inversible ample relativement à $S$.

Corollaire : pour $X/S$ propre et $n \neq 0$, la multiplication par $n$ dans le
schéma de Picard est de type fini.

\subsection*{Descente non plate}

La démonstration de (i) passe par la descente non plate --- « the fact that I
do not have any good effectivity criterion does not hamper » --- et une analyse
plus fine donne :

\emph{Théorème.} Soient $S$ intègre noethérien, $X$ et $X'$ propres sur $S$,
$f \colon X' \to X$ un $S$-morphisme surjectif, et
$f^* \colon \mathrm{Pic}_{X/S} \to \mathrm{Pic}_{X'/S}$. On suppose a) que le
problème d'existence A ci-dessous a toujours une solution pour $X/S$ (c'est le
cas si $X/S$ est projectif) ; b) que $f$ est, sur la fibre générique, un
morphisme de descente :
$\mathcal{O}_{X_\eta} \to f_*\mathcal{O}_{X'_\eta} \rightrightarrows h_*\mathcal{O}_{X''_\eta}$
exacte. Alors, après restriction à un ouvert non vide de $S$, $f^*$ est
représentable par des morphismes quasi-affines : dans la factorisation
$f^* = v u$ par le foncteur des données de descente, $u$ est affine, et le
monomorphisme $v$ est représentable par des sommes directes finies
d'immersions.\footnote{L'état antérieur (page 76) raye « and B » après
« problem A », que la copie au net (page 73) garde ; une note au crayon, en
face de « provided we replace $S$ … », demande « $S$ … needed ? ». La copie au
net écrit $f(\underline{O}_{X'_s})$ sans l'astérisque.}

\emph{Corollaires.} Sans b), en permettant de remplacer $X$ par des $X_1$ finis
sur $X$, la même conclusion vaut ; et, par le lemme de Chow et le théorème
d'existence principal, si pour tout $X'$ fini sur $X$ le problème A a une
solution génériquement sur $S$ (c'est le cas si $X$ est projectif), alors
au-dessus d'un ouvert non vide de $S$, $\mathrm{Pic}_{X/S}$ existe, est séparé,
et ses composantes connexes sont de type fini.

\emph{N.B.} Rien n'indique qu'on puisse remplacer « quasi-affine » par
« affine », sinon sur un corps, où un groupe algébrique quasi-affine est
affine. Un contre-exemple sur $k[t]$, avec $X$ et $X'$ projectifs et lisses et
$X' \to X$ birationnel, ou avec $X$, $X'$ projectifs normaux et $f$ fini, serait
intéressant ; dans le second cas, il fournirait un contre-exemple au problème
d'effectivité pour un morphisme fini posé dans le premier exposé sur la
descente.\footnote{L'état antérieur porte en marge, à la main, face au premier
contre-exemple : « impossible if $X$ smooth! ».}

\subsection*{Le problème A}

Étant donnés $X/S$ et un Module $F$ sur $X$, représenter le foncteur qui, à
$S' \to S$, associe un ensemble à un élément ou vide selon que
$F' = F \otimes_S S'$ sur $X' = X \times_S S'$ est plat sur $S'$ ou non. On
dit que le problème A a toujours une solution pour $X/S$ si, pour tout
cohérent $F'$ sur un $X'/S'$, ce foncteur est représentable par un $S'$-schéma
de type fini. L'étape principale de l'existence des schémas de Hilbert le
donne pour $X/S$ projectif, comme une somme disjointe de sous-schémas indexée
par les polynômes de Hilbert. Grothendieck s'attend à la représentabilité dès
que $X/S$ est propre, propose une réduction au cas projectif par le lemme de
Chow et une récurrence sur la dimension relative, et une approche par le cas
local complet et la descente fidèlement plate. Hartshorne, « a Harvard
student », a rencontré le problème, « but I doubt he will work seriously on
it » ; d'où l'appel cité dans le résumé.\footnote{C'est ce qu'on appelle
aujourd'hui la platification universelle. Pour $X/S$ propre de présentation
finie, Raynaud et Gruson (1971) ont montré que ce foncteur est représentable
par un monomorphisme de présentation finie --- et non nécessairement, comme
dans le cas projectif, par une somme de sous-schémas. Le « problème B », que le théorème de la copie au net
mentionne avec le problème A, n'est défini dans aucun des deux états.}

\pagerange{80}{93}
\section*{Fibrés à connexion et extension du dual d'une variété abélienne (pages 79 à 93)}

La chemise porte « Extension du dual d'une VA et faisceaux inversibles à
connexion ». Soit $f \colon X \to S$ plat de présentation finie, $R_{X/S}$ le
faisceau des fonctions méromorphes relatives et
$\mathrm{Div}_{X/S} = R^*_{X/S}/\mathcal{O}^*_X$ celui des diviseurs relatifs.

\subsection*{La construction par les diviseurs (pages 80 et 82)}

La dérivée logarithmique $\varphi \mapsto d\varphi/\varphi$ induit
\[
\mathrm{Div}_{X/S} \longrightarrow R_{X/S}(\Omega^1_{X/S})/\Omega^1_{X/S} ,
\]
injective si $S$ est de caractéristique nulle et $f$ normal. Le produit fibré
$E_{X/S}$ des couples $(D, \omega)$, $D$ un diviseur et $\omega$ une forme
méromorphe telle que $d\log D \equiv \omega$ modulo les formes régulières, est
une extension
\[
0 \to \Omega^1_{X/S} \to E_{X/S} \to \mathrm{Div}_{X/S} \to 0 .
\]
En prenant $f_*$, avec $\omega_{X/S} = f_*\Omega^1_{X/S}$, on obtient une suite
exacte
$0 \to \omega_{X/S} \to f_*E_{X/S} \to f_*\mathrm{Div}_{X/S} \xrightarrow{\partial} R^1 f_*\Omega^1_{X/S}$,
où $\partial$ est l'homomorphisme bien connu déduit de
$\mathrm{Pic}_{X/S} \to R^1 f_*\Omega^1_{X/S}$ --- la classe d'un faisceau
inversible en cohomologie de Hodge, ce qu'on appelle aujourd'hui sa classe
d'Atiyah. On pose
\[
\mathrm{Pic}^w_{X/S} = \operatorname{Ker}\bigl(\mathrm{Pic}_{X/S} \to R^1 f_*\Omega^1_{X/S}\bigr) ,
\]
le sous-foncteur des classes qui admettent, localement sur $S$, une connexion
relative. Si $X \to S$ est quasi-projectif, $\mathrm{Pic}_{X/S}$ est le
quotient de $f_*\mathrm{Div}_{X/S}$ par les diviseurs principaux
$P_{X/S} = f_*R^*_{X/S}/f_*\mathcal{O}^*_X$, et $\mathrm{Pic}^w_{X/S}$ celui
du noyau de $\partial$. Le morphisme $\varphi \mapsto (\operatorname{div}\varphi, d\varphi/\varphi)$
relève $P_{X/S}$ dans $f_*E_{X/S}$ ; il est nul sur $f_*\mathcal{O}^*_X$ dès que
$f_*\mathcal{O}_X \simeq \mathcal{O}_S$ universellement, ce qu'on suppose,\footnote{Sans cette hypothèse, note la marge, on n'obtient
qu'une extension par $\omega_{X/S}/\operatorname{Im} f_*\mathcal{O}^*_X$.} et
l'on obtient l'extension
\[
0 \to \omega_{X/S} \to \mathcal{M}_{X/S} \to \mathrm{Pic}^w_{X/S} \to 0, \qquad
\mathcal{M}_{X/S} = f_*E_{X/S}/P_{X/S} .
\]
Dans le « cas
favorable » où $\Omega^1_{X/S}$ est cohomologiquement plat en dimension $0$ ---
$f$ lisse et projectif en caractéristique nulle (?), ou $X$ un schéma abélien
--- $\omega_{X/S}$ et $R^1 f_*\Omega^1_{X/S}$ sont des fibrés vectoriels, et si
$\mathrm{Pic}_{X/S}$ est représentable, $\mathrm{Pic}^w_{X/S}$ et
$\mathcal{M}_{X/S}$ le sont aussi.

\subsection*{La définition directe (pages 84 et 86)}

Soit $\mathcal{M}^\#_{X/S}$ le faisceau associé au préfaisceau qui, à $T$,
associe les classes de faisceaux inversibles sur $X \times_S T$ munis d'une
connexion relative à $T$. La suite exacte de préfaisceaux
$H^0(\mathcal{O}^*_X) \xrightarrow{d\log} H^0(\Omega^1) \to \mathcal{Q}(T) \to \mathrm{Pic}(X_T) \to H^1(\Omega^1_{X/S})$
donne, par faisceautisation,
\[
0 \to \omega_{X/S}/\operatorname{Im} f_*\mathcal{O}^*_X \to \mathcal{M}^\#_{X/S} \to \mathrm{Pic}^w_{X/S} \to 0 ,
\]
et un morphisme d'extensions $\alpha \colon f_*E_{X/S} \to \mathcal{M}^\#_{X/S}$.
Pour le définir, on interprète $E_{X/S}$ : sur un ouvert, la donnée d'une
connexion $d_L$ sur $L = \mathcal{O}_X(D)$, c'est-à-dire d'un scindage de la
suite des parties principales
$0 \to L \otimes \Omega^1_{X/S} \to P^1(L) \to L \to 0$, équivaut à celle de
la forme méromorphe $\bar\omega = d_L\varphi_D/\varphi_D$, $\varphi_D$ la
section canonique, et une forme méromorphe provient d'une connexion si et
seulement si $\bar\omega - df/f$ est régulière pour toute équation locale $f$
de $D$. Ainsi $E_{X/S}$ est le faisceau des couples (diviseur, connexion sur le
faisceau associé), et $\alpha$ est évident ; il est un isomorphisme quand
$d\log$ est nul sur $f_*\mathcal{O}^*_X$.

La page 88 calcule la classe d'extension : une section $\xi$ de
$\mathrm{Pic}^w_{X/S}$ provenant d'un faisceau inversible $L$ donne, par image
inverse, l'élément de
$H^1(S, \omega_{X/S}) \simeq \operatorname{Ker}\bigl(H^1(X, \Omega^1_{X/S}) \to H^0(S, R^1 f_*\Omega^1_{X/S})\bigr)$
défini par $L$.\footnote{La page 85 est un brouillon barré, sans rapport
direct, sur les notations de descente et les torseurs $Z \times^G E$ ; la
page 87 porte une liste de choses à faire (« Demander Raynaud détails travaux
SGA 3 ; Prendre SGA 2 … ») et des notes au crayon, barrées, sur $d\log$ et les
résidus, où apparaît déjà l'extension $V(R^1\mathcal{O}_{A^\vee}) \to E(A^\vee) \to A^\vee$.}

\subsection*{L'extension vectorielle universelle (pages 90 à 93)}

Si $u \colon B \to \mathrm{Pic}_{X/S}$ est un homomorphisme d'un schéma abélien
et $\Omega^1_{X/S}$ est cohomologiquement plat en dimension $0$, $u$ se
factorise par $\mathrm{Pic}^w_{X/S}$, puisque le composé vers le fibré
vectoriel $R^1 f_*\Omega^1_{X/S}$ part d'un schéma propre à fibres connexes ;
en particulier $\mathrm{Pic}^{00}_{X/S} \subset \mathrm{Pic}^w_{X/S}$ quand il
existe. La donnée de $u$ équivaut à celle d'un torseur $A$ sous le dual
$B^\vee$ et d'un morphisme $X \to A$ (sous réserve d'existence de $B^\vee$ et
de possibilité de descente, par exemple $B$ projectif) --- c'est l'Albanese des
pages 59 à 63. Par fonctorialité, on obtient un diagramme d'extensions

\begin{tikzcd}[column sep=small, row sep=small]
0 \arrow[r] & \omega_{X/S} \arrow[r] & \mathcal{M}_{X/S} \arrow[r] & \mathrm{Pic}^w_{X/S} \arrow[r] & 0 \\
0 \arrow[r] & \omega_{A/S} \arrow[r] \arrow[u] & \mathcal{M}_{A/S} \arrow[r] \arrow[u] & \mathrm{Pic}^w_{A/S} \arrow[r] \arrow[u] & 0 \\
0 \arrow[r] & \omega_{A/S} \arrow[r] \arrow[u, no head, "\Vert" description] & \mathcal{M}^0_{A/S} \arrow[r] \arrow[u] & B \arrow[r] \arrow[u, hook] & 0
\end{tikzcd}

où $\mathcal{M}^0_{A/S}$ est l'image inverse de $\mathcal{M}_{A/S}$ par
$B = \mathrm{Pic}^0_{A/S} \hookrightarrow \mathrm{Pic}^w_{A/S}$ : l'extension de
$B$ par le fibré vectoriel $\omega_{A/S}$ des faisceaux inversibles
algébriquement équivalents à $0$ munis d'une connexion.\footnote{La page
surmonte l'$\omega$ de la première ligne d'un astérisque, peut-être une
rature ; elle note
$\omega_{A/S} \simeq V(R^1 f_{B^\vee *}\mathcal{O}_B)$ sous la dernière ligne.
« Fonctorialité doit être explicitée avec un peu de soin ! », dit la marge.}
Les translations de $A$ par $B^\vee$ opèrent trivialement sur
$\mathrm{Pic}^0_{A/S}$ et sur $\omega_{A/S}$, donc sur l'extension, qui descend :
$\mathcal{M}^0_{A/S}$ ne dépend que de $B$. Grothendieck veut y voir
l'\emph{extension universelle} de $B$ par un fibré vectoriel. Pour en calculer
la classe, il propose de prendre le faisceau de Poincaré sur $A \times_S B$ et
l'obstruction à le munir d'une connexion relative à $B$, classe dans
$H^1(A \times_S B, \Omega^1_{A \times B/B})$ ; la décomposition de Künneth
\[
H^{1,1}(A \times_S B) = H^{1,1}(A) \oplus H^{1,1}(B) \oplus H^{1,0}(A) \otimes H^{0,1}(B) \oplus H^{0,1}(A) \otimes H^{1,0}(B)
\]
et la fonctorialité par les deux inclusions montrent que la classe n'a de
composante que dans les termes mixtes, « ce qui peut se faire … ». Une question
barrée sur une connexion absolue canonique sur $\mathcal{M}_{X/S}$ termine la
page 93.

C'est ce qu'on appelle aujourd'hui l'extension vectorielle universelle d'un
schéma abélien, réalisée comme espace de modules des fibrés en droites à
connexion (Messing, Mazur et Messing, 1972--1974).

\pagerange{96}{100}
\section*{Critères d'équivalence par section hyperplane (pages 95 à 100)}

\subsection*{Le tableau (page 96)}

Soient $X$ une variété projective lisse de dimension $n$, $Y$ une section
hyperplane lisse et $X_m$ le $m$-ième voisinage infinitésimal de $Y$. La page
dresse le tableau suivant\footnote{La page note $TN$, $LN$, $N$ ce qu'on rend par torsion de
Néron--Severi, Néron--Severi modulo torsion, Néron--Severi ; « à $p$ près »
traduit son « exc.\ $p$ » (excepté la $p$-composante). Le tableau est
transcrit tel que la page le dresse ; nous n'en avons pas revérifié chaque
case.}, les conditions de cohomologie portant sur les
faisceaux $\mathcal{O}_Y(-m)$, $m > 0$, conormaux successifs de $Y$ dans
$X$.\footnote{La page écrit $\mathcal{O}_Y(m)$. Le noyau de
$\mathrm{Pic}(X_{m}) \to \mathrm{Pic}(X_{m-1})$ et l'obstruction sont dans
$H^1$ et $H^2$ de $\mathcal{J}^m/\mathcal{J}^{m+1} = \mathcal{O}_Y(-m)$ : ce
sont ces groupes qu'il faut annuler, comme dans le théorème de Lefschetz de
SGA~2, XII.}
\begin{itemize}
  \item $\mathrm{Pic}(X) \to \mathrm{Pic}(Y)$ : injectif si $n \geq 3$ et
    $H^1(Y, \mathcal{O}_Y(-m)) = 0$ ; surjectif si $n \geq 4$ et
    $H^2(Y, \mathcal{O}_Y(-m)) = 0$ (Lefschetz--Grauert) ; sans condition, de
    noyau un $p$-groupe fini si $n \geq 3$, nul si $Y$ est générale, nul si
    $n \geq 4$ ;
  \item $\mathrm{Pic}^\tau(X) \to \mathrm{Pic}^\tau(Y)$ : injectif à $p$ près
    si $n \geq 2$, bijectif à $p$ près si $n \geq 3$, injectif si $n \geq 2$
    et $Y$ générale (via $\pi_1$) ;
  \item $\mathrm{Pic}^0(X) \to \mathrm{Pic}^0(Y)$ : injectif à $p$ près si
    $n \geq 2$, surjectif si $n \geq 3$, bijectif si $n \geq 2$ et $Y$
    générale ;
  \item sur la torsion de Néron--Severi : bijectif à $p$ près si $n \geq 3$ ;
    sur Néron--Severi modulo torsion : injectif si $n \geq 3$ ; sur
    Néron--Severi : de noyau un $p$-groupe fini si $n \geq 3$, nul si
    $n \geq 4$ ou $Y$ générale ;
  \item $\mathrm{Pic}(X) \to \mathrm{Pic}(X_m)$, $m$ grand : injectif si
    $n \geq 3$, bijectif si $n \geq 4$.
\end{itemize}

La question posée : la surjectivité de $\mathrm{Pic}(X) \to \mathrm{Pic}(Y)$
pour $n \geq 4$ sans condition cohomologique --- équivalente à celle de
$\mathrm{Pic}(X_m) \to \mathrm{Pic}(Y)$ pour $m$ grand. Elle est vraie à $p$
près, les obstructions successives étant tuées par $p$ ; en caractéristique
$0$ la théorie kählérienne la donne. « Démonstration algébrique ??? » En marge,
le corollaire : $\rho(X) = \rho(Y)$ pour $n \geq 4$.\footnote{La note marginale
est en partie biffée et encadrée ; la lecture des conditions sur $n$ y est
douteuse.}

\subsection*{La section hyperplane générique (pages 98 à 100)}

\emph{Théorème.} Soient $k$ algébriquement clos, $X \subset \mathbf{P}^r_k$
normal intègre de dimension $\geq 2$ (donc de profondeur $\geq 2$), $G$ un
groupe algébrique commutatif affine, fini ou vectoriel, $K = k(t_1, \ldots, t_r)$,
$\overline{K}$ sa clôture algébrique, et $Y_{\overline{K}}$ la section
hyperplane générique. Alors $H^1(X, G) \to H^1(Y_{\overline{K}}, G)$ est
injectif.\footnote{La description de $G$ est raturée : « affine sans
composantes connexes (groupe … ou groupe algébrique fini) » ; la réduction qui
suit porte sur $\mathbf{G}_a$ et les groupes finis.}

Par la structure des groupes commutatifs finis, on se ramène à
$\mathbf{G}_a$, aux groupes finis étales, à $\mu_p$ et à $\alpha_p$, c'est-à-dire
aux énoncés : $H^1(X, \mathcal{O}_X) \to H^1(Y_{\overline{K}}, \mathcal{O})$
injectif (corollaire 1) ; $\pi_1(Y_{\overline{K}}) \to \pi_1(X)$ surjectif
(corollaire 2, Bertini et Lefschetz--Grauert, qui utilisent la profondeur) ;
injectivité sur $H^1(-, \alpha_p)$ et $H^1(-, \mu_p)$ (corollaires 3 et
4).\footnote{Le corollaire 3 est écrit pour
$H^1(X, \alpha_p) \to H^1(X_{\overline{K}}, \alpha_p)$, par lapsus pour
$Y_{\overline{K}}$.}

En caractéristique $p$, Frobenius opère sur $H^1(X, \mathcal{O}_X)$, qui se
décompose en une partie où il est bijectif, égale à
$H^1(X, \mathbf{Z}/p) \otimes k$ par la suite d'Artin--Schreier, et une partie
où il est nilpotent. Sur la première, l'injectivité vient du corollaire 2 ; sur
la seconde, il suffit de l'injectivité sur le noyau de $F$, qui est
$H^1(X, \alpha_p)$. Suivant Cartier, pour $X$ intègre et propre,
\[
H^1(X, \alpha_p) \simeq \operatorname{Ker}\bigl(F \mid H^1(X, \mathcal{O}_X)\bigr) \simeq H^0(X, \mathcal{O}_X/\mathcal{O}_X^p) \xrightarrow{\ d\ } H^0(X, \tilde{\Omega}^1_X),
\]
\[
H^1(X, \mu_p) \simeq \mathrm{Pic}(X)[p] \simeq H^0(X, \mathcal{O}^*_X/\mathcal{O}^{*p}_X) \xrightarrow{\ d\log\ } H^0(X, \tilde{\Omega}^1_X),
\]
les deux flèches étant injectives si $X$ est normal ; tout se ramène au lemme
« presque trivial » : $H^0(X, \tilde{\Omega}^1_{X/k}) \to H^0(Y_{\overline{K}}, \tilde{\Omega}^1)$
est injectif --- une forme qui s'annule sur toutes les sections hyperplanes
passant par un point lisse s'y annule, la dimension étant
$\geq 2$.\footnote{La page note $H^1(X, \mathcal{O}_X)^F$ et
$H^1(X, \mathcal{O}_X)^{1-F}$ les deux parties, et écrit la seconde ligne avec
$H^1(X, \mathcal{O}_X)^{1-F}$ et $\mathcal{O}_X/\mathcal{O}^*_X$ : ce n'est
qu'une recopie de la première ; le parallélisme demande
$\mathcal{O}^*_X/\mathcal{O}^{*p}_X$, et le premier membre est la $p$-torsion
de $\mathrm{Pic}$. $\tilde{\Omega}^1$ désigne les formes différentielles du
lieu lisse, prolongées.}

\emph{Corollaire.} Le noyau de
$\mathrm{Pic}^\tau_{X_{\overline{K}}} \to \mathrm{Pic}^\tau_{Y_{\overline{K}}}$
est un $p$-groupe fini unipotent. Si $X$ est normal, $\mathrm{Pic}^\tau_{X/k}$
est propre : par des sections hyperplanes générales successives, on se ramène
à une courbe normale, où c'est connu.\footnote{La page écrit
« $\mathrm{Pic}_{X/k}$ est propre » dans une phrase consacrée à
$\mathrm{Pic}^\tau$ ; $\mathrm{Pic}$ lui-même n'est pas de type fini. La phrase
qui suit sur la non-normalité n'est lue qu'en partie.}

La page 100 donne enfin une démonstration directe de la propreté de
$\mathrm{Pic}^0_{X/k}$ pour $X$ normal et projectif : par la structure des
groupes commutatifs, il suffit que tout morphisme
$\operatorname{Spec} k[t, t^{-1}] \to \mathrm{Pic}^0_{X/k}$ soit constant, ce
qui résulte du lemme : si $X$ est localement noethérien et normal,
$\mathrm{Pic}(X) \to \mathrm{Pic}(X[t, t^{-1}])$ est surjectif. Un faisceau
inversible sur $X[t, t^{-1}]$ est défini par un diviseur, dont l'adhérence est
localement principale « par raison de stabilité ».\footnote{La page barre
« bijectif » pour écrire « surjectif » ; l'injectivité vient d'une section.
La marge juge le lemme « certainement inutile, en procédant de façon
convenable ».}

\pagerange{101}{102}
\section*{Picard et profondeur (pages 101 et 102)}

\emph{Lemme.} Soient $X$ un schéma, $U$ un ouvert rétrocompact, $i \colon U \to X$,
$f \colon X' \to X$ fidèlement plat et quasi-compact, $U' = f^{-1}(U)$, et
supposons $\mathcal{O}_X \xrightarrow{\sim} i_*\mathcal{O}_U$ (ce qui équivaut à
la même condition sur $X'$).\footnote{La page écrit $U' = f^{-1}(X')$.} Notons
$P(-)$ l'ensemble des classes de faisceaux localement libres de type fini.
Alors le carré
\[
\begin{array}{ccc}
P(X) & \longrightarrow & P(X') \\
\downarrow & & \downarrow \\
P(U) & \longrightarrow & P(U')
\end{array}
\]
est cartésien, et de même pour $\mathrm{Pic}$. L'injectivité vient de ce que
la restriction des faisceaux localement libres de $X$ à $U$ est pleinement
fidèle. Pour la surjectivité, si $F$ sur $U$ et $E'$ sur $X'$ ont même image
sur $U'$, alors $E = i_*F$ convient : $f^*i_*F \simeq i'_*g^*F \simeq i'_*i'^*E' \simeq E'$
par changement de base plat, et l'on descend la locale liberté. D'où le
\emph{corollaire} : si $\mathrm{Pic}(U) \to \mathrm{Pic}(U')$ est injectif
(resp.\ bijectif), il en est de même de $\mathrm{Pic}(X) \to \mathrm{Pic}(X')$.
La marge ajoute qu'on peut prendre, plus généralement, les faisceaux de
profondeur $\geq 2$ hors de $U$.

\pagerange{104}{113}
\section*{Le théorème de Seshadri (pages 103 à 113)}

\emph{Théorème (Seshadri).} Soit $f \colon X \to S$ propre et plat, à fibres
géométriquement intègres et $(S_2)$, avec $S$ localement noethérien, normal et
irréductible de point générique $\eta$. On suppose que $\mathrm{Pic}_{X/S}$
existe et est essentiellement propre --- ce qui a lieu si $f$ est lisse, ou si
$X = S \times_k Y$ avec $Y$ normal sur un corps $k$. Soit $D$ un cycle de
codimension $1$ sur $X$ tel que a) le support $Z$ de $D$ ne contienne aucune
composante d'une fibre, et b) $D_\eta$ soit un diviseur de Cartier. Alors $D$
est un diviseur de Cartier.\footnote{La page corrige « diviseur » en
« 1-cycle » : il s'agit d'un cycle de codimension $1$, a priori de Weil. La
démonstration qui suit traite le cas $D \geq 0$, comme le note la marge
(« seulement si $D \geq 0$ ! ») ; la page 111 ramène le cas général au cas
positif quand $f$ est projectif.}

La démonstration (pages 104 à 109, $D \geq 0$) : $D_\eta$ définit une section
rationnelle de $\mathrm{Pic}_{X/S}$ ; si $T$ est l'adhérence de son domaine
de définition, $T \to S$ est birationnel et surjectif par la propreté
essentielle, et il suffit, par le théorème principal de Zariski, $S$ étant
normal, de voir que ses fibres sont finies. Le lemme : pour $t \in T$ au-dessus
de $s$, le faisceau $L_t$ que $t$ définit sur une clôture algébrique est
défini par un diviseur de Cartier positif contenu dans $Z$. On le voit sur un
anneau de valuation discrète $V$ dominant $\mathcal{O}_{T,t}$ : $Z$ est
universellement ouvert sur $S$ par le théorème de Chevalley, puisqu'aucune de
ses composantes ne contient de composante de fibre ; la section rationnelle
canonique, multipliée par une puissance convenable de l'uniformisante, devient
une pseudo-section régulière de diviseur $D'$ transversal aux fibres, et
$\operatorname{supp} D'_{s'} \subset (\overline{Z'_{\eta'}})_{s'} = Z'_{s'}$.
On construit ensuite un faisceau inversible $L$ sur $X$ prolongeant
$L_\eta$, une section de diviseur $D'$ transversal aux fibres avec
$D'_\eta = D_\eta$ ; la différence $D - D'$ est nulle sur $X_\eta$ et ne
contient aucune composante de fibre, donc est nulle, puisque ses composantes
seraient, au-dessus des points de codimension $1$ de $S$, des fibres
intègres.\footnote{Plusieurs passages sont barrés ou illisibles, en
particulier le bas de la page 108 ; les étapes sont celles que les énoncés
imposent, la marge demande à deux reprises « Précisez » et « explicitez ».}

Deux questions suivent : les conclusions restent-elles vraies si l'on suppose
seulement que le \emph{foncteur} de Picard est essentiellement propre, par
exemple pour $f$ lisse ? Et, pour $f$ une intersection complète relative de
dimension relative $\geq 3$ sur une base normale irréductible, sans aucune
hypothèse sur Picard ? Un corollaire (page 110) applique le théorème à une
relation d'équivalence propre et plate $R \rightrightarrows X$, à fibres
intègres et $(S_2)$, sur $X$ normal intègre : un cycle de codimension $1$ sur
$X$ dont l'image inverse vérifie a) et b) relativement à $R \to X$ est un
diviseur de Cartier. La démonstration (pages 111 à 113) passe par l'ouvert où
les deux projections sont lisses ; la page 111 est barrée de trois traits.

\pagerange{115}{123}
\section*{Séparation : unités et diviseurs à support dans les fibres (pages 114 à 123)}

La chemise porte « Séparation de Picard ». Le problème est celui de la lettre à
Murre : un faisceau trivial sur la fibre générique est défini par un diviseur
à support dans les fibres, et l'on veut que ce diviseur soit un multiple de la
fibre.

\emph{Lemme 1 (page 115).} Soient $A$ un anneau noethérien réduit, $f \in A$ non
diviseur de zéro, $\mathfrak{p}_i$ les idéaux premiers associés à $fA$, et
supposons $A/fA$ réduit, c'est-à-dire $fA = \bigcap \mathfrak{p}_i$, et la
normalisée $\tilde{A}$ égale à l'intersection des $A_{\mathfrak{p}}$, $\mathfrak{p}$
de hauteur $1$. Alors
\[
A^* = A_f^* \cap \bigcap_i A^*_{\mathfrak{p}_i} .
\]
Démonstration : $B = A_f \cap \bigcap A_{\mathfrak{p}_i}$ est contenu dans
chaque $A_{\mathfrak{p}}$ de hauteur $1$ --- un tel $\mathfrak{p}$ contenant $f$
est un $\mathfrak{p}_i$ --- donc dans $\tilde{A}$. Si $x = a/f^k \in B$ avec
$k \geq 1$ minimal, une équation de dépendance intégrale donne
$a^n \in f^kA \subset fA = \bigcap \mathfrak{p}_i$, d'où $a \in fA$, contre la
minimalité : $B = A$, et les unités suivent.\footnote{La fin de la démonstration
se poursuit hors des pages conservées ; la conclusion est celle que la page
encadre.}

Les corollaires (pages 117 et 120) : si $Z$ est affine sur $X$ réduit,
localement noethérien, à anneaux locaux « bons » (complétés satisfaisant une
condition de finitude et la condition des chaînes), $f$ une fonction non
diviseur de zéro avec $\mathcal{O}_X/f$ réduit, $Y = V(f)$ de composantes
$Y_i$ de points génériques $y_i$, alors une section rationnelle de $Z$ définie
sur $X_f$ et aux $y_i$ est définie partout ; une fonction rationnelle est
inversible si elle l'est sur $X_f$ et aux $y_i$ ; et, si une section $g$ de $\mathcal{O}_X$ sur $X_f$ est entière sur
$\mathcal{O}_X$, elle est définie partout (lemme 2) ; enfin l'on a une suite
exacte $0 \to A^* \to A_f^* \to \mathbf{Z}^I$, avec
$A_f^* \simeq A^* \times \mathbf{Z}$, $t = \varepsilon f^n$, quand $Y$ est
irréductible.\footnote{Le diagramme de la page 117 est très chargé et retracé
en partie ; on y voit, dans une boucle, une tête coiffée d'un chapeau à large
bord.}

\emph{Corollaires 3 et 4 (pages 121 à 123).} Sous ces conditions, un diviseur
$D$ sur $X$ nul sur $X_f$ et dont le coefficient sur chaque $Y_i$ est le même
entier $n$ est $n\,\operatorname{div}(f)$. Si $Y$ est « connexe par incidence »
--- deux composantes se joignent par une chaîne de composantes dont les
intersections deux à deux contiennent des composantes irréductibles communes
---, tout diviseur nul hors de $Y$ est un multiple de $\operatorname{div}(f)$.
La démonstration écrit localement $D = \frac1n \operatorname{div} t$ et, sur
chaque composante $X_i$, $t_i = \varepsilon_i f_i^{n_i}$ ; sur une composante
$Z_{ij}$ de $X_i \cap X_j$, si $n_i > n_j$, on aurait
$f_{ij}^{n_i - n_j} = \varepsilon_{ji}/\varepsilon_{ij}$, donc $f$ inversible
sur $Z_{ij}$, où il s'annule : $n_i = n_j$.\footnote{Les pages 121 et 122 sont
très raturées et la note marginale de la page 121 n'est lue que par
fragments.}

\emph{Théorème (page 123).} Soient $f \colon X \to Y$ propre et séparable,
$Y$ localement noethérien, et $L$ un faisceau inversible sur $X$. Il existe un
plus grand sous-schéma fermé $Y'$ de $Y$ au-dessus duquel $L$ provient d'un
faisceau sur la base, et sa formation commute au changement de base. La page
s'arrête sur cette formule, avant les hypothèses précises sur les fibres ;
l'énoncé, ainsi tronqué, est celui qui fait de la séparation de
$\mathrm{Pic}_{X/S}$ la fermeture du lieu où un faisceau devient
trivial.\footnote{Le sigle $\mathcal{S}(L/X/Y)$ de ce sous-schéma est lu sur
une lettre bouclée, de lecture incertaine. La condition sur le nombre de
composantes géométriques des fibres est illisible.}

\pagerange{125}{139}
\section*{Théorèmes d'existence relatifs (pages 124 à 139)}

\subsection*{Le principe : recoller par descente (pages 125 à 128)}

\emph{Proposition.} Soient $f \colon X \to S$ et $f' \colon X' \to S'$ se
factorisant en $X \xrightarrow{f_1} \Sigma \xrightarrow{f_2} S$ et
$X' \xrightarrow{f'_1} \Sigma' \to S$, munis de sections $g$, $g'$ sur $\Sigma$,
$\Sigma'$, et $p \colon X' \to X$ ; soient $Z = g(\Sigma)$ et $Z'$, $Z''$ ses
images inverses dans $X'$ et $X'' = X' \times_X X'$. On suppose : (i) $f_1$,
$f'_1$ et $Z'' \to S$ plats, $f_1$ et $p$ propres ; (ii)
$f_{1*}\mathcal{O}_X \simeq \mathcal{O}_\Sigma$ et
$f'_{1*}\mathcal{O}_{X'} \simeq \mathcal{O}_{\Sigma'}$ ; (iii) $Z' \to Z$ est un
morphisme de descente universel pour les Modules inversibles ; (iv)
$p \colon X' \to X$ est un morphisme de descente effective universel pour les
Modules inversibles ; (v) $g'$ est une section sur $Z'$. Alors
$\mathrm{Pic}_{X/S} \to \prod_{\Sigma'/S}\mathrm{Pic}_{X'/\Sigma'}$ est
représentable par des morphismes affines ; si $\mathrm{Pic}_{X'/\Sigma'}$
existe, $\mathrm{Pic}_{X/S}$ aussi, et le morphisme est affine.

Démonstration : grâce aux sections, $\mathrm{Pic}_{X/S}(T)$ est l'ensemble des
faisceaux inversibles sur $X_T$ $g$-rigidifiés. Étant donné $L'$ sur $X'$
rigidifié, le foncteur des $L$ rigidifiés sur $X_T$ tels que $p^*L \simeq L'$
est représentable par un schéma affine sur $S$. En effet $L$ équivaut à une
donnée de descente : un isomorphisme $\varphi \colon q_1^*L' \to q_2^*L'$ sur
$X''_T$ et une rigidification $\psi \colon \mathcal{O}_{Z'} \to i'^*L'$,
soumis à la transitivité $r_2^*\varphi \circ r_1^*\varphi = r_3^*\varphi$ sur
$X'''$ et à la compatibilité
$\bar q_2^*\psi = i''^*\varphi \circ \bar q_1^*\psi$ sur $Z''$, avec une
troisième condition sur $\Sigma'$ que le bord de la feuille coupe. Les
isomorphismes entre faisceaux inversibles sur des schémas propres et plats
sont paramétrés par des schémas affines $K$, $L$, $M$, $N$ sur $S$, et les
conditions disent que des couples de morphismes $K \rightrightarrows M$,
$L \rightrightarrows N$ s'égalisent : le foncteur est représenté par un
sous-schéma fermé d'un schéma affine. Les notations sont celles du diagramme
de la page 126 :\footnote{Diagramme recomposé comme dans la transcription, chaque faisceau de
flèches parallèles réduit à une flèche étiquetée ; trois arcs $s_1, s_2, s_3$
de $X'''$ vers $X'$ sont omis. La page 128, qui reprend la démonstration, est
barrée de trois traits.}

\begin{tikzcd}[column sep=small, row sep=small, nodes={font=\scriptsize}]
X & X' \arrow[l, "p"'] & X'' \arrow[l, "{q_1, q_2}"'] & X''' \arrow[l, "{r_1, r_2, r_3}"'] \\
Z \arrow[u, "i"] & Z' \arrow[u, "{i'}"] \arrow[l, "\bar{p}"] & Z'' \arrow[u, "{i''}"] \arrow[l, "{\bar{q}_1, \bar{q}_2}"'] & \\
\Sigma \arrow[u, "g"] & \Sigma' \arrow[u, "{\bar{g}'}"'] & &
\end{tikzcd}

\emph{Corollaire (page 128).} Si $X$ est séparable et complet sur un corps $k$,
de composantes irréductibles réduites $X_i$, alors
$\mathrm{Pic}_{X/k} \to \prod_i \mathrm{Pic}_{X_i/k}$ est représentable par
des morphismes affines ; si les $\mathrm{Pic}_{X_i/k}$ existent,
$\mathrm{Pic}_{X/k}$ existe. \emph{Lemme (page 130).} De même pour
$X = X_1 \cup X_2$ avec $X_1$, $X_2$ fermés séparables : on se ramène, par
changement de base fidèlement plat, au cas où $X_1 \cap X_2$ a un point
rationnel, et l'on applique la proposition à $X' = X_1 \amalg X_2$.

\subsection*{Le noyau pour un morphisme surjectif (pages 132 et 134)}

\emph{Théorème.} Soit $f \colon X \to Y$ un morphisme surjectif de schémas
propres sur un corps $k$. Le noyau de
$\mathrm{Pic}_{Y/k} \to \mathrm{Pic}_{X/k}$ est extension d'un groupe fini par
un groupe affine ; si les schémas de Picard existent, il est affine. Le
morphisme induit sur les groupes de Néron--Severi modulo torsion est injectif,
et les parties abéliennes (définies à isogénie près) se correspondent par un
morphisme à noyau fini.\footnote{Un point d'interrogation, de sa main, est en
marge de l'énoncé ; la page écrit « quasi-affine » dans le premier corollaire,
ce qui revient au même pour un groupe algébrique. Le sous-indice de
« NeuSev$_0$ » est incertain.}

La démonstration esquissée : a) on suppose $k$ algébriquement clos ; b) on se
ramène à $Y$ réduit, $\mathrm{Pic}_{Y} \to \mathrm{Pic}_{Y_{\mathrm{red}}}$
étant affine ; c) par la factorisation de Stein $X \to Y' \to Y$, avec
$Y' \to Y$ fini surjectif, produit d'épimorphismes stricts finis, pour lesquels
$\mathrm{Pic}_Y \to \mathrm{Pic}_{Y'}$ est affine « par Seshadri », on se
ramène au cas $f_*\mathcal{O}_X \simeq \mathcal{O}_Y$ ; d) où le morphisme est
un monomorphisme. C'est le point (i) de la lettre à Murre du 18 juillet 1962,
sur un corps.\footnote{La page 133 est un feuillet dactylographié d'un autre
texte, dont seule une annotation marginale de sa main est transcrite.}

\subsection*{$\mathrm{Pic}$ relatif et suite de Leray (pages 137 à 139)}

Soit $f \colon X \to S$ se factorisant en $X \xrightarrow{\varphi} \Sigma \xrightarrow{\psi} S$,
$\psi$ fini et $\varphi_*\mathcal{O}_X \simeq \mathcal{O}_\Sigma$. La suite
spectrale de Leray donne une suite exacte
\[
0 \to \mathrm{Pic}(\Sigma) \to \mathrm{Pic}(X) \to \mathrm{Pic}'(X/S) ,
\]
où $\mathrm{Pic}'(X/S)$ désigne les sections globales du faisceau
$R^1 f_*\mathbf{G}_m$ ; la dernière flèche est surjective si $X/\Sigma$ a une
section $g$. En effet, les faisceaux $g$-rigidifiés n'ont pas d'automorphisme
non trivial, et tout $L$ est équivalent au faisceau rigidifié
$L \otimes \varphi^*g^*L^{-1}$.\footnote{La page écrit
$L \otimes \varphi^*(g^*(L))$, sans l'inverse.} Pour $S' \to S$ fidèlement plat
quasi-compact, $\mathrm{Pic}'(X/S) \to \mathrm{Pic}'(X'/S')$ est injectif, et la
suite de descente
$\mathrm{Pic}'(X/S) \to \mathrm{Pic}'(X'/S') \rightrightarrows \mathrm{Pic}'(X''/S'')$
est exacte si $X/\Sigma$ a localement des sections ; d'où l'injectivité de
$\mathrm{Pic}'$ dans le faisceau associé, bijectif sous la même condition. La
page 138, écrite tête-bêche, est un premier essai abandonné.

\pagerange{141}{154}
\section*{Le Picard des schémas en courbes (pages 140 à 154)}

\subsection*{Une obstruction pour une courbe épointée (pages 141 à 144)}

Soient $f \colon X \to S$ projectif et lisse à fibres géométriques connexes de
dimension $1$, $Y \subset X$ un sous-schéma fermé étale sur $S$ (par exemple
une réunion disjointe de sections $g_i(S)$, $i \in I$), $U = X - Y$, et $n$
inversible sur $S$. La suite spectrale de Leray à supports propres
\[
E_2^{pq} = H^p(S, R^q_! f_U\, \mu_n) \Longrightarrow H^*_c(U, \mu_n)
\]
n'a que deux lignes, $R^1_! f_U\mu_n = E[n]$ et $R^2_! f_U\mu_n = \mathbf{Z}/n$,
et la différentielle $d_2 \colon E_2^{0,2} = H^0(S, \mathbf{Z}/n) \to E_2^{2,1}$
fournit un élément canonique $\xi_n \in H^2(S, E[n])$, d'où, pour $n = \ell^\nu$,
une classe $\xi(\ell) \in H^2(S, T_\ell(E))$ ; elle est nulle si $U$ a une
section sur $S$.

Ici $E$ est la jacobienne généralisée de $X$ relativement au module $Y$ : on
contracte $Y$ en $S$, en prenant le conoyau $Z$ de
$X \amalg (Y \times_S Y) \rightrightarrows X$, et l'on pose
$E = \mathrm{Pic}^0_{Z/S}$. C'est un schéma en groupes lisse à fibres
connexes, extension de la jacobienne $J = \mathrm{Pic}^0_{X/S}$ par le tore
$T = (\prod_{Y/S}\mathbf{G}_m)/\mathbf{G}_m$. Le torseur
$P = \mathrm{Pic}^1_{Z/S}$ sous $E$, qui reçoit $U$, a une classe
$\eta \in H^1(S, E)$, et la page écrit
$\xi_n = \partial_n(\eta)$, $\partial_n$ le cobord de la suite de Kummer
$0 \to E[n] \to E \to E \to 0$, avec en marge « à vérifier …… ». L'image de
$\eta$ dans $H^1(S, J)$ est la classe du torseur $\mathrm{Pic}^1_{X/S}$ ; si $X$
a une section, $\eta$ vient de $H^1(S, T)$, et la suite
$0 \to \mathbf{G}_m \to \prod_{Y/S}\mathbf{G}_m \to T \to 0$ donne, quand elle
est scindée, $H^1(S, T) \simeq \mathrm{Pic}(Y)/\operatorname{Im}\mathrm{Pic}(S)$ ;
pour $Y$ réunion de sections, $\mathrm{Pic}(S)^I/\mathrm{Pic}(S)$. Si
$\mathrm{Pic}(S) = 0$, $\eta$ et les $\xi_n$ sont nuls.

\subsection*{Le Picard d'une courbe réduite (pages 146 à 150)}

Soit $C$ une courbe complète, réduite et connexe sur $k$ algébriquement clos,
de composantes $C_i$ ($c_0$ d'entre elles), de normalisée
$\tilde{C} = \coprod \tilde{C}_i$, et de points singuliers $p_\alpha$. La suite
$0 \to \mathcal{O}^*_C \to \tilde{\mathcal{O}}^*_C \to \bigoplus_\alpha \tilde{\mathcal{O}}^*_{p_\alpha}/\mathcal{O}^*_{p_\alpha} \to 0$
donne
\[
0 \to \mathbf{G}_m \to \prod_{C_i}\mathbf{G}_m \to \prod_\alpha A_{p_\alpha} \to \mathrm{Pic}_{C/k} \to \prod_i \mathrm{Pic}_{\tilde{C}_i/k} \to 0 ,
\]
où $A_{p} = \tilde{\mathcal{O}}^*_{p}/\mathcal{O}^*_{p}$ est un groupe
algébrique de dimension $\delta_p = \dim_k \tilde{\mathcal{O}}_p/\mathcal{O}_p$.
La partie linéaire de $\mathrm{Pic}^0_{C/k}$ est
$G = \prod_\alpha A_{p_\alpha} / (\prod \mathbf{G}_m/\mathbf{G}_m)$, de
dimension $\sum_\alpha \delta_\alpha - c_0 + 1$.\footnote{La page note
$\mu_\alpha$ ce que nous notons $\delta_\alpha$, et écrit d'abord
$\mu_\alpha = 1 + \delta_\alpha$ dans un passage qu'elle barre.}

Elle se décompose en partie multiplicative et partie additive :
$\delta_\alpha = \delta'_\alpha + \delta''_\alpha$, où $\delta'_\alpha = r_\alpha - 1$,
$r_\alpha$ le nombre de branches formelles en $p_\alpha$. Alors
\[
\dim G_{\mathrm{mult}} = \sum_\alpha (r_\alpha - 1) - c_0 + 1, \qquad
\dim G_{\mathrm{add}} = \sum_\alpha \delta''_\alpha .
\]
La première est un invariant topologique : c'est le premier nombre de Betti du
graphe qui relie les composantes aux points singuliers par les branches, le
rang du conoyau de $\prod \pi_1(\tilde{C}_i) \to \pi_1(C)$, ce que la page
appelle « le nombre des cycles combinatoires » ; c'est le rang torique. On a
en tout cas
\[
g = \dim \mathrm{Pic}_{C/k} = \sum_i g_i + \rho + \sigma ,
\]
$g_i$ le genre de $\tilde{C}_i$, $\rho$ le rang torique et $\sigma$ la somme des
multiplicités additives.\footnote{La page 150 ajoute à $\rho$ « le nombre des
multiplications multiplicatives des $C_i$ », nul si les $C_i$ sont lisses ;
lecture incertaine.}

$\mathrm{Pic}^0_{C/k}$ n'a pas de partie additive si et seulement si
$\delta_\alpha = r_\alpha - 1$ en chaque point, singularité du type de $r$
droites concourantes en position générale ; il est abélien si de plus le graphe
est un arbre, et alors les $C_i$ sont lisses, puisque
$\mathrm{Pic}^0_{C_i}$ est un quotient de $\mathrm{Pic}^0_C$, et deux
composantes se coupent en au plus un point.\footnote{Les pages 150 et 154
formulent la condition avec des « points doubles ordinaires » ; pour $G = 0$
il suffit que chaque point singulier soit du type de $r$ axes de coordonnées,
$\delta = r - 1$, ce qui inclut les points doubles ordinaires. Le dernier
paragraphe de la page 148 n'est sûr que dans ses formules.}

\subsection*{Courbes stables (pages 152 et 154)}

Les pages 152 et 154 sont couvertes de croquis de courbes à composantes
rationnelles et elliptiques se coupant transversalement, et d'un décompte : la
somme, sur les composantes, de $3g_i - 3 + \nu_i$, $\nu_i$ le nombre de points
doubles sur $C_i$, avec une correction $\lambda$ pour les composantes
rationnelles à deux points doubles. Dans les variables de la page --- $c_0$
composantes, $c_1$ points doubles, $\rho$ cycles, $\sum \nu_i = c_0 - 1 + \rho + c_1$,
$\sum g_i = g - \rho$ --- le total s'écrit
\[
3g - 3 - \bigl[2(c_0 + \rho) - c_1 - \lambda - 2\bigr] = 3g - 3 - (c_1 - \lambda) ,
\]
puisque $\rho = c_1 - c_0 + 1$ pour une courbe à points doubles ordinaires :
chaque point double coûte une dimension. C'est le décompte qu'on fait
aujourd'hui pour la strate de l'espace des courbes stables de genre $g$
formée des courbes de graphe dual donné, de codimension le nombre de points
doubles ; la page ne nomme ni stabilité ni espace de modules, et rien ne date
ce feuillet.\footnote{La vérification de l'égalité est la nôtre ; la page
passe par les lignes $3g - 3\rho - 3c_0 + (c_0 - 1 + \rho) + c_1 + \lambda$ et
$3g - 2\rho - 2c_0 + c_1 + \lambda - 1$, qui sont correctes. Le calcul final
sur une « courbe elliptique $+$ $k$ points », $-6 + 2 - 1 = -5$, $9 - 5 = 4$,
n'est pas assez explicité pour être restitué. La page 153 est une épreuve des
\emph{Éléments} (chap.~I, §~6.4) portant une correction au crayon.}

\pagerange{157}{163}
\section*{Plongements projectifs et foncteur de Picard (pages 156 à 163)}

La chemise porte « Foncteurs de Picard, Diviseurs… ». Soient $X$ projectif sur
$S$, $L$ très ample avec $E = f_*L$ localement libre de rang $N + 1$, et
$i \colon X \to \mathbf{P}(E)$.\footnote{La page dit « de rang $N$ » et
travaille dans $\mathbf{P}^N_S$ ; il faut le rang $N + 1$.} Les pages 157 et
159 écrivent quatre équivalences de données :

A) une base de $E$ ; un isomorphisme $\mathbf{P}^N_S \simeq \mathbf{P}(E)$ et un
isomorphisme du faisceau $\mathcal{O}(1)$ ramené sur $X$ avec $L$ ; un morphisme
$g \colon X \to \mathbf{P}^N_S$ et un isomorphisme $g^*\mathcal{O}(1) \simeq L$,
tels que $H^0$ de $\mathcal{O}(1)$ s'envoie isomorphiquement sur $f_*L$ ---
condition qui ne dépend que de la classe de $g$ ;

A$'$) les mêmes données modulo $\Gamma(S, \mathcal{O}_S^*)$ : une classe de
bases, un isomorphisme $\mathbf{P}^N_S \simeq \mathbf{P}(E)$, un morphisme
$g$ tel que $g^*\mathcal{O}(1) \simeq L$ \emph{relativement à $S$} et
« spécial », c'est-à-dire induisant un isomorphisme sur les $H^0$ ;

A$''$) et A$'''$) les mêmes, quand le faisceau n'est donné que comme section
du foncteur de Picard : une $S$-donnée de faisceau inversible et une
trivialisation projective équivalent à un morphisme spécial
$X \to \mathbf{P}^N_S$ ; une donnée de faisceau « morphique et spécial »
équivaut à un fibré de Brauer--Severi associé à un fibré principal de groupe
$\mathrm{PGL}(N+1)$ et à un morphisme $GP(N)$-équivariant vers lui.

La conclusion, page 159 : $\mathrm{PGL}(N+1)$ opère sans points fixes sur
l'ouvert $\mathrm{Hom}^{\mathrm{spéc}}_S(X, \mathbf{P}^N)$ du schéma des
morphismes, et les données A$'''$) correspondent aux sections du quotient.
C'est la construction, annoncée dans la lettre à Murre des pages 19 à 22, des
grands ouverts de $\mathrm{Pic}_{X/S}$ comme quotients d'un schéma de
plongements par le groupe projectif.\footnote{La page 161 reprend en fragment
la correspondance entre fibrés principaux de groupe $GP(N)$ et morphismes
compatibles. Les marges des pages 157 et 159 portent de longues notes très
raturées. Les pages 158 et 160 sont des épreuves des \emph{Éléments} dont
seules les annotations au crayon, de main non attribuée avec certitude, sont
transcrites.}

\subsection*{Diviseurs et Picard : la suite tangente (page 163)}

Pour un diviseur $Y \subset X$ de faisceau $L = \mathcal{O}_X(Y)$, la suite
$0 \to \mathcal{O}_X \to L \to L|_Y \to 0$ donne
\[
0 \to \mathcal{O}_S \to f_*L \to f_*(L|_Y) \to R^1 f_*\mathcal{O}_X \to R^1 f_*L \to R^1 f_*(L|_Y) \to R^2 f_*\mathcal{O}_X ,
\]
où la page lit, sous les termes, le fibré projectif de $\mathrm{Div}$ sur
$\mathrm{Pic}$, l'espace tangent de $\mathrm{Div}$ en $Y$ et l'espace tangent
de $\mathrm{Pic}$, puis les obstructions de l'un et de l'autre.\footnote{La page
écrit $H^1_f(\mathcal{O}_X)$ pour le dernier terme ; c'est $R^2 f_*\mathcal{O}_X$
que la suite exacte donne. Les annotations sous les derniers termes ne se lisent
que par mots.}

\pagerange{165}{170}
\section*{Pseudo-morphismes et diviseurs (pages 165 à 170)}

La feuille porte en tête « AK », que rien ici n'explique, et « Pseudomorphismes ».
Un ouvert $U$ de $X$ est \emph{schématiquement dense} si aucune section non
nulle de $\mathcal{O}_X$ sur un ouvert $V$ ne s'annule sur $U \cap V$ ; la
notion est locale, stable par intersection finie, et coïncide avec la densité
si $X$ est réduit. Si $Y$ est séparé sur $S$, deux $S$-morphismes $X \to Y$
qui coïncident sur un ouvert schématiquement dense sont égaux. Un
\emph{pseudo-morphisme} est une classe de morphismes définis sur des ouverts
schématiquement denses ; vers un $Y$ séparé il a un plus grand domaine de
définition, avec $\mathrm{Dom}(f|_V) = \mathrm{Dom}(f) \cap V$, et les
pseudo-morphismes forment un faisceau. Pour $X$ réduit, ce sont les
applications rationnelles ; pour $Y = X[t]$, on obtient le faisceau d'anneaux
des pseudo-fonctions, qui contient $\mathcal{O}_X$.

Pages 168 à 170, sur une variété normale irréductible $X$ de dimension $n$,
de corps des fonctions $K$ : l'anneau local d'un diviseur premier $V$ est de
valuation discrète, d'où $\operatorname{div} f = \sum v_V(f)\,V$, somme finie ;
et $\operatorname{div} f \geq 0$ si et seulement si $f$ est partout définie,
parce qu'un anneau normal noethérien est l'intersection de ses localisés en
les premiers de hauteur $1$. L'ensemble des points où $f$ n'est pas définie est
la réunion des $V$ où $v_V(f) < 0$.\footnote{La page note que ces résultats
restent vrais si le lieu singulier est de codimension $\geq 2$. Il faut de plus
que $X$ soit $(S_2)$, c'est-à-dire normal : en pinçant deux points d'un plan,
on obtient une surface dont le lieu singulier est un point, et une fonction
qui vaut $1$ en l'un et $0$ en l'autre n'est pas définie au point pincé bien
que son diviseur soit positif.}

\emph{Théorème (page 169).} Une application rationnelle d'une variété normale
de dimension $n$ vers une variété complète est définie hors d'une partie de
dimension $\leq n - 2$. Deux démonstrations : a) le graphe donne un morphisme
birationnel propre $g$ vers $X$ ; le lieu $X'$ des points à fibre de dimension
$> 0$ est de dimension $\leq n - 2$, sinon $g^{-1}(X')$ serait de dimension
$n$, et hors de $X'$, $g$ est un isomorphisme par le théorème principal de
Zariski ; b) en chaque diviseur premier, l'anneau local est de valuation
discrète, et le critère valuatif des variétés complètes donne une localité de
la variété but dominée par lui. La seconde n'utilise que la régularité en
codimension $1$, et c'est pour elle que la généralisation notée en marge est
juste.

\pagerange{172}{182}
\section*{Le foncteur de Picard (pages 172 à 182)}

Ces pages définissent l'objet du dossier. Pour $X$ sur $S$ et
$X' = X \times_S S'$, on pose
\[
\mathrm{Pic}^{\mathrm{pré}}_{X/S}(S') = \mathrm{Pic}(X') = H^1(X', \mathcal{O}^*_{X'}),
\qquad
\mathrm{Pic}'_{X/S}(S') = \Gamma\bigl(S', R^1 f'_*\mathcal{O}^*_{X'}\bigr) ,
\]
et l'on passe du premier au second en remplaçant la topologie grossière par
celle des recouvrements de Zariski.\footnote{La page appelle ces foncteurs
$\underline{\underline{\text{Prépic}}}_{X/S}$ et
$\underline{\underline{\text{Prépic}}}'_{X/S}$ ; « Pic abs » y désigne le
groupe de Picard absolu de la base.}

\emph{Proposition 1.} Si $f_*\mathcal{O}_X = \mathcal{O}_S$, il y a un
monomorphisme canonique $\mathrm{Pic}(X)/\mathrm{Pic}(S) \to \mathrm{Pic}'(X/S)$,
isomorphisme si $X/S$ a une section. En effet, le foncteur $F \mapsto f^*F$
des faisceaux localement libres sur $S$ vers ceux de $X$ est pleinement
fidèle, d'inverse $f_*$ sur son image, et la suite exacte des termes de bas
degré de Leray s'écrit
\[
0 \to H^1(S, \mathcal{O}^*_S) \to H^1(X, \mathcal{O}^*_X) \to H^0(S, R^1 f_*\mathcal{O}^*_X) \to H^2(S, \mathcal{O}^*_S) \to H^2(X, \mathcal{O}^*_X) .
\]
Avec une section $s$, la catégorie des faisceaux inversibles $L$ munis d'un
isomorphisme $\mathcal{O}_S \simeq s^*L$ n'a pas d'automorphismes non triviaux,
et tout élément de $\mathrm{Pic}'(X/S)$ provient d'un tel $L$ unique à
isomorphisme unique près : on remplace $L_i$ par
$L_i \otimes f^*s^*L_i^{-1}$ sur les ouverts d'un recouvrement et l'on recolle.
« Question : donner une formulation cohomologique du raisonnement. »

\emph{Proposition 2.} Si de plus $X$ est quasi-compact et séparé sur $S$ et
$S' \to S$ fidèlement plat et quasi-compact, alors
$\mathrm{Pic}'(X/S) \to \mathrm{Pic}'(X'/S')$ est injectif, et la suite
$\mathrm{Pic}'(X/S) \to \mathrm{Pic}'(X'/S') \rightrightarrows \mathrm{Pic}'(X''/S'')$
est exacte si $X/S$ a localement des sections. Le point clé est qu'un faisceau
localement libre sur $X$ provient de $S$ si et seulement si son image inverse
sur $X'$ provient de $S'$, puisque $f_*$ commute au changement de base plat.
Si $f_*\mathcal{O}_X = \mathcal{O}_S$ universellement et que $X/S$ a localement
des sections, $\mathrm{Pic}'_{X/S}$ est un faisceau pour la topologie
fidèlement plate quasi-compacte.

On \emph{définit} alors $\mathrm{Pic}_{X/S}$ comme le faisceau associé à
$\mathrm{Pic}^{\mathrm{pré}}_{X/S}$ pour cette topologie, avec des morphismes
\[
\mathrm{Pic}^{\mathrm{pré}}_{X/S} \xrightarrow{\ \alpha\ } \mathrm{Pic}'_{X/S} \xrightarrow{\ \beta\ } \mathrm{Pic}_{X/S} .
\]
La page signale un « ennui » : un recouvrement de Zariski d'un schéma non
quasi-compact n'est pas majoré par un morphisme fidèlement plat
quasi-compact, et il faut parler de morphismes de localisation fidèlement
plate --- c'est la précaution qu'on prend aujourd'hui dans la définition de la
topologie fpqc. $\alpha$ n'est pas injectif en général (son noyau est le
Picard de la base), et surjectif si $X/S$ a une section ; $\beta$ est un
monomorphisme si $f_*\mathcal{O}_X = \mathcal{O}_S$ universellement, et
surjectif si $X$ a localement des sections. Exemple où $\beta$ n'est pas
surjectif : une variété de Brauer--Severi non triviale, l'obstruction étant
dans $H^2(S, \mathbf{G}_m)$. Ainsi $\mathrm{Pic}^{\mathrm{pré}}$ n'est jamais
représentable sauf cas triviaux, $\mathrm{Pic}'$ ne peut l'être que si $\beta$
est un isomorphisme, et $\mathrm{Pic}_{X/S}$ l'est « dans des cas suffisamment
importants ».

\subsection*{Diviseurs relatifs (pages 179 à 182)}

Un sous-schéma divisoriel $Y$ de $X$ transversal aux fibres est défini par un
faisceau inversible $L$ et une section de $L$ non diviseur de zéro sur les
fibres. Les $Y$ dont le faisceau est isomorphe à un $L$ donné correspondent à
ces sections modulo $\Gamma(S, f_*\mathcal{O}^*_X)$, donc modulo
$\mathcal{O}^*_S$ si $f_*\mathcal{O}_X = \mathcal{O}_S$ ; à équivalence près
sur $S$, on obtient les sections de $(f_*L)^\circ/\mathcal{O}^*_S$, et de même
pour la Picard-équivalence si $X$ est quasi-compact et séparé et
$f_*\mathcal{O}_X = \mathcal{O}_S$ universellement. Si $\mathrm{Pic}_{X/S}$ et
$\mathrm{Div}_{X/S}$ existent, avec $p \colon \mathrm{Div} \to \mathrm{Pic}$,
ces sections sont celles de la fibre $\mathrm{Div}^L_{X/S} = p^{-1}(\lambda)$
au-dessus de la section $\lambda$ définie par $L$ ; et si l'on dispose d'un
faisceau de Poincaré $\Lambda$ sur $\mathrm{Pic}_{X/S} \times_S X$ (par exemple
si $X/S$ a une section),
$\mathrm{Div}_{X/S} \simeq \mathrm{Div}^{\Lambda}_{\mathrm{Pic}_{X/S} \times_S X / \mathrm{Pic}_{X/S}}$,
ce qui permet d'étudier le schéma des diviseurs indépendamment de celui de
Picard.

\emph{Proposition (page 182).} Soit $X$ propre et plat sur $S$, à fibres
géométriquement intègres, et supposons
$f_*(L \otimes_{\mathcal{O}_S} M) \simeq \mathcal{H}om(Q, M)$ fonctoriellement
en $M$ quasi-cohérent, pour un $\mathcal{O}_S$-module cohérent $Q$ (ce qui est
vrai si $S$ est local et $X$ projectif). Alors
$\mathrm{Div}^L_{X/S} \simeq \mathbf{P}(Q) = \operatorname{Proj}(\operatorname{Sym} Q)$.\footnote{La
page corrige « intègres » en « primaires » ; c'est l'intégrité qui sert. Sur la
droite double $y^2 = 0$ du plan, fibre primaire, la section $y$ de
$\mathcal{O}(1)$ est nilpotente et ne définit pas de diviseur. La page note
$\underline{E}$ le module que nous notons $Q$.} D'où le \emph{corollaire} :
les conditions suivantes sont équivalentes : $\mathrm{Div}^L_{X/S}$ est plat
sur $S$ en $s$ ; il est lisse en $s$ ; $\operatorname{Sym}^n Q$ est libre en $s$
pour $n$ grand ; $Q$ est libre en $s$ ; $L$ est cohomologiquement plat
relativement à $S$ ; $f_*L \otimes k(s) \to H^0(X_s, L_s)$ est surjectif.

\pagerange{184}{193}
\section*{Le Picard des cônes et de leurs anneaux locaux (pages 183 à 193)}

\subsection*{Le cône épointé}

Soient $X \subset \mathbf{P}^r_k$, $C$ le cône affine de sommet $a$, et
$C_a = \operatorname{Spec}\mathcal{O}_{C,a}$. $C - a$ est un
$\mathbf{G}_m$-torseur sur $X$ de classe $\xi$, celle de $\mathcal{O}_X(1)$ ;
donc, si $X$ est régulier,
$\mathrm{Pic}(C - a) \simeq \mathrm{Pic}(X)/\mathbf{Z}\xi$. Grothendieck
affirme que $\mathrm{Pic}(C - a) \to \mathrm{Pic}(C_a - a) = \varinjlim_{U \ni a}\mathrm{Pic}(U - a)$
est un isomorphisme si $X$ est régulier ; cela revient à dire qu'un diviseur
$D$ de $C$ ne passant pas par $a$ est équivalent à $0$ sur $C - a$.

On passe au cône projectif $\hat{C}$, et à $E = \hat{C} - a$, l'espace total
d'un fibré en droites sur $X$. \emph{Lemme.} Si $X$ est régulier, $E$ l'espace
total d'un fibré en droites sur $X$ de section nulle $g$, $V = E - g(X)$, et
$D$ un diviseur sur $E$ à support propre sur $X$, alors $D|_V \sim 0$. Si $D$
est irréductible, défini par $X' \subset E$ régulier, alors $h \colon X' \to X$
est fini de degré $d$, le fibré $E' = E \times_X X'$ a une section
tautologique $q$, et, en calculant dans les groupes de classes de cycles,
\[
D = h'_*(q_*1) = h'_*\bigl(f'^*h^*\xi_E\bigr) = h'_*h'^*\bigl(f^*\xi_E\bigr) = d\, f^*\xi_E ,
\]
qui est nul sur $V$. Le cas où $X'$ est singulier se traite en retirant
l'image de son lieu singulier.\footnote{La fin de ce raisonnement est
illisible (« et on gagne … ») ; la réduction est seulement indiquée. La page
écrit « immersion canonique de $X'$ dans $X$ », lire « dans $E$ ». La marge
note que la méthode se généralise aux fibrés vectoriels quelconques.} D'où,
encadré,
\[
\mathrm{Pic}(C - a) \simeq \mathrm{Pic}(C_a - a) \simeq \mathrm{Pic}(X)/\mathbf{Z}\xi ,
\]
qui est le groupe des classes de diviseurs $\mathrm{Cl}(\mathcal{O}_{C,a})$
quand $C$ est normal.

\subsection*{Comparaison avec le complété}

Soient $S = C_a$ et $\hat{S}$ le spectre du complété. On se demande si
$\mathrm{Pic}(S - a) \to \mathrm{Pic}(\hat{S} - \hat{a})$ est un isomorphisme.
L'éclaté $\tilde{C}$ de l'origine est le fibré en droites
$E = \mathbb{V}(\mathcal{O}_X(1))$, dont la section nulle $X$ a pour idéal
conormal $\mathcal{J}/\mathcal{J}^2 = \mathcal{O}_X(1)$ ; par le théorème
d'existence, l'éclaté de $\hat{S}$ est le complété formel $\hat{E}$ de $E$ le
long de $X$, et l'homomorphisme étudié est induit, modulo les classes du
diviseur exceptionnel, par
\[
\mathrm{Pic}(X) \simeq \mathrm{Pic}(E) \longrightarrow \mathrm{Pic}(\hat{E}) = \varprojlim_n \mathrm{Pic}(E_n),
\]
$E_n$ le $n$-ième voisinage infinitésimal de $X$. Le composé avec
$\mathrm{Pic}(\hat{E}) \to \mathrm{Pic}(X)$ est l'identité : l'homomorphisme
est injectif, et bijectif si et seulement si $\mathrm{Pic}(\hat{E}) \to \mathrm{Pic}(X)$
est injectif. Le noyau de $\mathrm{Pic}(E_n) \to \mathrm{Pic}(E_{n-1})$ est
$H^1(X, \mathcal{J}^n/\mathcal{J}^{n+1}) = H^1(X, \mathcal{O}_X(n))$.\footnote{Le
diagramme de la page 188 est lu sur une page chargée ; le sens de ses flèches
reste douteux, et un troisième terme y est barré. L'ordre de lecture du passage
serré de la page 190 est une conjecture de la transcription.}

\emph{Proposition (page 191).} Pour que
$\mathrm{Pic}(S - a) \to \mathrm{Pic}(\hat{S} - \hat{a})$ soit bijectif, il
suffit que $H^1(X, \mathcal{O}_X(n)) = 0$ pour tout $n \geq 1$ ; si
$\dim X = 1$, c'est nécessaire, car alors les transitions sont surjectives.

Pour une courbe plane lisse de degré $d$, de genre
$g = 1 + \frac{d(d-3)}{2}$, $H^1(X, \mathcal{O}_X(1))$ est dual de
$H^0(X, \mathcal{O}_X(d - 4))$, non nul dès que $d \geq 4$ : la condition n'est
vérifiée que pour $d \leq 3$ --- en particulier pour la cubique, c'est-à-dire le
cône sur une courbe elliptique du tapuscrit des pages 2 à 12.\footnote{La page
donne un critère suffisant plus grossier : la condition échoue si
$\deg \mathcal{O}_X(1) < g - 1$, ce qui pour une courbe plane revient à
$d \geq 6$. C'est exact, mais l'échec commence à $d = 4$.}

\emph{Remarques (pages 192 et 193).} Si $\dim X = 1$, le noyau de
$\mathrm{Pic}(\hat{E}) \to \mathrm{Pic}(X)$ est connexe et unipotent, et
l'homomorphisme induit un isomorphisme sur les groupes des composantes, finis,
isomorphes à $\mathbf{Z}/d\mathbf{Z}$, $d$ le degré de $X$. Le résultat
devrait valoir en dimension quelconque en caractéristique $p > 0$ « à
pro-isomorphisme près », les obstructions dans les $H^2(X, \mathcal{O}_X(n))$
s'annulant pour $n$ grand et étant tuées par $p$ : le groupe de Néron--Severi
local serait le même pour $S$ et $\hat{S}$ modulo des $p$-groupes, « en accord
avec la théorie du groupe fondamental ». Enfin, avec le hensélisé $S^h$, on a
$\mathrm{Pic}(S - a) \to \mathrm{Pic}(S^h - a) \to \mathrm{Pic}(\hat{S} - \hat{a})$,
et la seconde flèche est peut-être toujours un isomorphisme, ce qu'on voit en
dimension $2$ par l'éclatement ; les courbes planes de degré $\geq 4$
donnent alors des exemples où le groupe de Picard local d'un anneau et celui de
son hensélisé diffèrent.\footnote{La page note le hensélisé d'un signe que la
transcription lit comme un tilde, déjà employé page 188 pour l'éclaté ; nous
écrivons $S^h$. La page dit « dans l'exemple précédent des courbes » ; nous
reportons sur $d \geq 4$, en accord avec la note précédente.}

\pagerange{195}{197}
\section*{Structure locale de $\mathrm{Pic}_{X/S}$ : les constituants de la torsion (pages 194 à 197)}

La chemise porte « $G^\tau$, $G^0$ etc » et « Structure locale de
$\mathrm{Pic}_{X/S}$ ». Soient $G$ un schéma en groupes commutatif de type fini
sur $S$ noethérien local, $s$ le point fermé, $s_i$ ses générisations
immédiates. Pour $t \in S$, on note $\sigma(t)$ l'ordre du groupe fini
$G_t/G_t^0$. Les conditions considérées sont : (i) $G$ est plat et $G^0$
contient un sous-schéma abélien $A$ de même ensemble sous-jacent ; (ii) un tel
$A$ existe et $G/A$, fini, est plat ; (iii) $\sigma$ est constante ; (iv)
$\sigma(s) = \sigma(s_i)$ pour tout $i$. On a (i) $\Leftrightarrow$ (ii) par la
théorie des quotients plats ($G \to G/A$ est fidèlement plat et projectif), et
(ii) $\Rightarrow$ (iii) $\Rightarrow$ (iv), avec (iii) $\Leftrightarrow$ (iv)
quand les fibres sont toutes de même dimension ; si de plus $S$ est réduit,
toutes sont équivalentes.

\emph{Lemme.} $\sigma$ est semi-continue supérieurement. Elle est
constructible, et sur un anneau de valuation discrète, si $N$ est
l'adhérence de la composante neutre générique --- un sous-groupe plat --- et
$Q = G/N$, fini, on a
\[
\sigma(s) \geq \operatorname{rg} Q_s \geq \operatorname{rg} Q_{s'} = \sigma(s') ,
\]
la première inégalité étant une égalité si et seulement si $N$ est un schéma abélien,
la seconde si et seulement si $Q$ est plat. Ceci prouve (iv) $\Rightarrow$ (ii)
sur un anneau de valuation discrète, et, par une note marginale qui passe par
le schéma de Hilbert et le critère valuatif, sur une base réduite.\footnote{La
page écrit dans (iv), entre crochets, « on sait $\sigma(s) \leq \sigma(s_i)$ »,
inégalité inverse de celle que le lemme démontre. C'est le lemme qui a raison :
au point spécial, le groupe des composantes ne peut que grossir. La note
marginale n'est lue que par fragments.}

\emph{Corollaire (page 197).} Si les fibres sont connexes et de même dimension,
ces conditions équivalent à l'égalité des rangs des $G[n]$ dans les fibres,
pour $n$ un multiple fixé convenable des $\sigma(t)^2$. Un encadré décompose
enfin, pour un schéma de Picard en caractéristique $p$, la longueur de la
$p^n$-torsion en trois termes, étale, multiplicatif et unipotent, interprétés
par $\mathrm{Hom}(\pi_1(\bar{X}_s), \mathbf{Z}/p^n)$, la $p^n$-torsion de
$\mathrm{Pic}(\bar{X}_s)$, et le noyau de $F^n$ sur $H^1(\bar{X}_s, W_n)$.\footnote{L'encadré
est lu avec ses renvois, mais le rattachement des trois interprétations aux
trois termes n'est pas sûr ; on n'en donne que la forme.}

\pagerange{200}{211}
\section*{Composante neutre, sous-groupe $G^\tau$, partie constructible (pages 200 à 211)}

Ces pages, au crayon et très rapides, ne se lisent que partiellement ; on
donne les énoncés que la transcription porte.

\subsection*{Sur un corps (pages 200 à 205)}

Soit $G$ un schéma en groupes localement de type fini sur un corps $k$.
Alors : ses composantes connexes sont de type fini et irréductibles ; la
composante neutre $G^0$ est géométriquement irréductible et c'est un
sous-groupe distingué ; une composante qui contient un point rationnel $a$ est
$G^0 a$ ; le quotient $G/G^0$ existe, étale, et $G \to G/G^0$ est fidèlement
plat ; plus généralement (corollaire de la page 202), $G/H$ existe pour tout
sous-groupe distingué $H$, et $G \to G/H$ est fidèlement plat. Démonstration, pour $k$ algébriquement clos : $G^0 \times G^0$ est connexe,
donc $G^0$ est un sous-groupe ; pour $U$ un ouvert dense de $G^0$ voisinage de
$e$, $U \cdot U = G^0$, car si $x \in G^0(k)$ les ouverts non vides $U$ et
$xU^{-1}$ de l'irréductible $G^0$ se rencontrent ; donc $G^0$ est
quasi-compact, de type fini ; il est irréductible, sinon tout point fermé
serait, par translation, sur deux composantes irréductibles.

Pour $G$ commutatif, soit $\Gamma = G/G^0$, groupe discret. Les noyaux
$\Gamma[n]$ sont ouverts et fermés, et leur réunion $\Gamma^\tau$ est un
sous-groupe ouvert, l'ensemble des points dont un multiple est neutre ; sur un
$S$ quelconque, $\Gamma^\tau(S)$ est la torsion de $\Gamma(S)$. Elle est la
réunion des sous-groupes finis de $\Gamma$, et le plus petit sous-groupe dont
le quotient soit sans torsion.\footnote{La page dit, pour $\Gamma$ commutatif,
que $\Gamma^\tau$ est « le plus grand sous-groupe … fini », avec deux mots
illisibles ; $\Gamma^\tau$ n'est pas fini en général (le groupe constant
$\mathbf{Q}/\mathbf{Z}$), d'où notre formulation.} On définit
$G^\tau$ comme l'image inverse de $\Gamma^\tau$ : c'est l'ensemble des points
dont un multiple est dans $G^0$, et le plus petit sous-groupe ouvert tel que
$G/G^\tau$ soit discret sans torsion. C'est le \emph{sous-groupe de torsion} de
$G$ au sens de la théorie de Picard. La page 205 note que la stationnarité
locale des $\Gamma[n]$ vaut pour tout groupe localement fini, séparable ou non,
mais que le cas radiciel est moins évident.

\subsection*{Un critère de quasi-compacité (pages 207 et 208)}

\emph{Proposition.} Soient $f \colon X \to Y$ un morphisme séparé, localement
de présentation finie et universellement ouvert, $Y$ localement noethérien,
dont les fibres sont propres. Alors $f$ est quasi-compact.\footnote{Les
hypothèses sont lues avec des lacunes, dont la première, illisible. Une
hypothèse d'ouverture est nécessaire : la réunion disjointe des points fermés
de la droite affine, au-dessus de cette droite, a des fibres propres (un point)
et n'est pas quasi-compacte. C'est l'ouverture qu'utilise la démonstration.} On
se ramène, par récurrence noethérienne, à trouver un ouvert non vide de $Y$
au-dessus duquel $f$ est quasi-compact, $Y$ intègre de point générique $\eta$.
Soit $V$ un ouvert quasi-compact contenant la fibre générique $X_\eta$, propre ;
quitte à restreindre $Y$, $V$ est propre sur $Y$, donc fermé dans $X$, et
$X - V$ est ouvert ; s'il était non vide, il rencontrerait la fibre générique
par ouverture, alors que $(X - V)_\eta = \emptyset$. Donc $X = V$. Le corollaire
de la page 208 énonce la même chose pour une partie localement constructible
$Z$ de $X$ dont chaque point a des générisations dans $Z$ au-dessus de toute
générisation de son image --- ce qui remplace l'ouverture --- et dont les fibres
sont propres.

\subsection*{La partie première à $p$ de $\mathrm{Pic}^\tau$ (pages 209 à 211)}

\emph{Proposition.} Soient $X$ projectif sur $S$ noethérien,
$f_*\mathcal{O}_X = \mathcal{O}_S$ universellement, à fibres primaires,
$P = \mathrm{Pic}_{X/S}$, et $P'$ la réunion des $P'_s$, où $P'_s$ est
l'ensemble des éléments de $P^\tau_s$ dont la classe dans $P^\tau_s/P^0_s$ est
d'ordre premier à la caractéristique de $k(s)$. Si l'adhérence de tout point
de $P^\tau$ est quasi-compacte (ou si les $P^0_s$ sont propres), $P'$ est une
partie constructible et quasi-compacte de $P$.

Démonstration : par récurrence noethérienne, il suffit de le voir au-dessus
d'un ouvert dense d'une base intègre de point générique $\eta$. $P'_\eta$ est un
sous-groupe ouvert de type fini de $P_\eta$ ; il s'étend en un sous-groupe
ouvert $H$ de $P^\tau$ au-dessus d'un voisinage, fermé quitte à restreindre la
base. Il reste à voir que $H = P'$ : comme $P^0_s$ est divisible par les
entiers premiers à la caractéristique, il suffit que $P_s[n] \subset H$ pour
$n$ inversible. Or $P_s[n] = H^1(X_s, \mu_n)$, et au-dessus de l'ouvert où $n$
est inversible, le foncteur $H^1(X/S, \mu_n)$ est étale sur $S$ : tout point de
sa fibre spéciale est spécialisation d'un point de $P_\eta[n] \subset H$, fermé.

\emph{Corollaires.} Si les $P^0_s$ sont propres et les groupes
$P^\tau_s/P^0_s$ d'ordre premier aux caractéristiques résiduelles, alors
$P^\tau$ est de type fini sur $S$. Si la réunion $P^0$ des $P^0_s$ est fermée,
$P^\tau$ est fermé sous la même condition d'ordre.

\pagerange{213}{223}
\section*{Deux exemples où le schéma de Picard se dégrade (pages 213 à 223)}

\subsection*{Une conique qui dégénère en deux droites (pages 213 à 218)}

Soient $S = \operatorname{Spec} k[t]$ et
$X = \operatorname{Proj} k[t][x, y, z]/(xy - tz^2)$, gradué en $x, y, z$.
L'idéal est premier, $X$ est régulier et plat sur $S$. Au-dessus de
$U = S - \{0\}$, $X_U \simeq \mathbf{P}^1_U$ ; la fibre spéciale $X_0$,
$xy = 0$, est la réunion de deux droites $Y_1$, $Y_2$ se coupant en un point. Les
fibres ont $H^i(\mathcal{O}) = 0$ pour $i > 0$, et pour $S'$ local sur $S$,
$\mathrm{Pic}(X') \to \mathrm{Pic}(X'_0)$ est bijectif ; ce groupe vaut
$\mathbf{Z}$ si le point fermé de $S'$ n'est pas au-dessus de $0$, et
$\mathbf{Z} \times \mathbf{Z}$ (les bidegrés) sinon.\footnote{La page écrit
la flèche $(*)$ vers $H^1(X'_0, \mathcal{O}_{X'_0})$ ; il faut lire
$\mathcal{O}^*_{X'_0}$, comme le confirme le calcul qui suit, et
l'injectivité est énoncée d'abord, la bijectivité ensuite (« Lemme e »).}

\emph{Lemme c).} Soient $D_i$ les diviseurs de Cartier $Y_i$ ($X$ est régulier)
et $L_i = \mathcal{O}_X(D_i)$. Alors $L_i|_{Y_i}$ est de degré $-1$ et
$L_i|_{Y_j}$ de degré $1$ pour $i \neq j$ : $D_1 + D_2 = \operatorname{div}(t) \sim 0$
et $D_1 \cdot D_2 = 1$, les deux droites se coupant transversalement au point où
$X$ est localement $xy = t$.\footnote{La page conclut « $L_{D_1}|D_1 = 1$ », là
où ce qui précède donne $-1$.}

Deux sections $s_1$, $s_2$ de $X/S$, données dans la carte affine
$xy = t$ par $(x, y) = (1, t)$ et $(t, 1)$, rencontrent chacune une seule des
deux droites ; leurs diviseurs donnent des faisceaux de bidegrés $(1, 0)$ et
$(0, 1)$. D'où une suite exacte
\[
0 \to \mathbf{Z} \xrightarrow{\ \alpha\ } \mathrm{Pic}(X) \xrightarrow{\ \beta\ } \mathbf{Z} \to 0 ,
\]
$\beta$ le degré sur la fibre générique et $\operatorname{Ker}\beta$ engendré
par la classe de $D_1$ (puisque $\mathrm{Pic}(X_U) = \mathrm{Pic}(\mathbf{P}^1_U) = \mathbf{Z}$,
$k[t, t^{-1}]$ étant principal), et la restriction
$\mathrm{Pic}(X) \to \mathrm{Pic}(X_0) \simeq \mathbf{Z} \times \mathbf{Z}$ est
bijective.

Le « Lemme f » de la page 216 affirme que le foncteur de Picard n'est pas
représentable par un schéma localement de type fini ; sa démonstration, qui
déborde sur la page 217, est barrée de diagonales. La page 218 fait le
contraire : soit, pour un ensemble $I$, le schéma $S_I$ obtenu en recollant
des copies $S_i$ de $S$, $i \in I$, le long de $U$ --- la droite dont l'origine
est dédoublée en autant de points qu'il y a d'indices. Pour $n \in \mathbf{Z}$,
soit $I_n$ l'ensemble des couples $(p, q)$ avec $p + q = n$, et
\[
\mathfrak{P} = \coprod_{n \in \mathbf{Z}} S_{I_n} .
\]
Sa fibre en $0$ est $\mathbf{Z} \times \mathbf{Z}$, sa fibre ailleurs
$\mathbf{Z}$, et $\mathfrak{P}$ \emph{représente} le foncteur de Picard. Le
schéma de Picard existe donc, localement de type fini, mais n'est pas séparé,
et ses composantes connexes ne sont pas quasi-compactes : un faisceau de degré
$n$ sur la conique générale a pour limites tous les bidegrés $(p, q)$ avec
$p + q = n$.\footnote{L'énoncé du Lemme f n'est pas barré ; seule sa
démonstration l'est. Nous retenons ce que la page 218 établit. La construction
de $S_I$ part d'un ensemble « $U \times \{I\} \amalg I$ » où l'on lit que la
partie $U$ est commune à tous les $S_i$.}

La page 219 (« Autre exemple ») calcule le Picard d'un cône projectif : pour
$X$ projectif sur $S$ local noethérien, engendré en degré $1$, avec
$f_*\mathcal{O}_X = \mathcal{O}_S$, et $\hat{C}$ le cône projectif, l'éclaté
$\hat{C}_X \to \hat{C}$ du sommet vérifie $g_*\mathcal{O} = \mathcal{O}$, donc
$\mathrm{Pic}(\hat{C}) \to \mathrm{Pic}(\hat{C}_X)$ est injectif ; une classe
dans l'image est triviale sur la section nulle, donc de la forme
$\mathcal{O}(n)$,\footnote{La page marque trois étapes « vérif ».} et
\[
\mathrm{Pic}(\hat{C}) \simeq \mathbf{Z}, \quad \text{engendré par } \mathcal{O}_{\hat{C}}(1) .
\]
La page 220 l'applique à la famille de quadriques
$x^2 + y^2 + z^2 - \lambda t^2 = 0$ sur la droite des $\lambda$, qui dégénère en
un cône quadratique en $\lambda = 0$ : sur un anneau local où $\lambda$ n'est
pas inversible, le Picard vaut $\mathbf{Z}$, engendré par $\mathcal{O}(1)$,
grâce à $H^1(X'_0, \mathcal{O}) = 0$ ; là où $\lambda$ est inversible, il vaut
$\mathbf{Z}^2$, engendré par les deux systèmes de génératrices (après
extension convenable). Le rang du Picard des fibres géométriques saute de $2$
à $1$ à la dégénérescence.\footnote{La conclusion de la page sur le schéma de
Picard ainsi obtenu, et sur sa partie de type fini, est illisible. En
marge : « le schéma de Picard d'une [quadrique] est $\mathbf{Z}$ ».}

\subsection*{Trois droites (pages 221 à 223)}

Sur $\mathbf{Z}[t]$, la famille de courbes planes $xy(x + y + tz) = 0$ :
trois droites formant un triangle pour $t$ inversible, concourantes pour
$t = 0$. La famille est plate, ses composantes $X_i$ sont des $\mathbf{P}^1_S$,
les intersections deux à deux sont des sections, et l'intersection triple est
la fibre $t = 0$. La page entreprend de calculer $H^1(X, \mathcal{O}_X)$ par
descente le long de $X_1 \amalg X_2 \amalg X_3 \to X$ et s'arrête là. Avec les
formules des pages 146 à 150, la fibre générique a pour $\mathrm{Pic}^0$ un
tore $\mathbf{G}_m$ (un cycle), la fibre spéciale un groupe additif
$\mathbf{G}_a$ (un point triple plan, $\delta = 3$, trois branches) : c'est
vraisemblablement ce que l'exemple doit montrer.\footnote{Cette dernière phrase
est notre ajout ; la page s'arrête sur « $\simeq$ ».}

\pagerange{226}{245}
\section*{Picard formel et analytique : la lettre à Hironaka (pages 225 à 245)}

La chemise porte « Schémas de Picard des schémas formels en car.~0, ou des
schémas propres sur une algèbre analytique sur $\mathbf{C}$ ». Elle contient
une lettre dactylographiée en anglais, « Neuilly July 6 1962 », corrigée à
l'encre, puis des notes.

\subsection*{La lettre}

Elle suit une conversation du mardi précédent. \emph{Théorème.} Soient
$f \colon X \to Y$ un morphisme propre d'espaces analytiques complexes, $y \in Y$,
$Y_n$ le spectre analytique de $\mathcal{O}_y/\mathfrak{m}_y^{n+1}$,
$X_n = X \times_Y Y_n$,
$\mathrm{Pic}(X_y) = (R^1 f_*\mathcal{O}^*_X)_y = \varinjlim_{U \ni y}\mathrm{Pic}(f^{-1}(U))$
et $\mathrm{Pic}(\hat{X}_y) = \varprojlim \mathrm{Pic}(X_n)$, avec
$u \colon \mathrm{Pic}(X_y) \to \mathrm{Pic}(\hat{X}_y)$ et
$v_n \colon \mathrm{Pic}(\hat{X}_y) \to \mathrm{Pic}(X_n)$. Alors (i) le système
$(\mathrm{Pic}(X_n))_n$ satisfait la condition de Mittag-Leffler, et même
uniformément, au sens d'Artin--Rees ; (ii) $\operatorname{Im} v_n = \operatorname{Im} v_n u$
est l'ensemble des images universelles de $\mathrm{Pic}(X_n)$ ; (iii) $u$ est un
isomorphisme si et seulement si $(R^1 f_*\mathcal{O}_X)_y$ est de longueur
finie. Plus précisément (i bis), dans le système des groupes analytiques
$\mathrm{Pic}_{X_n/\mathbf{C}}$, les groupes de Néron--Severi sont constants pour
$n$ grand, et noyaux et conoyaux des transitions sont des groupes vectoriels.

\emph{Corollaire 1.} $X$ est projectif au-dessus d'un voisinage de $y$ si et
seulement si chaque $X_n$ est projectif ; c'est le cas si $\dim X_0 \leq 1$.

La démonstration n'utilise que les théorèmes de Grauert de finitude et de
comparaison pour les images directes et la suite exponentielle
$0 \to \mathbf{Z} \to \mathcal{O} \to \mathcal{O}^* \to 0$ ; elle vaut pour tout
$H^i(\mathcal{O}^*)$. Pour un schéma formel propre sur un anneau local complet
de corps résiduel de caractéristique $0$, l'analogue de (i) et (i bis) est
vrai, avec une démonstration purement algébrique fondée sur l'absence de torsion
des noyaux et conoyaux de $\mathrm{Pic}_{X_{n+1}} \to \mathrm{Pic}_{X_n}$ et sur
le théorème de finitude de Néron ; tout cela tombe en caractéristique $> 0$.

Par le GAGA de Serre--Grauert--Remmert--Grothendieck, encore non écrit (« as yet
unwritten ! »), le théorème donne une propriété intrinsèque des algèbres
analytiques. \emph{Théorème.} Soient $A$ une algèbre analytique sur $\mathbf{C}$,
$\hat{A}$ son complété, $Y$, $\hat{Y}$ les spectres, $X$ propre sur $Y$ et
$\hat{X} = X \times_Y \hat{Y}$. Alors $(\mathrm{Pic}(X_n))$ satisfait
Mittag-Leffler--Artin--Rees ; $\mathrm{Pic}(X)$ et $\mathrm{Pic}(\hat{X})$ ont
même image dans $\mathrm{Pic}(X_n)$ ; et
$\mathrm{Pic}(X) \to \mathrm{Pic}(\hat{X})$ est un isomorphisme si et seulement
si $H^1(X, \mathcal{O}_X)$ est de longueur finie sur $A$. Corollaire : $X/Y$
est projectif si et seulement si $\hat{X}/\hat{Y}$ l'est, si et seulement si
chaque $X_n/Y_n$ l'est.

\emph{Corollaire 2.} Avec la résolution des singularités de Hironaka et la
méthode de Mumford qui relie le groupe de Picard local de $A$ à celui d'un
schéma régulier qui le domine : si $Y' = Y - \{y\}$ est régulier, alors
$\mathrm{Pic}(Y') \to \mathrm{Pic}(\hat{Y}')$ est un isomorphisme. Cela
explique pourquoi Mumford a pu mettre sur le groupe des classes de diviseurs
de $A$ normal une structure de groupe analytique, qui « algébriquement »
devrait plutôt exister pour le complété.

La lettre pose ensuite des questions. L'hypothèse de singularité isolée
pourrait être remplacée par « $A$ réduit » si l'on savait que $R^1 f_*\mathcal{O}_X = 0$
pour toute résolution $f$ ; Hironaka l'a démontré pour $Y$ régulier, « but I
wonder if this is \emph{really} needed » --- et une main a écrit au crayon en
regard : « of course it is ». Grothendieck s'attend à des énoncés analogues
pour $\pi_1$ et les autres invariants topologiques, liés à l'annulation de ces
invariants sur les fibres géométriques de $\operatorname{Spec}\hat{A} \to \operatorname{Spec} A$,
ce qui est facile pour $\pi_0$ et vrai pour tout anneau hensélien « bon ». Il
conjecture que la plupart de ces résultats valent pour tout « bon » anneau
hensélien, en particulier, pour l'anneau local $A$ d'une algèbre de type fini,
que tout faisceau inversible sur $\hat{Y}'$ provient d'un faisceau sur un
voisinage étale $Y_1$ de $Y$ à extension résiduelle triviale --- et demande ce
que donne l'exemple de Mumford, p.~16 de son article. Un « bon » anneau est un
quotient d'un anneau local régulier $B$ tel que les fibres de
$\operatorname{Spec}\hat{B} \to \operatorname{Spec} B$ soient universellement
régulières. La lettre se clôt sur des nouvelles personnelles.\footnote{« formal »
et « $\pi_0$ » sont de sa main ; les symboles que la machine n'a pas frappés
($\to$, $\in$, $\ni$) sont restitués par la transcription. Ces « bons »
anneaux sont, à peu de chose près, les anneaux excellents de EGA~IV.}

\subsection*{Le Picard d'un schéma formel en caractéristique $0$ (pages 232 à 244)}

\emph{Lemme.} Soit $\varphi \colon G \to H$ un homomorphisme de type fini de
groupes localement de type fini sur un corps $k$. (i) Le morphisme induit
$\pi_0(\varphi) \colon G/G^0 \to H/H^0$ est de type fini, de noyau fini.
(ii) Si $\operatorname{Ker}\varphi$ est connexe, le conoyau de $\pi_0(\varphi)$
et celui de $\varphi$ ne changent pas quand on remplace $\varphi$ par
l'immersion de son image. (iii) Si de plus $\operatorname{Coker}\varphi$ est
géométriquement sans torsion, $\pi_0(\varphi)$ est injectif et de conoyau sans
torsion. Pour (iii), on se ramène à une immersion fermée d'un $G$ discret ;
alors $H^0/(G \cap H^0)$ s'injecte dans le conoyau, sans torsion, ce qui force
$H^0$ à être vectoriel en caractéristique $0$, et réduit à l'élément neutre en
caractéristique $p$ ; le sous-groupe fini $G \cap H^0$ est donc
nul.\footnote{La formule d'abord écrite pour
$\operatorname{Ker}\pi_0(\varphi)$ dans (ii) est barrée et remplacée par
l'énoncé d'invariance ; l'étiquette « Corollaire » de (iii) est barrée. La
démonstration du cas a), caractéristique $p$, n'est lue qu'en partie ; la
réduction que nous donnons est celle que l'énoncé demande.}

\emph{Corollaire 1.} Soit $(G_n)$ un système projectif de groupes commutatifs
localement de type fini sur $k$, tel que $G_0/G_0^0$ soit de type fini et que
chaque transition $G_{n+1} \to G_n$ soit de type fini, de noyau connexe et de
conoyau géométriquement sans torsion. Alors le système $(\pi_0(G_n))$ est
essentiellement constant ; les systèmes $(G_n^0)$ et $(G_n)$ satisfont
Mittag-Leffler ; et pour $n \geq n_0$ le noyau de $G_n \to G_{n_0}$ est vectoriel
en caractéristique $0$, presque radiciel en caractéristique $p$.
\emph{Corollaire 2.} Si les noyaux des transitions sont unipotents et $k$
parfait, le système $(G_n(k))$ satisfait Mittag-Leffler, puisque
$H^1(k, N) = 0$ pour un groupe unipotent connexe $N$ sur un corps parfait.

\emph{Théorème (page 238).} Soit $\mathfrak{X}$ un schéma formel propre sur un
anneau local noethérien complet $A$ contenant un corps de caractéristique $0$,
de corps résiduel fini sur ce corps. Le système des $\mathrm{Pic}_{X_n/k}$
satisfait les conditions du corollaire 1, et le système des
$\mathrm{Pic}_{X_n/k}(k)$ satisfait Mittag-Leffler. On utilise l'existence des
$\mathrm{Pic}_{X_n/k}$ (Murre ; Grothendieck--Oort dans le cas projectif), la
théorie des obstructions d'Oort pour les noyaux et conoyaux des transitions, et
le théorème de Néron pour $\mathrm{Pic}_{X_0/k}$. Pages 242 et 244 : la suite
exacte
$0 \to \mathrm{Pic}(X_n) \to \mathrm{Pic}_{X_n/k}(k) \to \check{H}^2(k'/k, M_n)$,
$M_n$ le groupe des unités de $H^0(X_n, \mathcal{O}_{X_n})$, dont le système
satisfait Mittag-Leffler, montre que $(\mathrm{Pic}(X_n))$ satisfait lui aussi
Mittag-Leffler--Artin--Rees, et que les images universelles se correspondent.
\emph{Corollaire.} $\mathfrak{X}/Y$ est projectif si et seulement si chaque
$X_n/Y_n$ l'est.\footnote{La page 240 est très surchargée : deux insertions, en
partie en anglais (« becomes »), sont de place incertaine, et les dernières
lignes, barrées, ne se lisent qu'en fragments. La page 245, écrite tête-bêche
et sans lien apparent, pose une question sur les tores maximaux d'un
sous-groupe.}

\pagerange{247}{271}
\section*{Un brouillon d'article : « Un théorème de géométrie analytique » (pages 247 à 271)}

Le brouillon met en forme le théorème de la lettre à Hironaka. Il est écrit
au verso de tapuscrits étrangers au dossier, et sa prose de liaison est rapide
et largement illisible ; les énoncés et les formules passent.

\subsection*{§~1. Le théorème analytique}

\emph{Lemme 1.1.} Si $Y$ est une partie fermée d'un espace topologique $X$
ayant un système fondamental de voisinages paracompacts, et $F$ un faisceau
abélien, $\varinjlim_U H^*(U, F) \to H^*(Y, F|_Y)$ est un isomorphisme, $U$
parcourant les voisinages de $Y$. \emph{Corollaire 1.2.} Pour
$f \colon X \to Y$ propre, $Y$ localement compact, $(R^* f_*F)_y \simeq H^*(X_y, F|_{X_y})$.\footnote{La
page dit « application continue fermée », le second mot ajouté et lu sous
réserve ; il faut aussi les fibres compactes, c'est-à-dire $f$ propre, pour que
les $f^{-1}(V)$ forment un système fondamental de voisinages paracompacts de la
fibre.}

\emph{Lemme 1.3 (Grauert).} Pour $f$ propre entre espaces analytiques complexes
et $F$ cohérent : les $R^i f_*F$ sont cohérents ; avec $Y_n$, $X_n$, $F_n$ comme
plus haut, le système $(H^i(X_n, F_n))_n$ satisfait Mittag-Leffler, et
$(R^i f_*F)^\wedge_y \to \varprojlim_n H^i(X_n, F_n)$ est un isomorphisme.
\emph{Corollaire 1.4.} Si $(R^i f_*F)_y$ est artinien, c'est-à-dire si $y$ est
isolé dans le support, $(R^i f_*F)_y \simeq \varprojlim_n H^i(X_n, F_n)$.

La suite exponentielle $0 \to \mathbf{Z} \to \mathcal{O}_{X_n} \to \mathcal{O}^*_{X_n} \to 0$
donne un système projectif de suites exactes longues, avec
$H^i(X_n, \mathbf{Z}) = H^i(X_0, \mathbf{Z})$ constant et de type fini, et
$H^i(X_n, \mathcal{O}_{X_n})$ satisfaisant Mittag-Leffler ; on en déduit que
noyaux et images satisfont Mittag-Leffler, les suites de sous-groupes de type
fini d'un groupe de type fini dans un vectoriel complexe étant stationnaires.
\emph{Proposition 1.5.} $(H^i(X_n, \mathcal{O}^*_{X_n}))_n$ satisfait
Mittag-Leffler, et l'on a une suite exacte
\[
H^i(X_0, \mathbf{Z}) \to (R^i f_*\mathcal{O}_X)^\wedge_y \to \varprojlim_n H^i(X_n, \mathcal{O}^*_{X_n}) \to H^{i+1}(X_0, \mathbf{Z}) \to (R^{i+1} f_*\mathcal{O}_X)^\wedge_y .
\]
La comparaison avec la suite des tiges, par le diagramme (D),

\begin{tikzcd}[column sep=small, row sep=small, nodes={font=\scriptsize}]
H^i(X_0, \mathbf{Z}) \arrow[r] \arrow[d, no head, "="] & (R^if_*\mathcal{O}_X)_y \arrow[r] \arrow[d, "u^i"] & (R^if_*\mathcal{O}^*_X)_y \arrow[r] \arrow[d, "w^i"] & H^{i+1}(X_0, \mathbf{Z}) \arrow[r] \arrow[d, no head, "="] & (R^{i+1}f_*\mathcal{O}_X)_y \arrow[d, "u^{i+1}"] \\
H^i(X_0, \mathbf{Z}) \arrow[r] & (R^if_*\mathcal{O}_X)^{\wedge}_y \arrow[r] & \varprojlim_n H^i(X_n, \mathcal{O}^*_{X_n}) \arrow[r] & H^{i+1}(X_0, \mathbf{Z}) \arrow[r] & (R^{i+1}f_*\mathcal{O}_X)^{\wedge}_y
\end{tikzcd}

montre, par le lemme des cinq : a) $w^i$ est injectif, puisque $u^i$ et
$u^{i+1}$ le sont (théorème de Krull) ; b) $w^i$ est un isomorphisme si et
seulement si $u^i$ est surjectif, c'est-à-dire si et seulement si $y$ est
isolé dans le support de $R^i f_*\mathcal{O}_X$.\footnote{Le corollaire 1.6 qui
suit ne se lit pas, et un passage encadré de la page 255, barré de traits
ondulés, n'est lu que par bribes. La troisième ligne du diagramme (D) est
barrée.}

\emph{Remarque 1.7} (pages 255 et 257) : c'est le théorème de la lettre, pour
$\mathrm{Pic}(X_y) = (R^1 f_*\mathcal{O}^*_X)_y$ et
$\mathrm{Pic}(\hat{X}_y) = \varprojlim \mathrm{Pic}(X_n)$ ; la condition (iii)
est vérifiée en particulier si $f$ est un isomorphisme au-dessus de
$Y - \{y\}$. \emph{Corollaire 1.8.} $f$ est projectif au-dessus d'un voisinage
de $y$ si et seulement si chaque $X_n$ est un espace algébrique projectif.
\emph{Corollaire 1.9.} Si les fibres sont de dimension $\leq 1$, $f$ est
localement projectif, puisqu'un espace analytique compact de dimension
$\leq 1$ est projectif.\footnote{Le numéro 1.9 sert deux fois sur la page, au
corollaire et à ce lemme. Une rédaction barrée de la page 257 renvoie à
Grauert, 1960--1961.} \emph{Remarque 1.10} (page 261) : les conditions de 1.8
ne suffisent pas en général à rendre $f$ projectif ; on construirait des
contre-exemples par éclatement au-dessus d'une partie discrète $T$ de $Y$, en
des points à anneaux locaux factoriels. \emph{Remarque 1.11} : une
démonstration purement algébrique, par $\mathrm{Pic}$ et le théorème de Néron,
est annoncée.\footnote{Les remarques 1.10 et 1.11 ne sont lues que par
fragments.}

\subsection*{§~2. Application à la géométrie algébrique (pages 263 à 271)}

\emph{Théorème 2.1.} Soient $A$ une algèbre analytique sur $\mathbf{C}$,
$Y = \operatorname{Spec} A$, $X$ propre sur $Y$, et
$\hat{X} = X \times_Y \operatorname{Spf}\hat{A}$, de sorte que
$\mathrm{Pic}(\hat{X}) \simeq \varprojlim \mathrm{Pic}(X_n)$. Alors (i)
$(\mathrm{Pic}(X_n))$ satisfait Mittag-Leffler ; (ii) $\mathrm{Pic}(X)$ et
$\mathrm{Pic}(\hat{X})$ ont même image dans $\mathrm{Pic}(X_n)$, les images
universelles ; (iii) $\mathrm{Pic}(X) \to \mathrm{Pic}(\hat{X})$ est un
isomorphisme dès que $H^1(X, \mathcal{O}_X)$ est de longueur finie, c'est-à-dire
$\operatorname{Supp} R^1 f_*\mathcal{O}_X \subset \{y\}$. Il se déduit de 1.7
par GAGA.

\emph{Corollaire 2.2.} $X$ est projectif sur $Y$ si et seulement si $\hat{X}$
l'est sur $\operatorname{Spec}\hat{A}$, si et seulement si chaque $X_n$ l'est
sur $Y_n$. Pour la dernière implication : soit $n_0$ tel que l'image de
$\mathrm{Pic}(X_{n_0}) \to \mathrm{Pic}(X_0)$ soit universelle, et $L_{n_0}$
ample sur $X_{n_0}$ ; sa restriction à $X_0$ se relève à $\hat{X}$, donc, par
2.1 (ii), provient d'un $L$ sur $X$, ample par EGA~III. \emph{Corollaire 2.3.}
Si $\dim X_0 \leq 1$, $X$ est projectif sur $Y$.

\emph{Corollaire 2.4.} Si $A$ est réduite, normale, à singularité isolée en
$y$, et $Y' = Y - \{y\}$, $\hat{Y}'$ de même pour le complété,
$\mathrm{Pic}(Y') \to \mathrm{Pic}(\hat{Y}')$ est un
isomorphisme.\footnote{La page corrige « monomorphisme » en
« isomorphisme » ; les hypothèses « réduite » et « à singularité isolée » sont
des ajouts encadrés, d'ordre de lecture incertain.} L'injectivité vient d'une
réduction au cas de profondeur $\geq 2$ ; pour la surjectivité, une résolution
de Hironaka $X \to Y$, isomorphe au-dessus de $Y'$, avec $X$ régulier, donne
\[
\mathrm{Pic}(Y') \simeq \mathrm{Pic}(X)/\langle X_0^i \rangle, \qquad
\mathrm{Pic}(\hat{Y}') \simeq \mathrm{Pic}(\hat{X})/\langle X_0^i \rangle ,
\]
$X_0^i$ les composantes de la fibre exceptionnelle, et 2.1 (iii) donne
$\mathrm{Pic}(X) \simeq \mathrm{Pic}(\hat{X})$, puisque $R^1 f_*\mathcal{O}_X$
est supporté en $y$.\footnote{La page écrit « $\mathrm{Pic}(X)/I \simeq \mathbf{Z}^I$ »,
$I$ l'ensemble des composantes : le quotient est par le sous-groupe qu'elles
engendrent, et le « $\simeq \mathbf{Z}^I$ » ne peut pas se lire tel quel ; nous
ne le gardons pas.}

\emph{Questions (page 271).} Pour $A$ local noethérien hensélien et « bon »,
$X$ propre sur $A$ : les images universelles dans $\mathrm{Pic}(X_n)$ de
$\mathrm{Pic}(\hat{X})$ proviennent-elles de $\mathrm{Pic}(X)$ ? En particulier,
$\hat{X}$ projectif entraîne-t-il $X$ projectif, par exemple pour le hensélisé
d'un anneau local de la géométrie algébrique sur $\mathbf{C}$ ? Si
$\operatorname{Supp} R^1 f_*\mathcal{O}_X \subset \{y\}$, a-t-on
$\mathrm{Pic}(X) \simeq \mathrm{Pic}(\hat{X})$ ? Si $Y$ est à singularité isolée,
a-t-on $\mathrm{Pic}(Y - \{y\}) \simeq \mathrm{Pic}(\hat{Y} - \{\hat{y}\})$ ? La
deuxième et la résolution donneraient la troisième.\footnote{Une note marginale
en anglais demande ce que fait « Murre » dans les « Publications », p.~16 ; la
lettre à Hironaka cite pour la même page l'article de Mumford, ce qui suggère
une confusion de nom ou de lecture.}

\pagerange{274}{280}
\section*{Schémas de Néron--Severi et polarisations (pages 273 à 280)}

Les feuillets sont d'un autre papier et d'une autre encre que les précédents.

\emph{Définition.} Un élément de $\mathrm{Pic}(X/S)$ est \emph{numériquement
équivalent à $0$} si sa restriction à chaque fibre l'est ; la page précise :
si un multiple se déforme en $0$.\footnote{Ce que la parenthèse décrit est la
$\tau$-équivalence. Sur une fibre propre sur un corps, elle coïncide avec
l'équivalence numérique : c'est le théorème de Matsusaka, étendu par Kleiman,
que discutent les pages 67 à 71.} On note $\mathrm{Pic}^\tau(X/S)$ le
sous-groupe de ces éléments et
$\mathrm{NS}'(X/S) = \mathrm{Pic}(X/S)/\mathrm{Pic}^\tau(X/S)$. Pour
$S' \to S$ surjectif, $\mathrm{Pic}^\tau(X/S)$ est l'image inverse de
$\mathrm{Pic}^\tau(X'/S')$, donc $\mathrm{NS}'(X/S) \to \mathrm{NS}'(X'/S')$ est
injectif. En rendant le foncteur $S' \mapsto \mathrm{NS}'(X'/S')$ exact pour
les morphismes fidèlement plats localement de type fini, on obtient
$\mathrm{NS}(X/S)$, l'ensemble des « polarisations » de $X$ sur
$S$.\footnote{La page écrit « Nér » pour ce que nous notons $\mathrm{NS}$.}

Quand $\mathrm{Pic}_{X/S}$ existe, avec $\mathrm{Pic}^\tau_{X/S}$ ouvert, un
élément de $\mathrm{NS}(X/S)$ est un sous-schéma ouvert de $\mathrm{Pic}_{X/S}$
qui est, localement pour la topologie fidèlement plate, un torseur trivial sous
$\mathrm{Pic}^\tau_{X/S}$. \emph{Proposition.} Si $\mathrm{Pic}_{X/S}$ existe et
$\mathrm{Pic}^\tau_{X/S}$ est fidèlement plat de type fini sur $S$, le foncteur
$\mathrm{NS}_{X/S}$ est représentable si et seulement si le quotient
$\mathrm{Pic}_{X/S}/\mathrm{Pic}^\tau_{X/S}$ existe, la projection étant
fidèlement plate de type fini de noyau $\mathrm{Pic}^\tau_{X/S}$ ; il est alors
représenté par ce quotient. \emph{Corollaire 1.} Ce schéma de Néron--Severi
est alors non ramifié sur $S$, et séparé si et seulement si
$\mathrm{Pic}^\tau_{X/S}$ est fermé dans $\mathrm{Pic}_{X/S}$. \emph{Corollaire 2.}
Si $\mathrm{Pic}^\tau_{X/S}$ est propre et plat et si $\mathrm{Pic}_{X/S}$ est
réunion croissante d'ouverts quasi-projectifs stables par $\mathrm{Pic}^\tau$
--- ce qui découle, pour $X$ projectif et plat, de ce que $\mathrm{Pic}^\tau$
est de type fini ---, le schéma de Néron--Severi existe, séparé et non ramifié.
Il devrait exister dès que $\mathrm{Pic}^\tau_{X/S}$ est plat et fermé, ce qui
touche à la caractérisation des foncteurs représentables par des objets non
ramifiés ; « au commencement ! », et même l'existence de
$\mathrm{Pic}_{X/S}$ ne semble pas indispensable.

\subsection*{Schémas polarisés (pages 278 à 280)}

Soit $X$ projectif, plat et « steinien » ($f_*\mathcal{O}_X = \mathcal{O}_S$)
sur $S$ localement noethérien, avec $\mathrm{Pic}_{X/S}$
existant.\footnote{La marge juge l'hypothèse « excessive, par exemple pour
l'énoncé de K.\ et Matsusaka ».} Les classes de faisceaux inversibles modulo
$\tau$-équivalence sont les \emph{néronisations strictes} ; en rendant ce
foncteur compatible aux descentes fidèlement plates de type fini, on obtient
les \emph{néronisations} : une néronisation est donnée par un faisceau
inversible $L$ sur $X_T$, $T \to S$ fidèlement plat de type fini, tel que
$p_1^*L \otimes p_2^*L^{-1}$ soit $\tau$-équivalent à $0$ sur
$X_{T \times_S T}$. Quand le schéma quotient
$\mathrm{Pic}_{X/S}/\mathrm{Pic}^\tau_{X/S}$ existe « raisonnablement »,
par exemple si $\mathrm{Pic}^\tau_{X/S}$ est propre, les néronisations sont ses
sections sur $S$, ce qui justifie le nom.\footnote{La page écrit d'abord
« pseudo-polarisations » et « prépolarisations », qu'elle barre pour
« néronisations ».}

Les classes de faisceaux relativement amples définissent un sous-foncteur ; ses
sections sont les \emph{polarisations} de $X/S$. Elles forment un monoïde
commutatif sans élément neutre ; deux polarisations $\xi$, $\xi'$ sont
équivalentes s'il existe des entiers $n, n' > 0$ avec $n\xi = n'\xi'$. Les
classes de polarisations correspondent aux demi-droites rationnelles d'un cône
convexe $C$ dans le $\mathbf{Q}$-espace vectoriel qu'il engendre --- le cône
ample. Toute polarisation est équivalente à une polarisation très ample.

\pagerange{282}{291}
\section*{Exemples (pages 281 à 291)}

\subsection*{Dévissage pour un composé (page 282)}

Pour $X \xrightarrow{g} Y \xrightarrow{f} S$, avec $g$ quasi-compact séparé et
$g_*\mathcal{O}_X = \mathcal{O}_Y$ universellement, on a une suite exacte
\[
0 \to \mathrm{Pic}_{Y/S} \to \mathrm{Pic}_{X/S} \to \textstyle\prod_{Y/S}\mathrm{Pic}_{X/Y} ,
\]
$\prod_{Y/S}$ désignant la restriction de Weil, exacte à droite et scindée si
$X/Y$ a une section, d'où
$\mathrm{Pic}_{X/S} \simeq \mathrm{Pic}_{Y/S} \times \prod_{Y/S}\mathrm{Pic}_{X/Y}$.
Pour $X = Y \times_S Z$, $Y$ et $Z$ munis de sections :
$\mathrm{Pic}_{X/S} \simeq \mathrm{Pic}_{Y/S} \times \mathrm{Pic}_{Z/S} \times \mathrm{Hom}_{\text{pointés}}(Y, \mathrm{Pic}_{Z/S})$.

\subsection*{Les quotients $(A \times_S B)/\Gamma$ et l'exemple d'Igusa (pages 283 à 287)}

Soient $A$ un schéma abélien sur $S$, $B$ un $S$-schéma, $\Gamma$ un groupe
fini opérant sur $A$ par automorphismes de groupe et librement sur $B$, et
$X = (A \times_S B)/\Gamma$, fibré en formes de $A$ au-dessus de $Y = B/\Gamma$,
avec la section image de $0 \times B$. Le dévissage donne
$\mathrm{Pic}_{X/S} \simeq \mathrm{Pic}_{Y/S} \times \prod_{Y/S}\mathrm{Pic}_{X/Y}$,
et le dernier facteur est $\mathrm{Hom}_{S,\Gamma}(B, \mathrm{Pic}_{A/S})$. Si
$\Gamma$ opère sur $B$ par translations et que les coinvariants de $A$ sont
nuls, les homomorphismes $\Gamma$-équivariants de $A$ dans
$\mathrm{Pic}^\tau_{B/S}$, sur lequel $\Gamma$ opère trivialement, sont nuls, et
\[
\mathrm{Pic}_{X/S} \simeq \mathrm{Pic}_{Y/S} \times \mathrm{Pic}_{A/S}^{\Gamma}, \qquad
\mathrm{Pic}^\tau_{X/S} \simeq \mathrm{Pic}^\tau_{Y/S} \times (A^\vee)^{\Gamma} .
\]
\emph{Exemple.} $\Gamma = \mathbf{Z}/2$ opérant sur $A$ par $x \mapsto -x$, et sur
un second schéma abélien $B$ par translation par une section partout d'ordre
exactement $2$ ; alors $Y$ est un schéma abélien et
\[
\mathrm{Pic}^\tau_{X/S} \simeq Y^\vee \times A^\vee[2] .
\]
Il est plat sur $S$, lisse en caractéristique résiduelle $\neq 2$, et non lisse
en un point de caractéristique résiduelle $2$ : c'est l'exemple d'Igusa, qu'on
peut réaliser avec deux courbes elliptiques. En caractéristique $2$, $-1$
opère trivialement sur $H^1(A, \mathcal{O}_A)$, et
$\dim H^1(X, \mathcal{O}_X) = \dim B + \dim A$ alors que
$\dim \mathrm{Pic}_{X/k} = \dim B$.\footnote{La page n'écrit pas « en
caractéristique $2$ » ; hors de la caractéristique $2$, les invariants de $-1$
sur $H^1(A, \mathcal{O}_A)$ sont nuls et $h^1 = \dim B$.} Le $H^1$ saute donc
avec la caractéristique.

La page 286 calcule
$p_*\Omega^1_{X/S} = \pi_{A*}(\Omega^1_{A/S})^\Gamma \times \pi_{B*}(\Omega^1_{B/S})^\Gamma$,
d'où sur un corps, avec $t_A$ l'espace tangent,
\[
\dim H^1(X, \mathcal{O}_X) - \dim H^0(X, \Omega^1_X) = \dim (t_{A^\vee})^\Gamma - \dim (t_A^\vee)^\Gamma ,
\]
et annonce des exemples où cette différence est positive, négative, ou varie
dans une famille, où $(A^\vee)^\Gamma$ n'est pas plat, et où il n'y a pas de
« bon » sous-schéma abélien. On aurait ainsi : 1° la symétrie de Hodge fausse
dans les deux sens ; 2° la semi-continuité « continue » fausse pour
$H^0(X, \Omega^1)$ et $H^1(X, \mathcal{O})$ ; 3° $\mathrm{Pic}^\tau_{X/S}$ non
plat ; 4° « la formation du Pic réduit n'est pas compatible avec la
spécialisation ». Ces exemples sont annoncés, et ne sont pas construits dans le
dossier.

\subsection*{Un homomorphisme divisible par $p$ sur la fibre spéciale seulement (pages 288 et 289)}

Soient $V$ un anneau de valuation discrète d'inégales caractéristiques,
de caractéristique résiduelle $p$, $Y = \operatorname{Spec} V$, $X$ un schéma
abélien sur $Y$ de fibres $X_0$ (spéciale) et $X_1$ (générique), et $F$, $G$
deux sous-groupes finis et plats de $X$ tels que (i) $F_1 \not\subset G_1$ ;
(ii) $F_0 \subset G_0 \neq 0$ ; (iii) $pF_1 \subset G_1$. Soient
$\alpha \colon X \to A = X/F$ et $\beta \colon X \to B = X/G$. Alors $p\beta$
s'annule sur $F$ et se factorise en $u \colon A \to B$ avec
$u\alpha = p\beta$. Génériquement, $u$ n'est pas divisible par $p$ : si
$u = pv$, alors $v\alpha = \beta$, les groupes d'homomorphismes étant sans
torsion, et $F_1 \subset G_1$. Sur la fibre spéciale, $u_0$ est divisible par
$p$, puisque $F_0 \subset G_0$.\footnote{La page écrit (iii) « $pF_1 = G_1$ »,
l'indice surchargeant un autre chiffre ; jointe à (ii), cette égalité est
incompatible avec des sous-groupes plats de même ordre. La réalisation que la
page donne ensuite (deux sous-groupes d'ordre $p$) ne demande que
$pF_1 \subset G_1$, qui y est automatique ; la page 289 écrit d'ailleurs
« $pF = 0 \subset G$ ». La page écrit aussi $p\,\mathrm{Hom}(A_0, F_0)$ pour
$p\,\mathrm{Hom}(A_0, B_0)$.}

Réalisation : prendre $X_0$ sans point d'ordre $p$, $X_1$ avec des points
d'ordre $p$, et $F_1$, $G_1$ deux sous-groupes distincts d'ordre $p$ ; si la
$p$-torsion de la fibre spéciale n'a qu'un sous-groupe d'ordre $p$, leurs
adhérences coïncident en $F_0 = G_0$. C'est le cas pour une courbe elliptique
supersingulière, et l'on peut prendre $F$ et $G$ dans un facteur elliptique
supersingulier d'un produit de deux courbes elliptiques.\footnote{La page
prend « pour $X$ un produit de deux courbes elliptiques » et invoque l'unicité
du sous-groupe d'ordre $p$ dans la $p$-torsion de $X_0$ ; pour un produit de
deux courbes supersingulières, il y en a une infinité. L'unicité vaut dans un
facteur, et c'est là qu'il faut choisir $F$ et $G$.} Avec
$C = A \times B^\vee$, on a
$\mathrm{NS}_{C/Y} \simeq \mathrm{NS}_{A/Y} \times \mathrm{NS}_{B/Y} \times \mathrm{Hom}_{Y\text{-gr}}(A, B)$,
et $u$ définit une section $\bar{u}$ de $\mathrm{NS}_{C/Y}$, divisible par $p$
sur la fibre spéciale et non globalement : la divisibilité dans le schéma de
Néron--Severi ne se propage pas par générisation.

\subsection*{Groupes divisibles (pages 290 et 291)}

Les groupes de torsion divisibles sont les sommes directes de
$\mathbf{Q}_p/\mathbf{Z}_p$ ; un groupe $p$-primaire divisible est
$(\mathbf{Q}_p/\mathbf{Z}_p)^{(I)}$, $|I|$ étant la dimension sur $\mathbf{F}_p$
de $M[p]$. Sont équivalents : $M$ est $p$-primaire de $M[p]$ fini ; $M$ est
isomorphe à $F \oplus (\mathbf{Q}_p/\mathbf{Z}_p)^\rho$, $F$ un $p$-groupe fini,
$\rho$ et $F$ étant alors déterminés ($\rho$ est le rang de la $p$-torsion du
plus grand sous-groupe divisible $M^0$, et $F = M/M^0$). Le foncteur
contravariant $T_p(M) = \mathrm{Hom}(M, \mathbf{Q}_p/\mathbf{Z}_p)$ est une
antiéquivalence avec les $\mathbf{Z}_p$-modules de type fini ; le foncteur
covariant $T^p(M) = \varprojlim M[p^n]$ (le module de Tate) vérifie
$T^p(M) = \mathrm{Hom}(T_p(M), \mathbf{Z}_p)$, tue les groupes finis, est exact
modulo les groupes finis, et donne une équivalence, modulo les groupes finis,
avec les $\mathbf{Q}_p$-espaces de dimension finie. Pour un homomorphisme
$M \to N$ entre groupes $(\mathbf{Q}_p/\mathbf{Z}_p)^n$, sont équivalents :
$T_p(N) \to T_p(M)$ injectif de conoyau libre ; le même après réduction
modulo $p$ ; $T^p(M) \to T^p(N)$ surjectif ; $M \to N$ surjectif sur un
facteur direct ; $M[p] \to N[p]$ surjectif.\footnote{La marge renvoie aux
modules de Dieudonné.}

\pagerange{293}{303}
\section*{Notes, et contre-exemples pour le recollement (pages 293 à 303)}

\subsection*{Déformations (pages 293 et 296)}

Notes au verso ou en travers de feuillets dactylographiés étrangers au
dossier. La page 293 relie la platitude de $\mathrm{Pic}_{X/S}$ en un point à
une condition sur les déformations du premier ordre : l'obstruction à
prolonger un faisceau inversible $L$ le long d'une déformation de classe de
Kodaira--Spencer $\kappa \in H^1(X_0, T_{X_0})$ est le produit de $\kappa$ par
la classe $c(L) \in H^1(X_0, \Omega^1_{X_0})$, à valeurs dans
$H^2(X_0, \mathcal{O}_{X_0})$ : $\rho\partial(x) = 0 \iff \rho\, c(x) = 0$. La
condition est automatique pour les points de $\mathrm{Pic}^0$, « mais il est
douteux que $\mathrm{Pic}^0_{X/K}$ soit pour autant plat sur $S$ ». La page 296
fait les calculs correspondants pour les quotients des pages 283 à 287 :
$H^1(X, \mathcal{O}_X)$, $H^1(Y, T_Y)$ et une suite
$0 \to \mu_2 \to \mathrm{Pic}_Y \to \mathrm{Pic}_X^\Gamma \to \mu_2$.\footnote{Les
exposants de $\mathrm{Pic}$ sont douteux sur ces deux pages, et plusieurs flèches
sont accompagnées de mots biffés ; on n'en donne que le sens général.}

\subsection*{Recoller deux schémas le long d'un fermé (pages 297 à 301)}

Soient $\varphi_i \colon Y \to X_i$ deux immersions, et
$X = X_1 \amalg_Y X_2$ le recollement. \emph{Exemple 1.}\footnote{La page 297
dit seulement « exemple » ; c'est la page 301, qui commence par « Exemple 2 »,
qui le fait premier.} $X_1 = X_2 = \mathbf{A}^2 - \{0\}$
et $Y = \mathbf{G}_m = \operatorname{Spec} k[t, t^{-1}]$, plongé par
$x = t, y = 0$ dans le premier et $x = 1/t, y = 0$ dans le second. Une fonction
sur $X$ est un couple $(f_1, f_2)$ avec $f_1(t, 0) = f_2(1/t, 0)$, ce qui force
la restriction à $Y$ à être constante : les fonctions ne séparent pas les
points de $Y$, et $X$ n'est pas affine. La page passe ensuite au quotient par
$\mathbf{Z}/2$ opérant par symétrie, pour lequel
$(X_1 \amalg X_2)/(\mathbf{Z}/2) \simeq X_1$ ; elle n'est lue qu'en partie.

\emph{Exemple 2 (page 301).} $X_1 = X_2 = \mathbf{P}^2$, $Y$ une courbe
elliptique plongée deux fois, par des faisceaux $\xi_1 = \varphi_1^*\mathcal{O}(1)$
et $\xi_2 = \varphi_2^*\mathcal{O}(1)$ linéairement indépendants sur $\mathbf{Z}$
dans $\mathrm{Pic}(Y)$. Alors tout faisceau inversible sur $X_1 \amalg_Y X_2$
est trivial : il est donné par $\mathcal{O}(n_1)$, $\mathcal{O}(n_2)$ et un
isomorphisme $n_1\xi_1 \simeq n_2\xi_2$ sur $Y$, qui force $n_1 = n_2 = 0$, et
l'isomorphisme restant, une constante, se prolonge. Le recollement n'est donc
pas projectif, bien que sa normalisée $\mathbf{P}^2 \amalg \mathbf{P}^2$ le
soit.\footnote{La transcription porte « $\mathbf{P}^r$ » puis « $r = 2$ ».
Ces recollements le long de fermés sont ce qu'on appelle aujourd'hui des
pincements (Ferrand).}

\subsection*{Contracter un sous-schéma fini (page 303)}

Soient $X$ plat et quasi-projectif sur $S$, $Y$ un sous-schéma fermé fini et
plat sur $S$, et $X' = X \amalg_Y S$, obtenu en contractant $Y$ sur $S$, avec
$p \colon X \to X'$ et la section $\varepsilon \colon S \to X'$. Comparant les
suites $0 \to \mathcal{O}^* \to R^* \to \mathcal{D} \to 0$ des unités, des
fonctions méromorphes inversibles et des diviseurs sur $X$ et $X'$ :
$\mathcal{O}^*_{X'} \to p_*\mathcal{O}^*_X$ est injectif de conoyau
$\varepsilon_*(f_*\mathcal{O}^*_Y/\mathcal{O}^*_S)$, et
$R^*_{X'/S} \to p_*R^*_{X/S}$ est bijectif ; d'où les suites exactes
\[
0 \to f_*\mathcal{O}^*_Y/\mathcal{O}^*_S \to f'_*\mathcal{D}_{X'/S} \to f_*\mathcal{D}_{X/S} \to 0 ,
\qquad
0 \to f_*\mathcal{O}^*_Y/\mathcal{O}^*_S \to \mathrm{Pic}_{X'/S} \to \mathrm{Pic}_{X/S} ,
\]
la seconde si $f_*\mathcal{O}_X = \mathcal{O}_S$.\footnote{La page note la
seconde suite avec un $\mathcal{D}$ souligné et barré, que la transcription lit
« $\underline{\mathrm{Div}}$ » sous réserve ; déduite de la suite des unités,
c'est une suite de groupes de Picard, et c'est ainsi que nous la lisons. C'est
la jacobienne généralisée des pages 141 à 144 : contracter $Y$ ajoute un tore
au Picard.}

\pagerange{305}{305}
\section*{Picard et Albanese des variétés non complètes : Chevalley (pages 304 et 305)}

Soient $X \times Y \to Y$ une action d'un schéma algébrique $X$ sur $Y$
(intègres sur $k$ algébriquement clos) et $D$ un diviseur sur $Y$ tel que
$x \cdot D - D$ soit linéairement équivalent à $0$ pour tout $x \in X(k)$ ---
c'est le cas si $X$ est d'irrégularité nulle et qu'un point opère
trivialement. Alors il existe un système linéaire $P$ contenant $D$ et un
morphisme $X \times P \to P$ qui, ensemblistement, est $(x, D) \mapsto x \cdot D$.

\pagerange{308}{314}
\section*{Questions ouvertes (pages 307 à 314)}

La chemise porte « Questions ouvertes ». Les pages sont un bilan, en
plusieurs encres, où des réponses ont été ajoutées après coup.

\subsection*{Sur la structure de $\mathrm{Pic}_{X/S}$, $S$ noethérien (pages 308 à 310)}

\begin{itemize}
  \item $\mathrm{Pic}_{X/S}$ est-il somme disjointe de schémas de type fini sur
    $S$ ? « Non en général ; mais génériquement sur $S$. » L'adhérence d'une
    partie quasi-compacte est-elle quasi-compacte ? $\mathrm{Pic}^\tau_{X/S}$
    est-il fermé ? (Oui si $X$ est normal et projectif sur $S$, par la
    géométrie formelle.)
  \item Existence du quotient $\mathrm{NS}_{X/S} = \mathrm{Pic}_{X/S}/\mathrm{Pic}^\tau_{X/S}$.
  \item $\mathrm{Pic}^\tau_{X/S}$ est-il universellement ouvert, plat ? « Non »
    pour la platitude : contre-exemple de Mumford ; vrai si les fibres sont
    normales et la torsion première aux caractéristiques résiduelles.
  \item $\mathrm{Pic}^\tau_{X/S}$ est-il quasi-compact, donc de type fini ?
    « Oui, toujours. » De type fini équivaut-il à propre ? « Oui. »
  \item Le rang du groupe de Néron--Severi des fibres reste-t-il borné ?
    « Oui. »
  \item La multiplication par $n$ est-elle universellement ouverte, voire plate,
    pour $n$ grand ? Oui pour $n$ premier aux caractéristiques résiduelles,
    où elle est étale ; est-elle de type fini ? « Non », avec un croquis de
    deux droites marquées $p$.
  \item Les composantes connexes de $\mathrm{Pic}_{X/S}$ sont-elles de type
    fini ? « Non ? » --- c'est ce que montre l'exemple des pages 213 à 218.
  \item Semi-continuité des dimensions et de la lissité des fibres, continuité
    de la dimension des Picard des fibres (oui si $X$ est normal sur $S$, par
    le groupe fondamental), $\mathrm{Pic}^\tau$ d'un produit, injectivité sur
    $\mathrm{Pic}^\tau$ de la restriction à une section hyperplane,
    existence de $\mathrm{Pic}^{00}_{X/S}$ dans le cas artinien, continuité de
    la torsion de Néron--Severi, multiplicité de $\mathrm{Pic}^\tau$.\footnote{La
    page 309 porte « Howard » (lecture incertaine) devant cette dernière
    question. Plusieurs réponses marginales ne se lisent qu'en partie.}
\end{itemize}

\subsection*{Séparation (page 312)}

Le Picard est-il séparé si les fibres sont géométriquement intègres ?
$\mathrm{Pic}^\tau$ l'est-il si les fibres sont séparables ? Deux exemples
répondent en partie : un exemple de Mumford, par des familles de courbes
elliptiques, où $\mathrm{Pic}^0_{X/S}$ n'est pas séparé (« mais alors il n'est
pas représentable ») ; et un exemple, figuré par un croquis, où les fibres sont
séparables et $\mathrm{Pic}_{X/S}$ non séparé --- celui des pages 213 à 218 en
est un. Le préschéma $\mathrm{Corr}_S(X, Y)$ des classes de correspondances,
pour $X$, $Y$ propres à fibres géométriquement intègres, existe-t-il ? « Oui
par Raynaud. » Enfin, la dimension des Picard des fibres s'écrit
$\alpha_s - \mu_s + \lambda_s$, et l'on demande sa constance et le sens de
variation des invariants $\alpha$, $\lambda$, $\mu$.\footnote{La signification
de $\alpha$, $\lambda$, $\mu$ n'est pas donnée sur la page ; elle renvoie
vraisemblablement aux parties abélienne, additive et multiplicative.}

\subsection*{$\mathrm{Pic}^\tau$ est ouvert (pages 313 et 314)}

\emph{Proposition.} Soit $X$ propre, plat et de présentation finie sur $S$.
Alors $\mathrm{Pic}^\tau_{X/S}$ est un sous-foncteur ouvert de présentation
finie de $\mathrm{Pic}_{X/S}$ ; si $X \to S$ est projectif à fibres
géométriquement intègres, il est aussi fermé. \emph{Corollaire.} Pour $L$
inversible sur $X$, l'ensemble des $s$ tels que $L_s$ soit $\tau$-équivalent à
$0$ est ouvert et rétrocompact ; il est ouvert et fermé dans le cas projectif
à fibres géométriquement intègres. On se ramène à $S$ noethérien ; la
stabilité par générisation vient d'un critère (viii b$'$), en se ramenant au
cas d'un trait, et la fermeture dans le cas projectif, de la stabilité par
spécialisation, d'un critère (viii e) --- ces critères étant ceux de la liste
qu'invoquait le commentaire des pages 70 et 71, et que le dossier ne contient
pas.

\emph{Questions.} Éliminer l'hypothèse projective. Si $X \to S$ est normal,
$\mathrm{Pic}^0_{X/S}$ est-il fermé, voire propre sur $S$ ? Autrement dit, si
$X$ est normal et propre sur un trait, muni d'une section, tout faisceau
inversible algébriquement équivalent à $0$ sur la fibre générique se
prolonge-t-il en un faisceau sur $X$, encore algébriquement équivalent à $0$
sur la fibre spéciale ? C'est vrai si $\mathrm{Pic}_{X/S}$ est représentable,
ce qui est le cas si $X/S$ est projectif ; peut-être est-ce vrai dès que
$X \to S$ est propre et plat à fibres géométriquement normales, voire
seulement intègres. Le dossier s'arrête sur ce point d'interrogation.

\end{document}
